\documentclass[11pt]{article}

\usepackage[margin=1.1in]{geometry}

\usepackage{amssymb,mathtools,natbib}

% Anchored formal environments for causalsmith papers.
% First mandatory argument = obj_id (machine anchor joining the paper to the
% Lean crosswalk; invisible in the rendered PDF). Optional argument = name.
% Each environment also defines \label{obj:<obj_id>} so prose can use
% \cref{obj:T-1}; cleveref derives the printed kind and number from the target.
\usepackage{amsmath}
\usepackage{mathrsfs}
\usepackage{amsthm}
\usepackage{booktabs}
% Long displays may break across pages; let TeX stretch a little harder before an
% overfull \hbox so wide formulas/tables are less likely to run off the page.
\allowdisplaybreaks
\setlength{\emergencystretch}{3em}
\newtheorem{theorem}{Theorem}
\newtheorem{lemma}{Lemma}
\newtheorem{assumption}{Assumption}
\newtheorem{definition}{Definition}
\newtheorem{proposition}{Proposition}
\newtheorem{remark}{Remark}
\newtheorem{citedresult}{Cited result}
% The amsthm counter is `algorithmbox` (not `algorithm`) so a later `\usepackage{algorithm}` cannot clash.
\newcounter{algorithmbox}
\renewcommand{\thealgorithmbox}{\arabic{algorithmbox}}
\NewDocumentEnvironment{theoremv}{m o s}{\begin{theorem}\label{obj:#1}{\normalfont[\texttt{\detokenize{#1}}]\IfValueT{#2}{\ (#2)}\IfBooleanT{#3}{\verificationfootnotemark}\ }}{\end{theorem}}
\NewDocumentEnvironment{lemmav}{m o s}{\begin{lemma}\label{obj:#1}{\normalfont[\texttt{\detokenize{#1}}]\IfValueT{#2}{\ (#2)}\IfBooleanT{#3}{\verificationfootnotemark}\ }}{\end{lemma}}
\NewDocumentEnvironment{assumptionv}{m o s}{\begin{assumption}\label{obj:#1}{\normalfont[\texttt{\detokenize{#1}}]\IfValueT{#2}{\ (#2)}\IfBooleanT{#3}{\verificationfootnotemark}\ }}{\end{assumption}}
\NewDocumentEnvironment{definitionv}{m o s}{\begin{definition}\label{obj:#1}{\normalfont[\texttt{\detokenize{#1}}]\IfValueT{#2}{\ (#2)}\IfBooleanT{#3}{\verificationfootnotemark}\ }}{\end{definition}}
% A `propositionv` is a proved result tiered as a Proposition (theorem-grade, machine-verified).
\NewDocumentEnvironment{propositionv}{m o s}{\begin{proposition}\label{obj:#1}{\normalfont[\texttt{\detokenize{#1}}]\IfValueT{#2}{\ (#2)}\IfBooleanT{#3}{\verificationfootnotemark}\ }}{\end{proposition}}
% A `remarkv` holds NON-load-bearing interpretive / open-question content (e.g. a causalsmith `oeq:`
% open question). It is neither proved nor assumed; no theorem may depend on it.
\NewDocumentEnvironment{remarkv}{m o s}{\begin{remark}\label{obj:#1}{\normalfont[\texttt{\detokenize{#1}}]\IfValueT{#2}{\ (#2)}\IfBooleanT{#3}{\verificationfootnotemark}\ }}{\end{remark}}
% An `algorithmv` is a DEFINITION-kind object presented as a boxed, numbered-step procedure (an
% estimator / selection rule built in stages). Same anchor contract as `definitionv`; its body is
% ordinary LaTeX whose steps are `\begin{enumerate}` items, so tex2html needs no extra machinery.
% Presented in the CONVENTIONAL algorithm style readers expect from `algorithm`+`algorithmic`:
% a rule above, an "Algorithm N (Title)" caption line, a rule under the caption, the numbered
% steps, and a closing rule. It is NOT a float — the anchor contract requires the object to stay
% where the outline places it, and floats drift. The obj-id tag is suppressed here (it is
% bookkeeping, and a caption line carrying it reads as debug output in a finished paper).
\NewDocumentEnvironment{algorithmv}{m o s}%
  {\par\addvspace{\medskipamount}\noindent\rule{\linewidth}{0.8pt}\par\nobreak\vspace{2pt}%
   \refstepcounter{algorithmbox}\label{obj:#1}%
   \noindent\textbf{Algorithm \thealgorithmbox}\IfValueT{#2}{ \textnormal{(#2)}}%
   \IfBooleanT{#3}{\verificationfootnotemark}\par\nobreak\vspace{2pt}%
   \noindent\rule{\linewidth}{0.4pt}\par\nobreak\vspace{2pt}\normalfont}%
  {\par\nobreak\vspace{2pt}\noindent\rule{\linewidth}{0.8pt}\par\addvspace{\medskipamount}}
% A `citedv` is an IMPORTED external result the paper relies on but does NOT prove
% (a causalsmith `gate_class:"cited"` node). It is machine-checked against the cited
% source at F2.5, never proved here; the fixed provenance note makes that explicit
% so the environment can never be read as a contribution of this paper. The \citep
% attribution is supplied by the surrounding prose (P2), keyed to references.bib.
\NewDocumentEnvironment{citedv}{m o s}{\begin{citedresult}\label{obj:#1}{\normalfont[\texttt{\detokenize{#1}}]\IfValueT{#2}{\ (#2)}\IfBooleanT{#3}{\verificationfootnotemark}\ \textit{(imported result; machine-checked against the cited source, not proved here.)}\ }}{\end{citedresult}}
% \leanref{obj_id}{display text}: an inline first-mention link to a from-note object's
% verified Lean code. In the PDF it renders the display text plainly (the Lean link is a
% web-only affordance); tex2html turns it into a drawer-opening span.
\newcommand{\leanref}[2]{#2}
% For a theorem/lemma whose Lean declaration accepts a source-matched published proposition as an
% input, the starred anchored environment puts the mark in its heading; the generated text remains
% outside the frozen mathematical environment. Keep the combined macro for older paper bundles.
\newcommand{\verificationfootnotemark}{\footnotemark}
\newcommand{\verificationfootnotetext}[1]{\footnotetext{#1}}
\newcommand{\verificationfootnote}[1]{\verificationfootnotemark\verificationfootnotetext{#1}}
% xcolor before hyperref (hyperref would otherwise pull in plain `color` for colorlinks, and a
% later xcolor load clashes). The page-1 identity stamp at the end of this file needs a grey.
\usepackage{xcolor}
\usepackage{graphicx}
\usepackage{eso-pic}
% hyperref follows natbib and the environment macros; cleveref must follow hyperref.
\usepackage[colorlinks=true,linkcolor=blue,citecolor=blue,urlcolor=blue]{hyperref}
\usepackage[capitalise,nameinlink,noabbrev]{cleveref}
% Pin the names explicitly: labels are placed inside the underlying amsthm environments, and these
% declarations make the PDF contract independent of cleveref's theorem-name inference. House style:
% reference names always print with a capital first letter ("Theorem 1", "Definition 2"), in any
% sentence position — hence `capitalise` above and capitalized \crefname forms below.
\crefname{theorem}{Theorem}{Theorems}
\Crefname{theorem}{Theorem}{Theorems}
\crefname{lemma}{Lemma}{Lemmas}
\Crefname{lemma}{Lemma}{Lemmas}
\crefname{assumption}{Assumption}{Assumptions}
\Crefname{assumption}{Assumption}{Assumptions}
\crefname{definition}{Definition}{Definitions}
\Crefname{definition}{Definition}{Definitions}
\crefname{proposition}{Proposition}{Propositions}
\Crefname{proposition}{Proposition}{Propositions}
\crefname{remark}{Remark}{Remarks}
\Crefname{remark}{Remark}{Remarks}
\crefname{citedresult}{Cited result}{Cited results}
\Crefname{citedresult}{Cited result}{Cited results}
\crefname{algorithmbox}{Algorithm}{Algorithms}
\Crefname{algorithmbox}{Algorithm}{Algorithms}
% ---------------------------------------------------------------------------
% Page-1 identity stamp (arXiv style) and the pinned title-block date.
%
% Split of responsibility: THIS file owns the FORMAT (placement, font, colour, punctuation),
% the generated `paper_stamp.tex` owns only the VALUES (number+version, area, revised date,
% DOI, link target). P4 refreshes this template on every emit, so a layout fix reaches every
% bundle from a recompile; and because `paper_stamp.tex` is not one of the P2 assembly
% digest's authored inputs, a DOI that arrives after publication needs only that one small
% file rewritten plus a recompile — never a redraft, never a reassembly.
%
% Defaults first, so a bundle WITHOUT a stamp file compiles exactly as it always did:
% `\cspaperdate` falls back to `\today` and no stamp is drawn.
\providecommand*{\cspaperdate}{\today}  % the \date{} target
\providecommand*{\csstampid}{}          % e.g. CSWP-2026-008v3 (empty ⇒ no stamp at all)
\providecommand*{\csstamparea}{}        % e.g. Stat
\providecommand*{\csstampdate}{}        % e.g. 27 Aug 2026
\providecommand*{\csstampdoi}{}         % e.g. 10.5281/zenodo.123 (empty ⇒ DOI part omitted)
\providecommand*{\csstamplink}{}        % href target (DOI url, else the paper's page; may be empty)
\IfFileExists{paper_stamp.tex}{% GENERATED by CausalSmith P4 — paper identity stamp values. Do not hand-edit.
%
% Field order on the stamp (owned by paper_macros.tex):
%   <number>v<version> · doi:<doi> · [<Area>] <date>   — the DOI is never the last token, so a
%   text extractor cannot run it into whatever follows. Without a DOI that part is omitted.
% Format lives in paper_macros.tex; only the values are here, and this file is NOT one of
% the P2 assembly digest's authored inputs. A DOI arriving after publication therefore needs
% this file rewritten plus a `latexmk` recompile — never a redraft or a reassembly.
\renewcommand*{\csstampid}{CSWP-2026-024v4}
\renewcommand*{\csstamparea}{Stat}
\renewcommand*{\csstampdate}{29 Sep 2026}
\renewcommand*{\csstampdoi}{}
\renewcommand*{\csstamplink}{https://causalsmith.org/papers/stat_discrete_budgetvalue_curve_v1}
\renewcommand*{\cspaperdate}{29 September 2026}
}{}
\makeatletter
% Field order is deliberate: the DOI is NEVER the last token on the line. A trailing identifier
% runs into whatever a text extractor emits next, which makes it unsafe to match mechanically
% (the Zenodo stamp matcher reads exactly this line); the date is a safe terminator.
%   <number>v<version> · doi:<doi> · [<Area>] <date>      — with a DOI
%   <number>v<version> · [<Area>] <date>                  — without
\newcommand{\cs@stampsep}{\,\textperiodcentered\,}
\newcommand{\cs@stamptext}{%
  \csstampid
  \ifx\csstampdoi\@empty\else\cs@stampsep doi:\csstampdoi\fi
  \ifx\csstamparea\@empty\else\cs@stampsep[\csstamparea]\fi
  \ifx\csstampdate\@empty\else\quad\csstampdate\fi
}
\ifx\csstampid\@empty\else
  \definecolor{csstampgrey}{gray}{0.45}
  % Background shipout material: it occupies no space, so the text block and the \thanks
  % footnote are byte-identical to an unstamped compile. The starred form applies to the
  % NEXT page shipped only — i.e. page 1. Placed in the left margin (the text block starts
  % at 1.1in), reading bottom-to-top, in a Times-like face at 8pt.
  \AddToShipoutPictureBG*{%
    \AtPageLowerLeft{%
      \put(\LenToUnit{0.42in},\LenToUnit{1.1in}){%
        \rotatebox{90}{%
          \fontfamily{ptm}\fontsize{8}{10}\selectfont
          \ifx\csstamplink\@empty
            \textcolor{csstampgrey}{\cs@stamptext}%
          \else
            \expandafter\href\expandafter{\csstamplink}{\textcolor{csstampgrey}{\cs@stamptext}}%
          \fi
        }%
      }%
    }%
  }
\fi
\makeatother


\title{Minimax estimation of budgeted treatment value curves with discrete covariates}

\author{CausalSmith\thanks{Machine-generated by the CausalSmith research pipeline. Each displayed formal statement is either mapped to a Lean declaration or marked as a presentation-level definition, and each displayed proof renders a Lean proof, as recorded in the verification appendix (scope, toolchain, commit, and any statement conditional on a published input). Cited work otherwise supplies framing and positioning.}}

\date{\cspaperdate}

\begin{document}

\maketitle

\begin{abstract}
We study the population oracle-value curve indexed by treatment capacity for a binary treatment and outcome on a known discrete covariate alphabet. Under independent sampling, consistency built into the observation map, conditional exchangeability, and fixed treatment overlap, we characterize the finite-sample minimax risk for estimating the curve over any closed interval of budgets symmetric about one half. For every sample size \(n\geq1\), alphabet size \(d\geq2\), and interval endpoint \(b_0\in[0,1/2]\), the all-procedure expected squared supremum risk on \([b_0,1-b_0]\) is of order \(\min\{1,d/[n\log(ed)]\}\), with constants depending only on overlap. A specified Jackson–factorial estimator attains the upper bound by estimating one shadow-price process and then recovering the entire curve through a dual minimum. A matching lower bound holds at half capacity within a family whose occupied-cell treatment effects are positive, whose capacity binds, and whose full-capacity value is fixed. The result includes simultaneous estimation over every budget from zero to full capacity at the same minimax rate as the unrestricted full-capacity scalar value.
\end{abstract}

\section{Introduction}

A treatment capacity turns population value into a curve: each budget specifies the largest expected outcome attainable by a feasible assignment policy. Estimating that curve addresses questions about the value of additional capacity while allowing the most valuable recipients to change with the budget. This paper characterizes the finite-sample precision of that population target for binary treatment and outcome on a known discrete covariate alphabet that may grow with sample size.

Let \(d\) be the known number of covariate cells, \(n\) the number of independent observed records, and \(\epsilon\) a fixed treatment-overlap parameter. At budget \(b\), the oracle value \(V_b(P)\) is the maximum population outcome under policies whose average treatment share is at most \(b\), for a causal law \(P\) satisfying consistency, conditional exchangeability, and overlap. We measure estimation error by the expected squared largest error over \(I=[b_0,1-b_0]\), the budget interval with endpoint parameter \(b_0\in[0,1/2]\), and compare all measurable curve estimators. For every \(n\geq1\), \(d\geq2\), and \(b_0\in[0,1/2]\), \cref{obj:thm:matched-curve-frontier} establishes the risk order
\[
\min\left\{1,\frac{d}{n\log(ed)}\right\},
\]
with comparison constants depending only on \(\epsilon\). The statement includes the full interval \([0,1]\) and the single half-budget value.

The upper bound uses a dual representation of each budget value as the minimum of a common shadow-price process plus the capacity charge, as shown in \cref{obj:prop:dual-representation}. A pilot-local Jackson polynomial and factorial-moment statistics estimate that process across shadow prices; conditional averaging yields a curve estimator based on the observed sample in \cref{obj:def:jf-estimator}. The dual-representation proposition, \cref{obj:prop:dual-representation}, gives the factor-one contraction from uniform process error to curve error; \cref{obj:thm:curve-upper} gives the resulting estimator risk guarantee. This contraction lets estimation of the whole curve attain the minimax rate of the unrestricted full-capacity scalar value. The guarantee applies across the stated causal law class, including null cells, unequal occupied-cell masses, boundary outcome means, and treatment-effect ties.

The matching converse comes from a paired causal family. At half capacity, every occupied-cell treatment effect lies between \(1/8\) and \(3/8\), the available capacity is used, and the full-capacity value equals \(1/2\). In this family, the target obeys the affine identity \(V_{1/2}(P)=1/8+F_P(1/4)\), where \(F_P\) is the dual threshold process. An exact map sends paired samples from two unknown discrete distributions to observations from this causal family. Thus, an estimator of the half-budget value would yield an estimator of their \(L_1\) distance, giving the lower bound from \citet[Theorem 3 and the unknown-both-distributions reduction following Theorem 6]{JiaoHanWeissman2018}. The full-capacity coordinate, when included in the interval, is the unrestricted scalar oracle-value target studied by \citet{CausalSmithUnrestricted2026}.

The analysis connects constrained-allocation value processes \citep[Section 4 and Theorem 4.1.1]{FengHongNekipelov2026} and uniform inference for optimized value functions \citep{FirpoGalvaoParker2023} with large-alphabet functional estimation by polynomial approximation and moment matching \citep{JiaoVenkatHanWeissman2015,CaiLow2011}. Its risk criterion concerns the optimized population value across budgets; policy-learning work evaluates the welfare of an estimated assignment rule \citep{KitagawaTetenov2018,AtheyWager2021,Pellatt2023}. The formal statements and internal proofs are machine-checked in Lean 4, with verification scope recorded in the appendix.

The remainder first develops the connections to related work, then defines the causal class, budget value, and risk criterion. It next states the matched rate and summarizes the estimator, followed by discussion and limitations. The appendices give the supporting process bounds, lower-bound reduction, and verification scope.

\section{Related work}

The budget-indexed oracle value connects treatment allocation under a capacity constraint with inference for optimized value functions. \citet[Assumptions 4.1.1--4.1.3 and Theorem 4.1.1]{FengHongNekipelov2026} derive an asymptotic value-process law for binary constrained allocation using differentiability of the sorting operator, a regular decision boundary, and joint process convergence. Their conclusion concerns the limiting distribution of a plug-in process in that regular regime. Here, for a known \(d\)-cell alphabet and fixed overlap, the estimand is the same type of oracle capacity curve, but the conclusion is a finite-sample minimax bound for expected squared supremum loss over every symmetric budget interval. When \(d\) is fixed the risk is of order \(1/n\); when \(d\) grows it is of order \(\min\{1,d/[n\log(ed)]\}\). The causal law class permits positive-mass treatment-effect ties at decision thresholds. \citet[Section~4]{FirpoGalvaoParker2023} develop uniform tests for marginally optimized value functions using directional derivatives and resampling. Their optimization framework is more general, while the present result gives a rate for the particular causal capacity curve under fixed overlap. \citet[Sections~3.1 and~5]{SverdrupWuAtheyWager2025} evaluate targeting across budgets with multiple treatment arms using Qini curves and asymptotic pointwise confidence intervals at specified budgets. Our target optimizes the population policy separately at each budget and our loss compares an estimated curve with that oracle curve.

\citet[Section~5]{BhattacharyaDupas2012} study welfare from a data-based binary allocation under a budget constraint and give asymptotic confidence intervals for that policy's welfare. \citet[Sections~3 and~5]{LuedtkeVanDerLaan2016} analyze differentiability and inference for the mean outcome under an optimal individualized rule, including nonunique optima. \citet[Introduction]{Hirano2009} develop asymptotic comparisons of statistical treatment rules. These results frame allocation and optimal-value inference; here the estimand is the population oracle optimum at every budget in \(I=[b_0,1-b_0]\), and the guarantee is an all-procedure finite-sample minimax bound for expected squared supremum loss under fixed overlap and a known discrete alphabet. Policy-learning work instead assesses the welfare of a selected rule \citep{KitagawaTetenov2018,AtheyWager2021,Pellatt2023}. The unrestricted discrete scalar result \citep{CausalSmithUnrestricted2026} studies the \(b=1\) coordinate under the same causal class and has the same minimax order \(\min\{1,d/[n\log(ed)]\}\). The factor-one dual contraction in \cref{obj:prop:dual-representation} shows that estimating every coordinate in \(I\) attains that order simultaneously, including when \(I=[0,1]\).

The statistical tools also connect to estimation of functionals on large discrete alphabets. Polynomial approximation and factorial-moment statistics support estimation of nonsmooth functionals \citep{JiaoVenkatHanWeissman2015}, while moment matching supplies lower bounds for such problems \citep{CaiLow2011,JiaoHanWeissman2018}. Here, polynomial approximation is applied to cell contributions across shadow prices, and paired-cell moment matching captures the difficulty of estimating the capacity-indexed value.


\section{Setup and assumptions}

Let \leanref{sym:d}{\ensuremath{d}} denote the known number of covariate cells, with \(d\ge2\), and let \leanref{sym:[d]}{\ensuremath{[d]}} denote their index set. The sample size is \leanref{sym:n}{\ensuremath{n}}; we use \leanref{sym:j}{\ensuremath{j}} for a cell index and \leanref{sym:a}{\ensuremath{a}} for an arm index in \(\{0,1\}\). For nonnegative sequences, \(r_{n,d}\lesssim s_{n,d}\) means \(r_{n,d}\le C s_{n,d}\), where \(C\) is independent of \(n,d\) but may depend on fixed overlap; \(\asymp\) denotes the corresponding two-sided comparison. The norm \(\|h\|_\infty\) is the supremum norm on the indicated domain.

Let \leanref{sym:P}{\ensuremath{P}} be a full-data causal law for the covariate \leanref{sym:X}{\ensuremath{X}}, treatment \leanref{sym:A}{\ensuremath{A}}, and potential outcomes \leanref{sym:Y(a)}{\ensuremath{Y(a)}}. Treatment selects the observed outcome \leanref{sym:Y}{\ensuremath{Y}}, so an observed record is \leanref{sym:O}{\ensuremath{O=(X,A,Y)}}. The observation map makes the relation between the full-data law and its observed margin explicit. For the observed sample \leanref{sym:O^{(n)}}{\ensuremath{O^{(n)}=(O_1,\ldots,O_n)}}, we impose independent superpopulation sampling.

\begin{assumptionv}{ass:iid-sampling}[Independent observed records]
The observed records \(O_1,\ldots,O_n\) are independent and identically distributed under the observed margin of \(P\).
\end{assumptionv}

Independent superpopulation sampling makes the observed margin the common law of every record and supplies the product-law experiment used in the risk criterion \citep{vanDerVaart1998}. The identification conditions concern how that margin relates to potential outcomes. First, consistency connects each recorded outcome to the potential outcome under the received treatment.

\begin{assumptionv}{ass:consistency}[Consistency of observed outcomes]
\(Y=Y(A)\) almost surely.
\end{assumptionv}

Consistency is the standard link between observed and potential outcomes \citep{HernanRobins2020}. The observation map records this link atomwise.

\begin{assumptionv}{ass:conditional-exchangeability}[Conditional exchangeability of treatment]
Under \(P\), \((Y(0),Y(1))\perp A\mid X\).
\end{assumptionv}

Conditional exchangeability then relates treatment assignment to the two potential outcomes within each covariate cell.

\begin{definitionv}{synth_1}[Observations from potential outcomes]
For \(d\in\mathbb N\), let \(P\) be a full-data causal law on \((X,A,Y(0),Y(1))\), where \(X\) ranges over a \(d\)-element alphabet and \(A,Y(0),Y(1)\in\{0,1\}\). The associated full observation is defined atomwise by
\[
(X,A,Y(0),Y(1))\longmapsto
\bigl(X,A,Y(A),Y(0),Y(1)\bigr),
\qquad
Y(A)=
\begin{cases}
Y(1),& A=1,\\
Y(0),& A=0.
\end{cases}
\]
The observed margin of \(P\) is the law of \((X,A,Y(A))\).
\end{definitionv}

Conditional exchangeability permits arm-specific observed outcome means to identify the corresponding cell potential-outcome means \citep{RosenbaumRubin1983}. To state the remaining restriction precisely, we record cell masses and propensities.

\begin{definitionv}{synth_2}[Observed covariate-cell mass]
For a discrete law \(P\) on a covariate alphabet of size \(d\), the observed mass of cell \(j\) is
\[
p_j := P(X=j).
\]
\end{definitionv}

\begin{definitionv}{synth_3}[Product law and propensity]
For a discrete observed-data law \(P\) on a \(d\)-cell covariate alphabet and a sample size \(n\), the product law is the joint law of \(n\) independent records \(O_i=(X_i,A_i,Y_i)\) drawn from \(P\). For each occupied cell \(j\), write \(p_j=P(X=j)>0\) and define its treatment propensity by \(e_j:=P(A=1\mid X=j)\).
\end{definitionv}

The cell mass \leanref{sym:p_j}{\ensuremath{p_j}} determines whether a conditional quantity in cell \(j\) is relevant to population value.

\begin{assumptionv}{ass:overlap}[Fixed treatment overlap]
For every \(j\in[d]\) with \(p_j>0\), the treatment propensity satisfies
\[
\epsilon\le e_j\le 1-\epsilon.
\]
\end{assumptionv}

The known alphabet fixes the cells over which policies are optimized; cell probabilities may be arbitrarily small, so the law class includes rare cells and cells absent from a sample. Binary outcomes give the observed arm-and-outcome categories used by the factorial-moment estimator. For each occupied cell, the treatment propensity \leanref{sym:e_j}{\ensuremath{e_j}} lies in the band set by the fixed overlap parameter \leanref{sym:\epsilon}{\ensuremath{\epsilon}}, where \(0<\epsilon<1/2\). This condition supports identification of both cell means \citep{RosenbaumRubin1983} and places no lower bound on a cell’s probability. The risk constants depend only on fixed overlap; the paired lower-bound family uses propensity \(1/2\) in every occupied cell. Consistency is encoded by the observation map, which records the observed outcome as \(Y(A)\).

\begin{definitionv}{def:model-class}[Causal law class \(\mathcal M_{d,\epsilon}\)]
\(\mathcal M_{d,\epsilon}\), for each integer \(d\ge2\), is the class of laws \(P\) satisfying:
\begin{itemize}
\item \textbf{(Covariate alphabet.)} \(X\in[d]\).
\item \textbf{(Binary variables.)} \(A,Y(0),Y(1)\in\{0,1\}\).
\item \textbf{(Identification conditions.)} \cref{obj:ass:consistency,obj:ass:conditional-exchangeability,obj:ass:overlap}.
\end{itemize}
\end{definitionv}

The class \leanref{sym:\mathcal M_{d,\epsilon}}{\ensuremath{\mathcal M_{d,\epsilon}}} fixes the covariate alphabet, binary variables, and identification conditions for the minimax comparison. We next specify how a population treatment capacity restricts randomized policies.

\begin{definitionv}{def:budget-policy-class}[Budget policy class \(\Pi_b(P)\)]
\[
\Pi_b(P)=\left\{\pi\in[0,1]^d:\sum_{j=1}^d p_j\pi_j\le b\right\}.
\]
\end{definitionv}

A policy assigns each cell a treatment probability \leanref{sym:\pi}{\ensuremath{\pi_j}}. At population capacity \(b\), its average treatment share is constrained by the cell masses; the resulting feasible class is \leanref{sym:\Pi_b(P)}{\ensuremath{\Pi_b(P)}}. The value of assigning every cell to control supplies the baseline for evaluating these policies.

\begin{definitionv}{synth_7}[Control-policy value]
For a full-data law \(P\), let \(p_j=P(X=j)\) and let \(\mu_{0j}=E_P[Y(0)\mid X=j]\) on occupied cells, with \(\mu_{0j}=0\) on null cells. The population value of assigning every cell to control is
\[
B(P)=\sum_{j=1}^d p_j\mu_{0j}.
\]
\end{definitionv}

Population value depends on the arm-specific potential-outcome mean \leanref{sym:\mu_{aj}}{\ensuremath{\mu_{aj}}} in each occupied cell and on the cell treatment effect \leanref{sym:\tau_j}{\ensuremath{\tau_j}}. The control-policy value \leanref{sym:B(P)}{\ensuremath{B(P)}} provides the baseline for evaluating feasible policies.

\begin{definitionv}{synth_4}[Cell potential-outcome means]
For a full-data causal law \(P\) on a \(d\)-cell covariate alphabet, an arm \(a\in\{0,1\}\), and a cell \(j\), let \(p_j=P(X=j)\). On occupied cells, where \(p_j>0\), define
\[
\mu_{aj}=\mathbb E_P\!\left[Y(a)\mid X=j\right].
\]
On null cells, where \(p_j=0\), set \(\mu_{aj}=0\).
\end{definitionv}

\begin{definitionv}{synth_5}[Cell treatment effect]
For a full-data causal law \(P\) on a \(d\)-cell covariate alphabet and a cell \(j\), define the cell treatment effect by
\[
\tau_j := \mu_{1j}-\mu_{0j},
\]
where \(\mu_{aj}\) is the potential-outcome mean for arm \(a\) in cell \(j\).
\end{definitionv}

The zero convention gives \(\tau_j=0\) on null cells and leaves every population value unchanged.

These cell quantities determine the maximum population value attainable at each budget.

\begin{definitionv}{def:budget-value}[Budget value $V_b(P)$ and curve $V(P)$]
The maximum population value at budget \(b\) and the budget-indexed oracle-value curve on \(I\) are
\[
V_b(P)=B(P)+\max_{\pi\in\Pi_b(P)}\sum_{j=1}^d p_j\pi_j\tau_j,
\qquad
V(P)=\bigl(V_b(P):b\in I\bigr).
\]
\end{definitionv}

Fix an endpoint parameter \leanref{sym:b_0}{\ensuremath{b_0}}\(\in[0,1/2]\) and write \leanref{sym:I}{\ensuremath{I=[b_0,1-b_0]}} for the closed budget interval. Maximizing over feasible randomized policies gives the budget value \leanref{sym:V_b(P)}{\ensuremath{V_b(P)}} and its curve on \(I\).

To express this value through the observed margin, we use treatment-outcome masses within each cell and their arm totals.

\begin{definitionv}{synth_9}[Arm and total masses of a four-vector]
For \(u=(u_{ay}:(a,y)\in\{0,1\}^2)\in\mathbb R_+^4\), define the mass in treatment arm \(a\) and the total mass by
\[
s_a(u)=u_{a0}+u_{a1},\qquad
s(u)=s_0(u)+s_1(u)=\sum_{a,y\in\{0,1\}}u_{ay}.
\]
\end{definitionv}

For each cell, let the observed four-mass vector \leanref{sym:q_j}{\ensuremath{q_j}} have coordinates \(q_{ay,j}=P(X=j,A=a,Y=y)\), for \(a,y\in\{0,1\}\). Its entries encode the observed treatment-outcome distribution without conditioning on a possibly null cell. For a generic nonnegative four-vector \leanref{sym:u}{\ensuremath{u}}, write \leanref{sym:s_a(u)}{\ensuremath{s_a(u)}} for its arm mass and \leanref{sym:s(u)}{\ensuremath{s(u)}} for its total mass.

The finite full-data support can equivalently be indexed by potential atoms.

\begin{definitionv}{synth_6}[Potential atom space]
For a covariate alphabet of size \(d\in\mathbb N\), a potential atom is an element of
\[
\{j\in\mathbb N:j<d\}\times\{0,1\}\times\{0,1\}\times\{0,1\}.
\]
\end{definitionv}

A shadow price \leanref{sym:\lambda}{\ensuremath{\lambda}}\(\in[0,1]\) converts the capacity constraint into a threshold comparison within each cell. The stabilized arm contribution \leanref{sym:g_a(u)}{\ensuremath{g_a(u)}}, threshold contribution \leanref{sym:f_\lambda(u)}{\ensuremath{f_\lambda(u)}}, and aggregate process \leanref{sym:F_P(\lambda)}{\ensuremath{F_P(\lambda)}} express that comparison in observed masses.

\begin{definitionv}{def:dual-process}[Dual threshold process $F_P(\lambda)$]
For \(u\in\mathbb R_+^4\), define the stabilized arm contributions, cell contribution, and dual threshold process by
\[
g_a(u)=\frac{s(u)u_{a1}}{\max\{s_a(u),\epsilon s(u)\}},
\qquad
f_\lambda(u)=g_0(u)+\bigl[g_1(u)-g_0(u)-\lambda s(u)\bigr]_+,
\qquad
F_P(\lambda)=\sum_{j=1}^d f_\lambda(q_j).
\]
When all four cell masses vanish (\(u=0\)), the quotient defining \(g_a(u)\) takes the value zero: \(g_a(0)=0\).
\end{definitionv}

\begin{definitionv}{synth_12}[Binding half-budget capacity]
We say the half-budget constraint is binding for \(P\) when there exists an optimal policy \(\pi^\star\in\Pi_{1/2}(P)\) such that
\[
V_{1/2}(P)=B(P)+\sum_{j=1}^d p_j\pi^\star_j\tau_j
\quad\text{and}\quad
\sum_{j=1}^d p_j\pi^\star_j=\frac12.
\]
\end{definitionv}

The zero convention in \cref{obj:def:dual-process} gives \(f_\lambda(0)=0\) and covers empty covariate cells.

Under \cref{obj:ass:conditional-exchangeability,obj:ass:overlap}, the stabilized arm contribution at an occupied cell equals \(p_j\mu_{aj}\). The threshold contribution therefore combines control value with the positive part of the cell treatment gain after charging the shadow price for treatment capacity. The following identity turns those cellwise contributions into the population budget value.

\begin{propositionv}{prop:dual-representation}[Dual representation of budget value]
Let \(P\) be a full-data causal law with observed margin. Suppose:
\begin{itemize}
\item \textbf{Alphabet size.} The known covariate alphabet has size \(d\ge 2\).
\item \textbf{Overlap parameter.} The fixed overlap parameter satisfies \(0<\epsilon<1/2\).
\item \textbf{Model membership.} \(P\in\mathcal M_{d,\epsilon}\) as defined in \cref{obj:def:model-class}.
\end{itemize}
For \(0\le\lambda\le1\), write \(F_P(\lambda)=\sum_{j=1}^{d}f_\lambda(q_j)\) for the dual threshold process formed from the observed margin of \(P\). Then, for every \(b\in[0,1]\), the budget value \(V_b(P)\) of \cref{obj:def:budget-value} satisfies
\[
V_b(P)=\inf_{0\le\lambda\le1}\{b\lambda+F_P(\lambda)\}.
\]
Moreover, for every process \(\widehat F\) bounded on \([0,1]\),
\[
\sup_{b\in[0,1]}
\left|
\inf_{0\le\lambda\le1}\{b\lambda+\widehat F(\lambda)\}
-V_b(P)
\right|
\le
\sup_{\lambda\in[0,1]}
\left|\widehat F(\lambda)-F_P(\lambda)\right|.
\]
\end{propositionv}

The identity represents every budget value as the lower envelope of the same shadow-price process plus \(b\lambda\). Its second assertion transfers a uniform process error bound to a uniform curve error bound with factor one, so a single estimated process serves the entire budget interval. The identity and contraction concern the optimized value and permit multiple maximizing policies or minimizing shadow prices.

The observed sample collects the records on which a curve estimator may depend.

\begin{definitionv}{synth_11}[Finite observed sample]
For \(n,d\in\mathbb N\), the observed sample \(O^{(n)}=(O_i)_{i=1}^n\) consists of \(n\) observed records \(O_i=(X_i,A_i,Y_i)\), each with \(X_i\) in the \(d\)-element covariate alphabet.
\end{definitionv}

We assess curve estimation by the all-procedure expected squared supremum risk \leanref{sym:\mathfrak R_{n,d,\epsilon,b_0}^{\mathrm{curve}}}{\ensuremath{\mathfrak R_{n,d,\epsilon,b_0}^{\mathrm{curve}}}}. The criterion compares bounded-function-valued estimators over the causal law class and the closed interval \(I\).

\begin{definitionv}{def:curve-risk}[Curve risk $\mathfrak R_{n,d,\epsilon,b_0}^{\mathrm{curve}}$]
For \(b_0\in[0,1/2]\), set \(I=[b_0,1-b_0]\). The all-procedure expected squared supremum risk is
\[
\mathfrak R_{n,d,\epsilon,b_0}^{\mathrm{curve}}
=
\inf_{\widehat V}
\sup_{P\in\mathcal M_{d,\epsilon}}
E_P\!\left[
\sup_{b\in I}
\left|\widehat V_b(O^{(n)})-V_b(P)\right|^2
\right],
\]
where the infimum ranges over all measurable maps from the observed sample into bounded functions on \(I\).
\end{definitionv}

Both endpoints belong to \(I\). In particular, \(b_0=0\) gives the full interval \([0,1]\), while \(b_0=1/2\) gives the single budget \(1/2\).


\section{Minimax risk for the budget-value curve}

The shadow-price representation in \cref{obj:prop:dual-representation} reduces estimation of the budget-value curve to estimation of one process. The resulting minimax risk is matched for every interval \(I=[b_0,1-b_0]\), including the full budget range and the single half-budget value.

\begin{theoremv}{thm:matched-curve-frontier}[Matched curve minimax frontier]*
Fix an overlap parameter \(\epsilon\in(0,1/2)\). There exist constants \(0<c_\epsilon\le C_\epsilon<\infty\) such that, for every \(n\ge1\), \(d\ge2\), and \(b_0\in[0,1/2]\), the curve minimax risk of \cref{obj:def:curve-risk} satisfies
\[
c_\epsilon\min\left\{1,\frac{d}{n\log(ed)}\right\}
\le \mathfrak R_{n,d,\epsilon,b_0}^{\mathrm{curve}}
\le C_\epsilon\min\left\{1,\frac{d}{n\log(ed)}\right\}.
\]
For every measurable scalar estimator \(\widehat V_{1/2}\) based on the \(n\) observed records, there is a full-data law \(P\) with the following properties:
\begin{itemize}
\item \textbf{(Causal model.)} \(P\in\mathcal M_{d,\epsilon}\), and each occupied cell has \(p_j>0\), \(\tau_j\in[1/8,3/8]\), and \(\tau_j>0\).
\item \textbf{(Values and binding capacity.)} \(V_1(P)=1/2\), and there is a policy \(\pi\in[0,1]^d\) such that
\[
\sum_{j=1}^{d}p_j\pi_j=\frac12,
\qquad
V_{1/2}(P)=B(P)+\sum_{j=1}^{d}p_j\pi_j\tau_j.
\]
\item \textbf{(Scalar risk.)} Under \(n\) independent observations from the observed margin of \(P\),
\[
\mathbb E_P\!\left[\bigl(\widehat V_{1/2}-V_{1/2}(P)\bigr)^2\right]
\ge c_\epsilon\min\left\{1,\frac{d}{n\log(ed)}\right\}.
\]
\end{itemize}
When \(b_0=0\), every \(P\in\mathcal M_{d,\epsilon}\) also satisfies the full-budget identity
\[
V_1(P)=B(P)+\sum_{j=1}^{d}p_j\max\{\tau_j,0\},
\]
where \(V_b(P)\) is the budget value of \cref{obj:def:budget-value}.
\end{theoremv}
% CAUSALSMITH-CITED-SCOPE-BEGIN thm:matched-curve-frontier
\verificationfootnotetext{\textbf{Formalization scope.} This result uses the published conclusion \citep{JiaoHanWeissman2018} (Theorem 3 and the unknown-both-distributions reduction following Theorem 6). The cited source proof is not formalized here: Lean verifies its use and all remaining steps. If that cited conclusion is false, the portions of this result that depend on it are not certified.}
% CAUSALSMITH-CITED-SCOPE-END thm:matched-curve-frontier

The upper bound applies to all budgets in the chosen interval simultaneously. In the paired family, the half-budget coordinate attains the stated risk order with positive occupied-cell effects, binding capacity, and a full-capacity value fixed at \(1/2\). This family makes the allocation problem visible through the ranking of cell benefits. Simultaneous curve estimation therefore attains the unrestricted full-capacity scalar minimax order under the stated conditions. The reduction uses the bound for estimating the \(\ell_1\) distance between two unknown discrete distributions from \citet[Theorem~3 and the unknown-both-distributions reduction following Theorem~6]{JiaoHanWeissman2018}.

The proof has three steps. First, the shadow-price identity reduces all budget values to one threshold process, and taking a minimum is a contraction in supremum norm. Second, the estimator below controls that process cell by cell: empirical masses suffice for a small alphabet, while pilot localization and factorial moments handle a growing alphabet. Third, the paired family converts the half-budget value into a discrete two-sample \(L_1\) functional, giving the matching lower bound. The appendix contains the polynomial and process estimates behind the second step.

The exact estimator and its uniform upper-bound theorem are stated in the technical appendix, \cref{obj:def:jf-estimator,obj:thm:curve-upper}. Its statistical construction is summarized here.



The estimator first estimates one threshold process across shadow prices, then takes a shadow-price minimum to obtain each budget value. For a small alphabet it inserts empirical cell masses into that process; with a larger alphabet and adequate sample size, pilot counts localize each mass and evaluation counts estimate a local Jackson polynomial by factorial moments. Clipping controls the cell estimates before they are summed, and the sparse-sample branch reports a constant curve. The displayed construction uses an exact conditional average over auxiliary permutations and markings. Its polynomial degree and pilot-width constants serve the risk proof: at ordinary alphabet sizes the degree remains two and the pilot intervals are broad. The rate is a theoretical benchmark for this exact estimator.

In the Jackson--factorial branch, write \(L=\log(ed)\). Pilot counts set cell-specific localization widths; on each resulting rectangle, the Jackson approximation error is bounded by radius-dependent terms divided by the polynomial order \(K\) and \(K^2\) \cref{obj:lem:jackson-boundary-adaptive}. Independent evaluation counts turn the polynomial terms into factorial-moment estimates. The resulting cell bias is bounded by \(C_\epsilon\{\sqrt{p_j/(nL)}+1/(nL)\}\) uniformly over shadow prices, while the centered good-pilot process has expected squared supremum error of order \(d/(nL)\); the bad-pilot envelope has an exponentially small factor in \(L\) \cref{obj:thm:integrated-process-certificate}. Summing the bias over cells and combining it with these fluctuation bounds yields the same process-risk order. The empirical branch covers small alphabets, and the constant branch covers the sparse-sample regime.



Conditional averaging makes the exact curve-valued estimator a function of the observed sample. By \cref{obj:prop:dual-representation}, the error in each budget value is bounded by the uniform error of the estimated threshold process. The same process estimate therefore controls the supremum over \(I\).



The bound measures the squared largest error across the interval before taking expectation, so it also applies when a budget is selected after viewing the estimated curve. The process bound behind it is given in the appendix; the dual identity transfers that bound to budget values without an additional cost for considering several budgets.

The matching converse can be stated directly for estimation at half capacity. It retains positive effects throughout the occupied cells and fixes the value under full provision.

\begin{theoremv}{thm:capacity-active-lower}[Capacity-active risk lower bound]*
Fix an overlap parameter \(0<\epsilon<1/2\). There exists \(c_\epsilon>0\) such that, for every \(n\geq1\), \(d\geq2\), every collection of observed-data laws satisfying \cref{obj:ass:iid-sampling} for each full-data law, and every measurable scalar estimator \(T\) based on the \(n\) observed records, there is a full-data law \(P\) satisfying:
\begin{itemize}
\item \textbf{(Model class.)} \(P\in\mathcal M_{d,\epsilon}\).
\item \textbf{(Full-budget value.)} \(V_1(P)=1/2\), with \(V_b(P)\) as in \cref{obj:def:budget-value}.
\item \textbf{(Effects in occupied cells.)} For every cell \(j\), \(p_j>0\) implies \(1/8\leq\tau_j\leq3/8\).
\item \textbf{(Binding capacity.)} The half-budget constraint is binding for \(P\).
\end{itemize}
For this law,
\[
E_P\!\left[(T-V_{1/2}(P))^2\right]
\geq c_\epsilon\min\left\{1,\frac{d}{n\log(ed)}\right\}.
\]
\end{theoremv}
% CAUSALSMITH-CITED-SCOPE-BEGIN thm:capacity-active-lower
\verificationfootnotetext{\textbf{Formalization scope.} This result uses the published conclusion \citep{JiaoHanWeissman2018} (Theorem 3 and the unknown-both-distributions reduction following Theorem 6). The cited source proof is not formalized here: Lean verifies its use and all remaining steps. If that cited conclusion is false, the portions of this result that depend on it are not certified.}
% CAUSALSMITH-CITED-SCOPE-END thm:capacity-active-lower

Positive treatment effects make the half-budget allocation use its available capacity, while differences across cells determine which recipients are most valuable. The theorem establishes the lower risk benchmark in this capacity-active setting even with the full-budget value fixed. Its reduction again invokes the discrete \(\ell_1\) bound of \citet[Theorem~3 and the unknown-both-distributions reduction following Theorem~6]{JiaoHanWeissman2018}.

Two special cases put the general rate and the budget curve in familiar terms: a fixed covariate alphabet and a common positive effect across occupied cells.

\begin{propositionv}{prop:fixed-alphabet-reduction}[Fixed alphabet rate and identity]*
Fix the overlap parameter \(\epsilon\) and the known covariate alphabet size \(d\) under the following conditions:
\begin{itemize}
\item \textbf{(Overlap.)} \(0<\epsilon<1/2\).
\item \textbf{(Alphabet.)} \(d\ge 2\).
\end{itemize}
For the curve minimax risk in \cref{obj:def:curve-risk}, there are constants \(c,C\) with \(0<c\le C\) such that, for every \(n\ge1\) and every \(b_0\in[0,1/2]\),
\[
\frac{c}{n}\le
\mathfrak R_{n,d,\epsilon,b_0}^{\mathrm{curve}}
\le\frac{C}{n}.
\]

Moreover, let \(P\) belong to the causal model class \(\mathcal M_{d,\epsilon}\) in \cref{obj:def:model-class}. Write \(p_j\) for the mass of cell \(j\), \(\mu_{aj}\) for its arm-\(a\) potential-outcome mean, and \(\tau_j=\mu_{1j}-\mu_{0j}\) for its treatment effect. If some \(c>0\) satisfies \(\tau_j=c\) in every cell with \(p_j>0\), then, for every \(b_0\in[0,1/2]\) and every \(b\in[b_0,1-b_0]\), the budget value in \cref{obj:def:budget-value} satisfies
\[
V_b(P)=B(P)+cb,
\qquad
B(P)=\sum_{j=1}^{d}p_j\mu_{0j}.
\]
\end{propositionv}
% CAUSALSMITH-CITED-SCOPE-BEGIN prop:fixed-alphabet-reduction
\verificationfootnotetext{\textbf{Formalization scope.} This result uses the published conclusion \citep{JiaoHanWeissman2018} (Theorem 3 and the unknown-both-distributions reduction following Theorem 6). The cited source proof is not formalized here: Lean verifies its use and all remaining steps. If that cited conclusion is false, the portions of this result that depend on it are not certified.}
% CAUSALSMITH-CITED-SCOPE-END prop:fixed-alphabet-reduction

With the alphabet and overlap parameter fixed, the matched rate is \(1/n\). Under a common positive effect, every unit of treatment capacity adds the same population value, giving a linear budget-value curve.
In the common-effect clause of \cref{obj:prop:fixed-alphabet-reduction}, the symbol \(c\) denotes an arbitrary positive treatment effect; it is independent of the lower comparison constant also written \(c\) in the risk bound.


\section{Discussion and limitations}

The risk in \cref{obj:def:curve-risk} measures expected squared supremum error over the budget interval \(I\). Thus, \cref{obj:thm:curve-upper} gives a simultaneous guarantee for one estimated value curve. When \(b_0=0\), that guarantee covers every budget from zero to full population capacity. Together with \cref{obj:thm:capacity-active-lower}, the upper bound identifies the minimax rate under the stated conditions.

When \(d/[n\log(ed)]\) is small, the minimax risk has that order; when the ratio is large, it has constant order. The minimax rate contains a logarithmic factor, and the upper bound uses polynomial estimation of the discrete threshold process. In the paired family, differences among positive cell benefits determine their ranking and hence the half-budget allocation value, even as the full-budget value stays fixed. Under \cref{obj:thm:capacity-active-lower}, this single coordinate attains the curve’s risk order. The factor-one dual contraction carries the process estimate to all budgets, and the whole curve attains the unrestricted full-capacity scalar minimax order.

\subsection{Limitations and future work}

The global supremum criterion covers the endpoint when \(b_0=0\); a budget-local analysis could characterize risk as the budget shrinks with sample size. Adaptive simultaneous bands call for a coverage criterion, while policy regret calls for evaluating an estimated allocation. The theorem's estimator is the exact conditional average in \cref{obj:def:jf-estimator}; approximating that average by Monte Carlo draws would require a separate draw-count and approximation-error guarantee before claiming the same rate. Other directions include covariate alphabets learned from data and allocation across multiple treatment arms.


\appendix

\section*{Appendices}

\section{Proofs and verification scope}

\subsection{Estimator and uniform upper bound}

The observed category-cell masses specify the inputs to the exact Jackson--factorial construction.

\begin{definitionv}{synth_8}[Observed category-cell masses]
For a covariate cell \(j\) and a treatment-outcome category \(\zeta=(a,y)\in\{0,1\}^2\), define
\[
q_{\zeta j}=P(X=j,A=a,Y=y).
\]
These four masses are the components of the observed cell vector \(q_j=(q_{\zeta j}:\zeta\in\{0,1\}^2)\).
\end{definitionv}

\begin{algorithmv}{def:jf-estimator}[Jackson–factorial estimator $\widehat V_b^{\mathrm{JF}}$]
Given \(O^{(n)}\), \(d\), and \(I\), the estimator \(\widehat V_b^{\mathrm{JF}}\) uses the tuning quantities \(H_0,\kappa,D_0,L,m,K\) defined below.
\begin{enumerate}
\item Set
\[
H_0=4096,\qquad
\kappa=\frac1{512},\qquad
D_0=16,\qquad
L=\log(ed),\qquad
m=\frac n8,\qquad
K=\max\left\{2,\left\lfloor\frac L{512}\right\rfloor\right\},
\qquad
\mathcal Z=\{0,1\}^2.
\]
Index treatment-outcome categories by \(\zeta=(a,y)\in\mathcal Z\).

\item If \(d<D_0\), form the empirical category-cell masses and threshold process:
\[
\widehat q_{\zeta j}
=
\frac1n\sum_{i=1}^n
\mathbf1\{X_i=j,\ (A_i,Y_i)=\zeta\},
\qquad
\widehat q_j=(\widehat q_{\zeta j})_{\zeta\in\mathcal Z},
\qquad
\widehat F(\lambda)=\sum_{j=1}^d f_\lambda(\widehat q_j).
\]
Proceed to the final step.

\item If \(d\geq D_0\) and \(n<d/L\), output
\[
\widehat V_b^{\mathrm{JF}}=\frac12,\qquad b\in I.
\]
In the remaining branch, \(d\geq D_0\) and \(n\geq d/L\), draw \(M\sim\operatorname{Pois}(n/4)\) independently of \(O^{(n)}\). On \(\{M>n\}\), set
\[
\widetilde F^{\mathrm{cap}}(\lambda)=0,\qquad \lambda\in[0,1].
\]
On \(\{M\leq n\}\), uniformly permute the records, retain the first \(M\), and independently mark each retained record pilot or evaluation with probability \(1/2\) each. Define the pilot and evaluation counts by
\[
N'_{\zeta j}
=
\sum_{i=1}^{M}
\mathbf1\{X_i^*=j,\ (A_i^*,Y_i^*)=\zeta,\ i\text{ is marked pilot}\},
\qquad
N_{\zeta j}
=
\sum_{i=1}^{M}
\mathbf1\{X_i^*=j,\ (A_i^*,Y_i^*)=\zeta,\ i\text{ is marked evaluation}\},
\]
where \(O_i^*=(X_i^*,A_i^*,Y_i^*)\) denotes the \(i\)th record in the uniform permutation.

\item On \(\{M\leq n\}\), for every \(j\) and \(\zeta\), set
\[
\delta=\frac Lm,\qquad
c_{\zeta j}=\frac{N'_{\zeta j}}m,\qquad
h_{\zeta j}=H_0\left\{\sqrt{c_{\zeta j}\delta}+\delta\right\},
\]
\[
\ell_{\zeta j}=\max\{0,c_{\zeta j}-h_{\zeta j}\},
\qquad
u_{\zeta j}=c_{\zeta j}+h_{\zeta j},
\qquad
c^0_{\zeta j}=\frac{\ell_{\zeta j}+u_{\zeta j}}2,
\qquad
R_{\zeta j}=\frac{u_{\zeta j}-\ell_{\zeta j}}2,
\qquad
S_j=\sum_{\zeta\in\mathcal Z}R_{\zeta j}.
\]
All radii \(R_{\zeta j}\) are strictly positive. Let \(c_j^0=(c^0_{\zeta j})_{\zeta\in\mathcal Z}\) and \(R_j=(R_{\zeta j})_{\zeta\in\mathcal Z}\). Define the normalized positive order-four Jackson kernel by
\[
J_K(t)=C_K\left\{\frac{\sin(Kt/2)}{\sin(t/2)}\right\}^{4},
\qquad -\pi\leq t\leq\pi,
\qquad
\int_{-\pi}^{\pi}J_K(t)\,dt=1,
\]
where the quotient takes its continuous value at zero and \(C_K>0\) is the unique normalizing constant. Set \(D=2(K-1)\). For each \(j\) and \(\lambda\), define the unique tensor polynomial \(P_{j,\lambda}\), of coordinatewise degree at most \(D\), through
\[
P_{j,\lambda}(c_j^0+R_j\odot\cos\vartheta)
=
\int_{[-\pi,\pi]^4}
f_\lambda\!\left(c_j^0+R_j\odot\cos(\vartheta+t)\right)
\prod_{\zeta\in\mathcal Z}J_K(t_\zeta)\,dt,
\]
and define its coefficients \(\beta_{j,\alpha}(\lambda)\) by the unique expansion
\[
P_{j,\lambda}(u)-f_\lambda(c_j^0)
=
\sum_{\alpha\in\{0,\ldots,D\}^4}
\beta_{j,\alpha}(\lambda)
\prod_{\zeta\in\mathcal Z}
\left(\frac{u_\zeta-c^0_{\zeta j}}{R_{\zeta j}}\right)^{\alpha_\zeta}.
\]
For integers \(h\geq0\), define
\[
U_h(N;z)
=
\sum_{t=0}^h {h\choose t}(-z)^{h-t}\frac{(N)_t}{m^t},
\qquad
(N)_t=N(N-1)\cdots(N-t+1),
\qquad
(N)_0=1.
\]
The cell estimate, its clipped version, and the auxiliary threshold process are
\[
Z_j(\lambda)
=
f_\lambda(c_j^0)+
\sum_{\alpha\in\{0,\ldots,D\}^4}
\beta_{j,\alpha}(\lambda)
\prod_{\zeta\in\mathcal Z}
\frac{U_{\alpha_\zeta}(N_{\zeta j};c^0_{\zeta j})}
{R_{\zeta j}^{\alpha_\zeta}},
\]
\[
T_j(\lambda)
=
f_\lambda(c_j^0)+
\Pi_{[-d^{1/4}S_j,d^{1/4}S_j]}
\!\left(Z_j(\lambda)-f_\lambda(c_j^0)\right),
\qquad
\widetilde F^{\mathrm{cap}}(\lambda)
=
\sum_{j=1}^d T_j(\lambda).
\]
Here
\[
\Pi_{[r,s]}(x)=\min\{s,\max\{r,x\}\}.
\]
In this branch, average the auxiliary process conditional on the observed sample:
\[
\widehat F(\lambda)
=
E\!\left\{\widetilde F^{\mathrm{cap}}(\lambda)\mid O^{(n)}\right\}.
\]

\item In the empirical and Jackson–factorial branches, report
\[
\widehat V_b^{\mathrm{JF}}
=
\Pi_{[0,1]}
\left(
\min_{0\leq\lambda\leq1}
\{b\lambda+\widehat F(\lambda)\}
\right),
\qquad b\in I.
\]
The output uses the minimum value and requires no optimizer selection. If an implementation stores an optimizer, it stores the least element of the nonempty compact argmin set. These branch rules, the zero process on \(\{M>n\}\), and the least-minimizer convention apply under every data or treatment-effect tie.
\end{enumerate}
\end{algorithmv}

The conditional average in \cref{obj:def:jf-estimator} is an observed-sample statistic. The dual contraction transfers a uniform process bound to every budget in \(I\).

\begin{theoremv}{thm:curve-upper}[Uniform budget-curve upper bound]
Fix the overlap parameter \(\epsilon\). Suppose:
\begin{itemize}
\item \textbf{(Overlap.)} \(0<\epsilon<1/2\).
\item \textbf{(Sample and alphabet.)} \(n\geq 1\) and \(d\geq 2\).
\item \textbf{(Budget interval.)} \(b_0\in[0,1/2]\), and \(I=[b_0,1-b_0]\).
\item \textbf{(Sampling.)} For each causal law \(P\in\mathcal M_{d,\epsilon}\), the \(n\) observed records follow \cref{obj:ass:iid-sampling} for the observed margin of \(P\).
\end{itemize}
There is a constant \(C_\epsilon>0\), depending only on \(\epsilon\), such that the estimator in \cref{obj:def:jf-estimator} and the budget value in \cref{obj:def:budget-value} satisfy
\[
\sup_{P\in\mathcal M_{d,\epsilon}}
E_P\!\left[
\sup_{b\in I}
\left|\widehat V_b^{\mathrm{JF}}-V_b(P)\right|^2
\right]
\leq
C_\epsilon\min\left\{1,\frac{d}{n\log(ed)}\right\}.
\]
\end{theoremv}

The estimator controls the threshold process over the full interval of shadow prices before taking the dual minimum. Since the minimum value is Lipschitz in the supremum norm of that process, one process bound controls every budget in \(I\), including zero and full capacity when \(b_0=0\).

\subsection{Process estimates}

The upper-bound analysis begins with regularity of the cell contribution in the dual representation of \cref{obj:def:dual-process}. Bounded variation controls its dependence on the full interval of shadow prices.

\begin{lemmav}{lem:bv-lipschitz}[Bounded variation Lipschitz bound]
For \(0<\epsilon<1/2\), there is a finite constant \(C_\epsilon\ge 0\) such that, for every \(u,v\in\mathbb R_+^4\),
\[
\|f_{\cdot}(u)-f_{\cdot}(v)\|_{BV([0,1])}
\le C_\epsilon\|u-v\|_1,
\]
where \(\|h\|_{BV([0,1])}=|h(0)|+\operatorname{TV}_{[0,1]}(h)\).
\end{lemmav}

\begin{proof}[Proof of \cref{obj:lem:bv-lipschitz}]
Fix \(u,v\in\mathbb R_+^4\), and write
\[
s_\eta(u)=\sum_{\omega=0}^1u_{\eta\omega},\qquad
s(u)=\sum_{\eta=0}^1s_\eta(u),\qquad
\Delta_{uv}=\sum_{\eta,\omega=0}^1|u_{\eta\omega}-v_{\eta\omega}|.
\]
In particular, \(|s(u)-s(v)|\leq\Delta_{uv}\). For either arm \(\eta\), the stabilized contribution in \cref{obj:def:dual-process} satisfies
\[
0\leq g_\eta(u)\leq s(u),\qquad
|g_\eta(u)-g_\eta(v)|
 \leq(1+\epsilon^{-1})\Delta_{uv}.
\]
Here is a derivation of the second bound, including zero masses. Put
\(m_u=s_\eta(u)\) and \(z_u=u_{\eta1}\), and define \(m_v,z_v\) analogously.
When \(s(u)>0\), the two branches are
\[
g_\eta(u)=
\begin{cases}
z_u/\epsilon,&m_u\leq\epsilon s(u),\\
s(u)z_u/m_u,&m_u\geq\epsilon s(u).
\end{cases}
\]
They agree at equality; when \(s(u)=0\), \(g_\eta(u)=0\). If both vectors use the first branch, their difference is at most
\(\epsilon^{-1}|z_u-z_v|\). If both use the second, \(m_u,m_v>0\), and
\[
\left|\frac{z_u}{m_u}-\frac{z_v}{m_v}\right|
\leq
\frac{|z_u-z_v|+
 |(m_u-z_u)-(m_v-z_v)|}{m_u}.
\]
Since \(z_v/m_v\leq1\) and \(s(u)/m_u\leq\epsilon^{-1}\), splitting the difference of the two products bounds it by
\(|s(u)-s(v)|+\epsilon^{-1}\Delta_{uv}\).

For the mixed case \(m_u\leq\epsilon s(u)\) and
\(m_v\geq\epsilon s(v)\), the identity
\[
\frac{z_u}{\epsilon}-s(v)\frac{z_v}{m_v}
=
\frac{z_u-z_v}{\epsilon}
+\frac{z_v}{m_v}\frac{m_v-\epsilon s(v)}{\epsilon}
\]
and the bounds
\[
0\leq m_v-\epsilon s(v)
\leq(m_v-m_u)-\epsilon\bigl(s(v)-s(u)\bigr),
\qquad 0\leq\frac{z_v}{m_v}\leq1
\]
give the upper bound below. The rearrangement
\[
(z_u-z_v)+\frac{z_v}{m_v}(m_v-m_u)
=\left(1-\frac{z_v}{m_v}\right)(z_u-z_v)
 +\frac{z_v}{m_v}\bigl((m_v-z_v)-(m_u-z_u)\bigr)
\]
shows that its left-hand side has absolute value at most \(\Delta_{uv}\), since \(0\leq z_v/m_v\leq1\). Consequently,
\[
\frac{z_u}{\epsilon}-s(v)\frac{z_v}{m_v}
\leq \frac{(z_u-z_v)+(z_v/m_v)(m_v-m_u)}{\epsilon}
 -\frac{z_v}{m_v}\bigl(s(v)-s(u)\bigr)
\leq(1+\epsilon^{-1})\Delta_{uv}.
\]
For the lower bound, return to the first mixed-case identity: its correction term \((z_v/m_v)(m_v-\epsilon s(v))/\epsilon\) is nonnegative. Hence the difference is at least \((z_u-z_v)/\epsilon\geq-\epsilon^{-1}\Delta_{uv}\). Interchanging \(u,v\) handles the other mixed case. Finally, if either total mass is zero, all its coordinates and its arm contribution vanish; \(0\leq g_\eta\leq s\) and \(|s(u)-s(v)|\leq\Delta_{uv}\) give \(|g_\eta(u)-g_\eta(v)|\leq(1+\epsilon^{-1})\Delta_{uv}\) directly. The same inequalities also establish \(0\leq g_\eta(u)\leq s(u)\) in every branch. % lean: CausalSmith.Stat.DiscreteBudgetvalueCurve.budget_armValue_lipschitz

Define the positive gain and its hinge by
\[
\gamma_u=[g_1(u)-g_0(u)]_+,\qquad
\Phi_u(\lambda)=[\gamma_u-\lambda s(u)]_+,
\]
and similarly for \(v\). The arm bounds give
\(0\leq\gamma_u\leq s(u)\) and \(0\leq\gamma_v\leq s(v)\).
Because \(\lambda s(u)\geq0\) on \([0,1]\),
\[
f_\lambda(u)=g_0(u)+\Phi_u(\lambda).
\]
Taking positive parts and using the arm contribution bound yields
\[
|\gamma_u-\gamma_v|
\leq |g_1(u)-g_1(v)|+|g_0(u)-g_0(v)|
\leq2(1+\epsilon^{-1})\Delta_{uv}.
\]
% lean: CausalSmith.Stat.DiscreteBudgetvalueCurve.budget_positiveGain_lipschitz

Set
\[
\rho_u=
\begin{cases}\gamma_u/s(u),&s(u)>0,\\0,&s(u)=0,\end{cases}
\qquad
\rho_v=
\begin{cases}\gamma_v/s(v),&s(v)>0,\\0,&s(v)=0,\end{cases}
\qquad
\Psi_{uv}(\lambda)=\Phi_u(\lambda)-\Phi_v(\lambda).
\]
Both breakpoints lie in \([0,1]\), including when a total mass vanishes. The positive-part inequality gives, throughout that interval,
\[
|\Psi_{uv}(\lambda)|
\leq |\gamma_u-\gamma_v|+|s(u)-s(v)|
=:\Xi_{uv}.
\]
If \(\rho_u\leq\rho_v\), the three successive affine expressions for
\(\Psi_{uv}\) are
\[
\gamma_u-\gamma_v+\lambda\bigl(s(v)-s(u)\bigr),\qquad
-\gamma_v+\lambda s(v),\qquad 0
\]
on \([0,\rho_u]\), \([\rho_u,\rho_v]\), and \([\rho_v,1]\), respectively.
If \(\rho_v\leq\rho_u\), they are
\[
\gamma_u-\gamma_v+\lambda\bigl(s(v)-s(u)\bigr),\qquad
\gamma_u-\lambda s(u),\qquad 0
\]
on \([0,\rho_v]\), \([\rho_v,\rho_u]\), and \([\rho_u,1]\).
The formulas agree at their shared endpoints, also for coincident breakpoints. Variation on each affine piece equals the absolute difference of its endpoint values, which is at most \(2\Xi_{uv}\) by the preceding pointwise bound. Adding the three pieces gives
\[
\operatorname{TV}_{[0,1]}(\Psi_{uv})
\leq6\bigl(|\gamma_u-\gamma_v|+|s(u)-s(v)|\bigr).
\]
% lean: CausalSmith.Stat.DiscreteBudgetvalueCurve.hinge_difference_variation_le

The constant \(g_0(u)-g_0(v)\) in
\(f_\lambda(u)-f_\lambda(v)=g_0(u)-g_0(v)+\Psi_{uv}(\lambda)\)
leaves variation unchanged. At \(\lambda=0\), the positive-part inequality and the arm contribution bound give
\[
|f_0(u)-f_0(v)|
\leq 2|g_0(u)-g_0(v)|+|g_1(u)-g_1(v)|
\leq\bigl(3(1+\epsilon^{-1})+1\bigr)\Delta_{uv}.
\]
Thus the positive-gain and variation bounds, together with \(|s(u)-s(v)|\leq\Delta_{uv}\), imply
\[
\|f_{\cdot}(u)-f_{\cdot}(v)\|_{BV([0,1])}
\leq\bigl(15(1+\epsilon^{-1})+7\bigr)\Delta_{uv}.
\]
The displayed finite, nonnegative coefficient supplies \(C_\epsilon\).
% lean: CausalSmith.Stat.DiscreteBudgetvalueCurve.thresholdFun_bv_lipschitz
\end{proof}

The bound in \cref{obj:lem:bv-lipschitz} supplies a common regularity measure across shadow prices. The process criterion below collects the corresponding approximation, bias, and fluctuation bounds for the estimator in \cref{obj:def:jf-estimator}.

\begin{definitionv}{def:process-handle}[Uniform process handle $\mathsf H_\epsilon$]
For fixed \(0<\epsilon<1/2\), the uniform process handle \(\mathsf H_\epsilon\) is the set of constants \(C>0\) satisfying the five groups of bounds below for the single pilot-local Jackson polynomial and clipping rule of \(\widehat V^{\mathrm{JF}}\). Membership requires the bounds simultaneously for every \(n,d\) in the Jackson–factorial branch, every observed categorical law \(P\) on the known \(d\)-cell alphabet, every pilot realization, every cell \(j\), and every \(\lambda\in[0,1]\).

Set
\[
L=\log(ed),\qquad
m=\frac n8,\qquad
K=\max\left\{2,\left\lfloor\frac L{512}\right\rfloor\right\},
\qquad
S_j=\sum_{\zeta\in\mathcal Z}R_{\zeta j},
\qquad
v_j=\sqrt{\frac{p_jL}{m}}+\frac Lm.
\]
For the expectations and \(L_2\) norms below, use the uncapped marked-Poisson experiment in which all pilot and evaluation counts \(N'_{\zeta j}\) and \(N_{\zeta j}\) are mutually independent with distribution \(\operatorname{Pois}(m q_{\zeta j})\). Define
\[
G_j
=
\bigcap_{\zeta\in\mathcal Z}
\left\{
|c_{\zeta j}-q_{\zeta j}|
\leq \frac{h_{\zeta j}}4
\right\},
\qquad
\|h\|_{BV([0,1])}
=
|h(0)|+\operatorname{TV}_{[0,1]}(h).
\]
Evaluate \(T_j\), \(\beta_{j,\alpha}\), and \(R_{\zeta j}\) in this experiment. The membership bounds are:
\begin{itemize}
\item \textbf{Coefficient envelope and exponential bound.}
\[
\sum_\alpha\|\beta_{j,\alpha}\|_{BV([0,1])}
\leq C S_j e^{12K},
\qquad
S_j^2 e^{24K+16K^2/L}
\leq C d^{1/16}S_j^2.
\]

\item \textbf{Centered process on the pilot event.}
\[
E\left\|
\sum_{j=1}^d
\left[
\mathbf1_{G_j}\{T_j-f_\cdot(q_j)\}
-
E\mathbf1_{G_j}\{T_j-f_\cdot(q_j)\}
\right]
\right\|_\infty^2
\leq C\frac d{mL}.
\]

\item \textbf{Cellwise bias.}
\[
\sup_{0\leq\lambda\leq1}
\left|ET_j(\lambda)-f_\lambda(q_j)\right|
\leq
C\left\{\sqrt{\frac{p_j}{mL}}+\frac1{mL}\right\}.
\]

\item \textbf{Cellwise envelope outside the pilot event.}
\[
\left\|
\mathbf1_{G_j^c}
\sup_{0\leq\lambda\leq1}
|T_j(\lambda)-f_\lambda(q_j)|
\right\|_{L_2}
\leq
C d^{1/4}e^{-10L}v_j.
\]

\item \textbf{Centered process outside the pilot event.}
\[
E\left\|
\sum_{j=1}^d
\left[
\mathbf1_{G_j^c}\{T_j-f_\cdot(q_j)\}
-
E\mathbf1_{G_j^c}\{T_j-f_\cdot(q_j)\}
\right]
\right\|_\infty^2
\leq C\frac d{mL}.
\]
\end{itemize}
\end{definitionv}

The criterion in \cref{obj:def:process-handle} treats pilot localization and its complement within the same process bound. A second-moment bound for the pilot radii supplies an integrability input for that criterion.

\begin{lemmav}{lem:pilot-radius-second-moment}[Second moment of pilot radii]
Let \(n\geq 1\), \(d\geq 16\), and \(m=n/8\). Let \(P\) be a categorical law on the \(d\)-cell alphabet, with category-cell masses \(q_{\zeta j}\) and cell masses \(p_j=\sum_{\zeta\in\mathcal Z}q_{\zeta j}\), and fix a cell \(j\). In the ideal Poisson experiment, the pilot counts \(N'_{\zeta j}\) are independent with means \(m q_{\zeta j}\). Let \(S_j=\sum_{\zeta\in\mathcal Z}R_{\zeta j}\) be the sum of the four pilot radii in \cref{obj:def:jf-estimator}, and set
\[
L=\log(ed),\qquad H_0=4096,\qquad
v_j=\sqrt{\frac{p_jL}{m}}+\frac{L}{m}.
\]
Then \(S_j^2\) is integrable under the ideal Poisson law, and
\[
\mathbb E[S_j^2]\leq 32H_0^2v_j^2.
\]
\end{lemmav}

\begin{proof}[Proof of \cref{obj:lem:pilot-radius-second-moment}]
Since \(n\geq1\) and \(d\geq16\), we have \(m>0\) and \(L>0\). For each \(\zeta\in\mathcal Z\), the localization intervals in \cref{obj:def:jf-estimator} give
\[
c_{\zeta j}=\frac{N'_{\zeta j}}m,\qquad
h_{\zeta j}
=H_0\left(\sqrt{\frac{N'_{\zeta j}L}{m^2}}+\frac Lm\right),\qquad
R_{\zeta j}
=\frac{c_{\zeta j}+h_{\zeta j}
-\max\{0,c_{\zeta j}-h_{\zeta j}\}}2.
\]
Both \(c_{\zeta j}\) and \(h_{\zeta j}\) are nonnegative. The maximum in the last expression lies between \(c_{\zeta j}-h_{\zeta j}\) and \(c_{\zeta j}+h_{\zeta j}\); hence \(0\leq R_{\zeta j}\leq h_{\zeta j}\). Applying \((x+y)^2\leq2(x^2+y^2)\) yields
\[
R_{\zeta j}^2
\leq 2H_0^2\left\{
\frac{N'_{\zeta j}L}{m^2}
+\left(\frac Lm\right)^2
\right\}.
\]
% lean: CausalSmith.Stat.DiscreteBudgetvalueCurve.pilotRadius_sq_le_count_envelope

There are four categories in \(\mathcal Z\). Cauchy–Schwarz and the coordinate bound therefore give, pointwise,
\[
S_j^2
=\left(\sum_{\zeta\in\mathcal Z}R_{\zeta j}\right)^2
\leq4\sum_{\zeta\in\mathcal Z}R_{\zeta j}^2
\leq8H_0^2\left\{
\frac L{m^2}\sum_{\zeta\in\mathcal Z}N'_{\zeta j}
+4\left(\frac Lm\right)^2
\right\}.
\]
% lean: CausalSmith.Stat.DiscreteBudgetvalueCurve.pilotRadius_sum_sq_integral_le_cellScale

Each pilot count is integrable under its Poisson law, with
\[
\mathbb E\!\left[\sum_{\zeta\in\mathcal Z}N'_{\zeta j}\right]
=\sum_{\zeta\in\mathcal Z}m q_{\zeta j}
=mp_j.
\]
The pointwise bound thus proves integrability of \(S_j^2\) and gives
\[
\mathbb E[S_j^2]
\leq8H_0^2\left\{
\frac{p_jL}{m}+4\left(\frac Lm\right)^2
\right\}.
\]
Finally, \(p_jL/m\geq0\) and \(L/m\geq0\), so
\[
\frac{p_jL}{m}+4\left(\frac Lm\right)^2
\leq4\left(\sqrt{\frac{p_jL}{m}}+\frac Lm\right)^2
=4v_j^2.
\]
Combining the two inequalities proves \(\mathbb E[S_j^2]\leq32H_0^2v_j^2\).
% lean: CausalSmith.Stat.DiscreteBudgetvalueCurve.pilotRadius_sum_sq_integral_le_cellScale
\end{proof}

The bound in \cref{obj:lem:pilot-radius-second-moment} controls the random localization scale by cell mass and the logarithmic alphabet factor. On a localization rectangle, the Jackson approximation has the following coordinatewise bound, including at its boundary.

\begin{lemmav}{lem:jackson-boundary-adaptive}[Boundary-adaptive Jackson approximation]
Let \(P_\lambda\) be the local polynomial for \(f_\lambda\) of \cref{obj:def:dual-process}, formed by the normalized order-four Jackson construction in \cref{obj:def:jf-estimator} with pilot center and radii replaced by \(c^0\) and \(R\), and with the value \(f_\lambda(c^0)\) added to the normalized polynomial evaluation. Suppose that:
\begin{itemize}
\item \textbf{(Parameters.)} The overlap parameter satisfies \(\epsilon>0\), the shadow price satisfies \(\lambda\in[0,1]\), and \(K\) is a positive integer.
\item \textbf{(Rectangle.)} For every \(\zeta\in\mathcal Z\), \(R_\zeta>0\) and \(c^0_\zeta-R_\zeta\geq0\).
\item \textbf{(Evaluation point.)} For every \(\zeta\in\mathcal Z\), \(|y_\zeta-c^0_\zeta|\leq R_\zeta\).
\end{itemize}
Then
\[
\bigl|P_\lambda(y)-f_\lambda(y)\bigr|
\leq
32\bigl\{3(1+\epsilon^{-1})+1\bigr\}
\sum_{\zeta\in\mathcal Z}
\left\{
\frac{\sqrt{(y_\zeta-c^0_\zeta+R_\zeta)(c^0_\zeta+R_\zeta-y_\zeta)}}{K}
+\frac{R_\zeta}{K^2}
\right\}.
\]
\end{lemmav}

\begin{proof}[Proof of \cref{obj:lem:jackson-boundary-adaptive}]
\begin{enumerate}
\item Set
\[
\mathcal R=\prod_{\zeta\in\mathcal Z}
[c^0_\zeta-R_\zeta,c^0_\zeta+R_\zeta],
\qquad
\Delta(\mathbf w,\mathbf w')
=\sum_{\zeta\in\mathcal Z}|w_\zeta-w'_\zeta|,
\qquad
\Lambda_\epsilon=3(1+\epsilon^{-1})+1
\]
for \(\mathbf w,\mathbf w'\in\mathcal R\). The rectangle assumption makes both vectors nonnegative. We first establish
\[
|f_\lambda(\mathbf w)-f_\lambda(\mathbf w')|
\leq\Lambda_\epsilon\Delta(\mathbf w,\mathbf w').
\]

For each arm \(a\in\{0,1\}\), the contribution in \cref{obj:def:dual-process} satisfies
\[
|g_a(\mathbf w)-g_a(\mathbf w')|
\leq(1+\epsilon^{-1})\Delta(\mathbf w,\mathbf w').
\]
Here is the scalar estimate behind this bound. At zero total mass, every coordinate vanishes. We evaluate the quotient defining \(g_a\) in \cref{obj:def:dual-process} with the convention \(0/0=0\), so \(g_a(\mathbf 0)=0\). Since \(0\leq g_a(\mathbf w)\leq s(\mathbf w)\), and likewise for \(\mathbf w'\), the claim follows whenever either vector has zero total mass. For positive total masses, the defining maximum gives \(g_a(\mathbf w)=w_{a1}/\epsilon\) when \(s_a(\mathbf w)\leq\epsilon s(\mathbf w)\), and \(g_a(\mathbf w)=s(\mathbf w)w_{a1}/s_a(\mathbf w)\) in the other regime. When both vectors are in the first regime, their difference is bounded by \(\epsilon^{-1}\Delta(\mathbf w,\mathbf w')\). When both are in the second, \(0\leq w'_{a1}/s_a(\mathbf w')\leq1\) and
\[
\left|
\frac{w_{a1}}{s_a(\mathbf w)}
-\frac{w'_{a1}}{s_a(\mathbf w')}
\right|
\leq
\frac{|w_{a1}-w'_{a1}|+|w_{a0}-w'_{a0}|}
{s_a(\mathbf w)}.
\]
Expanding the difference of the two products, using
\(|s(\mathbf w)-s(\mathbf w')|\leq\Delta(\mathbf w,\mathbf w')\) and
\(s(\mathbf w)/s_a(\mathbf w)\leq\epsilon^{-1}\), proves the displayed arm bound.

For the mixed regime in which \(\mathbf w\) is in the first regime and \(\mathbf w'\) in the second, the difference equals
\[
\frac{w_{a1}-w'_{a1}}{\epsilon}
+
\frac{w'_{a1}}{s_a(\mathbf w')}
\frac{s_a(\mathbf w')-\epsilon s(\mathbf w')}{\epsilon}.
\]
The second term is nonnegative. Its ratio lies in \([0,1]\), and the two regime inequalities imply
\[
0\leq s_a(\mathbf w')-\epsilon s(\mathbf w')
\leq s_a(\mathbf w')-s_a(\mathbf w)
       -\epsilon\{s(\mathbf w')-s(\mathbf w)\}.
\]
Consequently the difference lies between
\(-\epsilon^{-1}\Delta(\mathbf w,\mathbf w')\) and
\((1+\epsilon^{-1})\Delta(\mathbf w,\mathbf w')\).
Interchanging the vectors covers the remaining mixed regime.

Finally, \(x\mapsto[x]_+\) is one-Lipschitz. Applying this to the hinge in \cref{obj:def:dual-process} gives
\[
|f_\lambda(\mathbf w)-f_\lambda(\mathbf w')|
\leq
2|g_0(\mathbf w)-g_0(\mathbf w')|
+|g_1(\mathbf w)-g_1(\mathbf w')|
+|\lambda|\,|s(\mathbf w)-s(\mathbf w')|
\leq\Lambda_\epsilon\Delta(\mathbf w,\mathbf w').
\]
Since \(\Delta(\mathbf w,\mathbf w')\leq4\|\mathbf w-\mathbf w'\|_\infty\), \(f_\lambda\) is continuous on \(\mathcal R\). % lean: CausalSmith.Stat.DiscreteBudgetvalueCurve.budget_thresholdFunReal_pointwise_lipschitz

\item Write \(\beta_\alpha^{c^0,R}(\lambda)\) for the centered coefficients in \cref{obj:def:jf-estimator} with its pilot center and radii replaced by \(c^0,R\). Explicitly, they are the unique coefficients satisfying
\[
P_\lambda(c^0+R\odot\mathbf z)-f_\lambda(c^0)
=
\sum_{\alpha\in\{0,\ldots,2(K-1)\}^{\mathcal Z}}
\beta_\alpha^{c^0,R}(\lambda)
\prod_{\zeta\in\mathcal Z}z_\zeta^{\alpha_\zeta},
\qquad \mathbf z\in\mathbb R^{\mathcal Z}.
\]
The normalized polynomial is therefore
\[
\widetilde P_\lambda(\mathbf z)
=
f_\lambda(c^0)+
\sum_{\alpha\in\{0,\ldots,2(K-1)\}^{\mathcal Z}}
\beta_\alpha^{c^0,R}(\lambda)
\prod_{\zeta\in\mathcal Z}z_\zeta^{\alpha_\zeta},
\qquad \mathbf z\in\mathbb R^{\mathcal Z}.
\]
The continuity just proved applies to
\(\mathbf z\mapsto f_\lambda(c^0+R\odot\mathbf z)\) on \([-1,1]^4\).
The normalized construction, including the added center value, therefore gives, for every \(\boldsymbol\vartheta\in\mathbb R^4\),
\[
\widetilde P_\lambda(\cos\boldsymbol\vartheta)
=
\int_{[-\pi,\pi]^4}
f_\lambda\!\left(c^0+R\odot
\cos(\boldsymbol\vartheta+\boldsymbol\omega)\right)
\prod_{\zeta\in\mathcal Z}J_K(\omega_\zeta)
\,d\boldsymbol\omega.
\]
The order-four kernel has trigonometric degree \(2(K-1)\) in each coordinate, so \(\widetilde P_\lambda\) has coordinatewise polynomial degree at most \(2(K-1)\). Positive radii permit the affine substitution
\[
P_\lambda(\mathbf w)
=
\widetilde P_\lambda\!\left(
\left(\frac{w_\zeta-c^0_\zeta}{R_\zeta}\right)_{\zeta\in\mathcal Z}
\right),
\qquad \mathbf w\in\mathbb R^4,
\]
which preserves those coordinatewise degree bounds and identifies this polynomial with the one in the statement. % lean: Causalean.Mathlib.Analysis.JacksonApproximation.mvPolynomial_affine_substitution_four

\item Put
\[
\vartheta_\zeta
=
\arccos\!\left(\frac{y_\zeta-c^0_\zeta}{R_\zeta}\right),
\qquad \zeta\in\mathcal Z.
\]
The evaluation-point assumption makes every angle well defined in \([0,\pi]\), and \(y=c^0+R\odot\cos\boldsymbol\vartheta\). Positivity and unit integral of each \(J_K\), followed by the Lipschitz estimate from the first step, yield
\[
|P_\lambda(y)-f_\lambda(y)|
\leq
\Lambda_\epsilon\sum_{\zeta\in\mathcal Z}R_\zeta
\int_{[-\pi,\pi]^4}
\left|\cos(\vartheta_\zeta+\omega_\zeta)
-\cos\vartheta_\zeta\right|
\prod_{\xi\in\mathcal Z}J_K(\omega_\xi)
\,d\boldsymbol\omega.
\]
The cosine addition formula and
\(|\sin t|\leq|t|\), \(1-\cos t\leq t^2/2\) give
\[
|\cos(\vartheta_\zeta+\omega_\zeta)-\cos\vartheta_\zeta|
\leq|\sin\vartheta_\zeta|\,|\omega_\zeta|
+\frac{\omega_\zeta^2}{2}.
\]

For completeness, the kernel moments used here obey
\[
\int_{-\pi}^{\pi}|t|J_K(t)\,dt\leq\frac{32}{K},
\qquad
\int_{-\pi}^{\pi}t^2J_K(t)\,dt\leq\frac{64}{K^2}.
\]
Indeed, writing \(D_K(t)=\sum_{r=0}^{K-1}e^{irt}\), the unnormalized fourth power in \cref{obj:def:jf-estimator} is \(|D_K(t)|^4\). Orthogonality of the Fourier modes gives
\[
\int_{-\pi}^{\pi}|D_K(t)|^4\,dt
=2\pi\sum_{h=-(K-1)}^{K-1}(K-|h|)^2
=\frac{2\pi}{3}K(2K^2+1)
\geq\frac{4\pi}{3}K^3.
\]
For \(t\in[-\pi,\pi]\), the inequality
\(t^2\leq\pi^2\sin^2(t/2)\), together with
\(\int_{-\pi}^{\pi}\{\sin(Kt/2)/\sin(t/2)\}^2dt=2\pi K\),
bounds its weighted second integral by \(2\pi^3K\).
Normalization and \(\pi^2\leq16\) thus give the sharper second-moment bound \(24/K^2\). Integrating
\[
|t|\leq\frac K4t^2+\frac1K
\]
gives a first-moment bound \(7/K\); the displayed bounds follow for every positive integer \(K\).

Integrating the cosine estimate coordinatewise now gives
\[
|P_\lambda(y)-f_\lambda(y)|
\leq
32\Lambda_\epsilon
\sum_{\zeta\in\mathcal Z}
\left\{
\frac{R_\zeta|\sin\vartheta_\zeta|}{K}
+\frac{R_\zeta}{K^2}
\right\}.
\]
Finally,
\[
R_\zeta|\sin\vartheta_\zeta|
=
\sqrt{(y_\zeta-c^0_\zeta+R_\zeta)
(c^0_\zeta+R_\zeta-y_\zeta)}.
\]
Substitution proves the asserted bound, including points on the faces of the rectangle. % lean: Causalean.Mathlib.Analysis.Approximation.Chebyshev.AffineFourBoundary.affineJackson_boundary_adaptive_four
\end{enumerate}
\end{proof}

The endpoint factors in \cref{obj:lem:jackson-boundary-adaptive} make the approximation bound sensitive to the evaluation point’s position within the pilot rectangle. To control observations outside the good-pilot event, the next lemma bounds moments of the pilot discrepancy and width together.

\begin{lemmav}{lem:pilot-bad-score-moment}[Bad pilot score moment]
Let \(n,d\geq 1\), let \(P\) be an observed categorical law on the \(d\)-cell alphabet, and fix a cell \(j\). Write \(q_{\zeta j}\) for its four category-cell masses and \(p_j=\sum_{\zeta\in\mathcal Z}q_{\zeta j}\). Set \(m=n/8\) and \(L=\log(ed)\).

Suppose the four pilot counts are independent, with \(N'_{\zeta j}\sim\operatorname{Pois}(mq_{\zeta j})\) for each \(\zeta\in\mathcal Z\). Form \(c_{\zeta j}\) and \(h_{\zeta j}\) as in \cref{obj:def:jf-estimator}, and let \(G_j\) be the quarter-halfwidth pilot event of \cref{obj:def:process-handle}. Define
\[
B_j=\sum_{\zeta\in\mathcal Z}\bigl(|c_{\zeta j}-q_{\zeta j}|+h_{\zeta j}\bigr),
\qquad
v_j=\sqrt{\frac{p_jL}{m}}+\frac{L}{m}.
\]
For each \(t\in\{1,2,4\}\),
\[
\mathbb E\!\left[\mathbf 1_{G_j^c}B_j^{\,t}\right]
\leq
4^t\,10^{80(t+1)}e^{-20L}v_j^{\,t}.
\]
\end{lemmav}

\begin{proof}[Proof of \cref{obj:lem:pilot-bad-score-moment}]
\begin{enumerate}
\item For each \(\zeta\in\mathcal Z\) and \(w\in\mathbb N\), define
\[
\rho_\zeta=mq_{\zeta j},\qquad
\gamma_\zeta=\sqrt{\rho_\zeta L}+L,\qquad
\sigma_\zeta(w)=|w-\rho_\zeta|+1024\bigl(\sqrt{wL}+L\bigr),
\]
and let
\[
\mathscr E_\zeta
=\left\{|N'_{\zeta j}-\rho_\zeta|
>256\bigl(\sqrt{N'_{\zeta j}L}+L\bigr)\right\}.
\]
Here \(m>0\) and \(L=\log(ed)\geq1\), including when \(d=1\). The following scalar estimates hold also when \(q_{\zeta j}=0\):
\[
\Pr(\mathscr E_\zeta)\leq 2e^{-80L},\qquad
\mathbb E\,\sigma_\zeta(N'_{\zeta j})^\nu
\leq 10^{4\nu+4}\gamma_\zeta^\nu
\quad(0\leq\nu\leq8).
\]
For the first estimate, a count \(w\) in \(\mathscr E_\zeta\) with \(w\geq\rho_\zeta\) satisfies
\[
w-\rho_\zeta>\sqrt{160\rho_\zeta L}+160L;
\]
one with \(w<\rho_\zeta\) satisfies
\[
\rho_\zeta-w>\sqrt{160\rho_\zeta L}.
\]
For the upper tail, \(\rho_\zeta\leq w\) gives \(\sqrt{160\rho_\zeta L}\leq16\sqrt{wL}\), so the defining threshold exceeds the stated Bernstein threshold. For the lower tail, the bad-event condition gives \(\rho_\zeta-w>256(\sqrt{wL}+L)\). Since \(L\geq1\),
\[
[256(\sqrt{wL}+L)]^2-160L\,256(\sqrt{wL}+L)-160wL
=65376wL+90112L\sqrt{wL}+24576L^2>0.
\]
Moreover,
\[
\begin{aligned}
&\bigl[(\rho_\zeta-w)^2-160L(\rho_\zeta-w)-160wL\bigr]
-\bigl[[256(\sqrt{wL}+L)]^2-160L\,256(\sqrt{wL}+L)-160wL\bigr]\\
&=\bigl(\rho_\zeta-w-256(\sqrt{wL}+L)\bigr)
\bigl(\rho_\zeta-w+256(\sqrt{wL}+L)-160L\bigr)>0.
\end{aligned}
\]
Thus \((\rho_\zeta-w)^2>160L\{w+(\rho_\zeta-w)\}=160\rho_\zeta L\); positivity of the deviation gives \(\rho_\zeta-w>\sqrt{160\rho_\zeta L}\). The upper and lower Poisson Bernstein bounds at level \(80L\) give \(e^{-80L}\) for each tail.

For the moment estimate, put \(\xi_\zeta=1/(2\gamma_\zeta)\). Since \(\gamma_\zeta\geq1\) and \(\rho_\zeta\xi_\zeta^2\leq1/4\),
\[
\sigma_\zeta(w)
\leq1536\bigl(\gamma_\zeta+|w-\rho_\zeta|\bigr),
\qquad
\mathbb E\exp\!\left\{
\frac{\sigma_\zeta(N'_{\zeta j})}{4096\gamma_\zeta}
\right\}
\leq
e^{1/2}\!\sum_{\pm}
\exp\!\left\{\rho_\zeta
\bigl(e^{\pm\xi_\zeta}-1\mp\xi_\zeta\bigr)\right\}
\leq 2e^{3/4}<8.
\]
The middle expression is the centered Poisson exponential moment; the last inequality uses
\(e^{\pm\xi_\zeta}-1\mp\xi_\zeta\leq\xi_\zeta^2\).
Applying \(x^\nu\leq\nu!e^x\) to the nonnegative exponent and
\(8\nu!4096^\nu\leq10^{4\nu+4}\) for \(0\leq\nu\leq8\) yields the stated moment estimate.

For \(t\in\{1,2,4\}\), the elementary square inequality gives
\[
\mathbf1_{\mathscr E_\zeta}\sigma_\zeta(N'_{\zeta j})^t
\leq\frac12\left[
e^{-40L}\gamma_\zeta^{-t}\sigma_\zeta(N'_{\zeta j})^{2t}
+e^{40L}\gamma_\zeta^t\mathbf1_{\mathscr E_\zeta}
\right].
\]
Taking expectations and using the two estimates above gives
\[
\mathbb E\!\left[
\mathbf1_{\mathscr E_\zeta}\sigma_\zeta(N'_{\zeta j})^t
\right]
\leq
\frac{10^{8t+4}+2}{2}e^{-40L}\gamma_\zeta^t
\leq10^{8t+8}e^{-40L}\gamma_\zeta^t.
\]
% lean: Causalean.Stat.Concentration.PoissonSelfNormalized.poisson_weighted_bad_moment

\item Define
\[
\mathscr E_j=\bigcup_{\zeta\in\mathcal Z}\mathscr E_\zeta,\qquad
\Sigma_j=\sum_{\zeta\in\mathcal Z}\sigma_\zeta(N'_{\zeta j}),
\qquad
\Gamma_j=\sqrt{mp_jL}+L.
\]
For \(\zeta,\eta\in\mathcal Z\), the preceding truncated moment applies when
\(\zeta=\eta\). When \(\zeta\ne\eta\), independence factors the expectation, and the probability and ordinary moment estimates apply. Thus, in both cases,
\[
\mathbb E\!\left[
\mathbf1_{\mathscr E_\zeta}
\sigma_\eta(N'_{\eta j})^t
\right]
\leq
\bigl(10^{8t+8}+2\cdot10^{4t+4}\bigr)
e^{-40L}\gamma_\eta^t.
\]
All scores are nonnegative, so the four-term power inequality and the union bound yield
\[
\mathbb E[\mathbf1_{\mathscr E_j}\Sigma_j^t]
\leq
4^{t-1}\sum_{\zeta,\eta\in\mathcal Z}
\mathbb E\!\left[
\mathbf1_{\mathscr E_\zeta}
\sigma_\eta(N'_{\eta j})^t
\right].
\]
Since \(\sum_\eta\rho_\eta=mp_j\), Cauchy–Schwarz and nonnegativity give
\[
\sum_\eta\gamma_\eta
\leq2\sqrt{mp_jL}+4L\leq4\Gamma_j,
\qquad
\sum_\eta\gamma_\eta^t
\leq4\left(\sum_\eta\gamma_\eta\right)^t.
\]
Consequently,
\[
\mathbb E[\mathbf1_{\mathscr E_j}\Sigma_j^t]
\leq
4^{2t+1}
\bigl(10^{8t+8}+2\cdot10^{4t+4}\bigr)
e^{-40L}\Gamma_j^t
\leq10^{80(t+1)}e^{-20L}\Gamma_j^t.
\]
The final coefficient comparison holds for each of \(t=1,2,4\), and \(L\geq1\) gives \(e^{-40L}\leq e^{-20L}\).
% lean: Causalean.Stat.Concentration.PoissonSelfNormalized.independent_poisson_badAny_moment_four

\item By \cref{obj:def:jf-estimator}, for every \(\zeta\),
\[
c_{\zeta j}=\frac{N'_{\zeta j}}m,\qquad
h_{\zeta j}=\frac{4096}{m}
\bigl(\sqrt{N'_{\zeta j}L}+L\bigr).
\]
The event in \cref{obj:def:process-handle} therefore satisfies
\(G_j^c\subseteq\mathscr E_j\): failure of its quarter-halfwidth condition requires a deviation greater than \(1024(\sqrt{N'_{\zeta j}L}+L)\), which exceeds the threshold defining \(\mathscr E_\zeta\). Pointwise,
\[
B_j
=\frac1m\sum_{\zeta\in\mathcal Z}
\left\{|N'_{\zeta j}-\rho_\zeta|
+4096\bigl(\sqrt{N'_{\zeta j}L}+L\bigr)\right\}
\leq\frac4m\Sigma_j.
\]
Finally, \(\Gamma_j/m=\sqrt{p_jL/m}+L/m=v_j\). Applying the preceding product moment bound proves
\[
\mathbb E[\mathbf1_{G_j^c}B_j^t]
\leq\left(\frac4m\right)^t
\mathbb E[\mathbf1_{\mathscr E_j}\Sigma_j^t]
\leq4^t10^{80(t+1)}e^{-20L}v_j^t.
\]
% lean: CausalSmith.Stat.DiscreteBudgetvalueCurve.pilotCell_empiricalBadScore_moment
\end{enumerate}
\end{proof}

The exponential factor in \cref{obj:lem:pilot-bad-score-moment} controls the contribution of bad pilot localizations. Across cells, bounded variation gives a maximal inequality for the centered paths on either part of the pilot split.

\begin{lemmav}{lem:bv-maximal-cell-paths}[Maximal bound for centered cell paths]
Let \(n,d\) be nonnegative integers, let \(\epsilon\in\mathbb R\), and let \(P\) be a discrete law on the \(d\)-cell alphabet with category-cell masses \(q_{\zeta j}\). In the ideal count experiment of \cref{obj:def:process-handle}, the pilot and evaluation counts \(N'_{\zeta j}\) and \(N_{\zeta j}\) are mutually independent with \(\operatorname{Pois}(mq_{\zeta j})\) laws, where \(m=n/8\). Choose the pilot events uniformly across cells: either \(E_j=G_j\) for every \(j\), or \(E_j=G_j^c\) for every \(j\), with \(G_j\) as in \cref{obj:def:process-handle}. For each cell \(j\), define the path on \([0,1]\) by
\[
W_j(\lambda)=\mathbf 1_{E_j}\bigl\{T_j(\lambda)-f_\lambda(q_j)\bigr\},
\]
where \(T_j\) is the clipped cell estimate of \cref{obj:def:jf-estimator} and \(f_\lambda\) is defined in \cref{obj:def:dual-process}. Assume:

\begin{itemize}
\item \textbf{(Regularity.)} For every count table and cell \(j\), \(W_j\) is continuous in \(\lambda\). For each \(j\), the count-table map into \(C([0,1])\) is measurable and Bochner integrable.
\item \textbf{(Finite variation.)} For each \(j\), \(\operatorname{TV}_{[0,1]}(W_j)<\infty\) almost surely.
\item \textbf{(Second moments.)} For each \(j\), \(\bigl(\|W_j\|_\infty+\operatorname{TV}_{[0,1]}(W_j)\bigr)^2\) is integrable.
\item \textbf{(Independence.)} The paths \(W_1,\ldots,W_d\) are independent.
\end{itemize}

Then
\[
\mathbb E\left\|\sum_{j=1}^{d}\bigl(W_j-\mathbb EW_j\bigr)\right\|_\infty^2
\leq 16384\sum_{j=1}^{d}
\mathbb E\left[\bigl(\|W_j\|_\infty+\operatorname{TV}_{[0,1]}(W_j)\bigr)^2\right].
\]
\end{lemmav}

\begin{proof}[Proof of \cref{obj:lem:bv-maximal-cell-paths}]
\noindent\textit{1. Set up the count law and path envelopes.}
Let \(\mathbb P_{\mathrm{id}}\) be the joint law of the pilot and evaluation count tables in \cref{obj:def:process-handle}. For a count table \(\omega\), write \(W_j(\omega)\) for the path in the statement and set
\[
\rho_j(\omega)
=\|W_j(\omega)\|_\infty+
  \operatorname{TV}_{[0,1]}(W_j(\omega)).
\]
The second-moment assumption makes both \(\rho_j^2\) and
\(\|W_j\|_\infty^2\) integrable. The regularity assumptions give
Bochner-integrable paths, and the independence assumption applies to their
joint law. These observations also cover \(d=0\), when every sum below is
empty. % lean: CausalSmith.Stat.DiscreteBudgetvalueCurve.idealPilotErrorCellPath_centered_maximal

\noindent\textit{2. Reduce the centered sum to signed differences.}
Take independent count tables \(\omega,\omega'\) with law
\(\mathbb P_{\mathrm{id}}\), and define, for \(j=1,\ldots,d\) and
\(\lambda\in[0,1]\),
\[
\Delta_j(\lambda)
=W_j(\lambda;\omega)-W_j(\lambda;\omega').
\]
Let \(\eta_1,\ldots,\eta_d\) be independent uniform signs, independent of
the count tables, and define their signed supremum energy by
\[
\mathcal E_\eta(\Delta)
=\mathbb E_\eta\left\|
  \sum_{j=1}^d\eta_j\Delta_j
 \right\|_\infty^2.
\]
Jensen's inequality applied to the independent copy gives
\[
\mathbb E\left\|
 \sum_{j=1}^d(W_j-\mathbb EW_j)
 \right\|_\infty^2
\le
\mathbb E_{\omega,\omega'}\left\|
 \sum_{j=1}^d\Delta_j
 \right\|_\infty^2.
\]
The pairs \((W_j(\omega),W_j(\omega'))\) are independent across \(j\).
Swapping either entry of any pair preserves their joint law and reverses
the sign of its difference. Averaging these equalities over all sign
choices therefore yields
\[
\mathbb E_{\omega,\omega'}\left\|
 \sum_{j=1}^d\Delta_j
 \right\|_\infty^2
=
\mathbb E_{\omega,\omega'}\mathcal E_\eta(\Delta).
\]
All these expectations are finite by the second-moment assumption.
% lean: Causalean.Stat.Concentration.BoundedVariation.independent_copy_difference_energy_eq_sign_energy

\noindent\textit{3. Bound the signed supremum for fixed count tables.}
Fix a pair \((\omega,\omega')\) for which every \(W_j\) has finite
variation. Let \(\mathcal V_j(t)\), \(t\in[0,1]\), be the cumulative
variation of \(\Delta_j\) from \(0\) to \(t\). Continuity and finite
variation give a continuous, nondecreasing control satisfying
\[
\mathcal V_j(0)=0,\qquad
\mathcal V_j(1)=\operatorname{TV}_{[0,1]}(\Delta_j),\qquad
|\Delta_j(t)-\Delta_j(s)|
 \le \mathcal V_j(t)-\mathcal V_j(s)\quad(s\le t).
\]
Define the common control on \([0,1]\) by
\[
\mathcal U(t)=\sum_{j=1}^d
  \mathcal V_j(1)\mathcal V_j(t).
\]
Since \(0\le\mathcal V_j(t)-\mathcal V_j(s)\le\mathcal V_j(1)\),
\begin{equation}
\sum_{j=1}^d\bigl(\Delta_j(t)-\Delta_j(s)\bigr)^2
 \le \mathcal U(t)-\mathcal U(s)
 \quad(s\le t).
\label{eq:obj:bv-maximal-cell-paths-control}
\end{equation}

Choose ordered, compatible dyadic control quantiles \(t_{k,i}\), for
\(k\ge0\) and \(0\le i\le2^k\), so that
\[
\mathcal U(t_{k,i})=\frac{i}{2^k}\mathcal U(1).
\]
Such choices follow from continuity of \(\mathcal U\). For each sign
vector, let
\[
M_k(\eta)=
\max_{0\le i<2^{k+1}}
\left|
\sum_{j=1}^d\eta_j
 \bigl(\Delta_j(t_{k+1,i+1})-\Delta_j(t_{k+1,i})\bigr)
\right|.
\]
By \cref{eq:obj:bv-maximal-cell-paths-control}, every edge at level \(k+1\) has squared coefficient sum at most
\(\mathcal U(1)/2^{k+1}\). For \(k\ge0\), \(0\le i<2^{k+1}\), and \(1\le j\le d\), write
\[
\gamma_{k,i,j}
=\Delta_j(t_{k+1,i+1})-\Delta_j(t_{k+1,i}),
\qquad
\mathcal Y_k(\eta)=M_k(\eta)^2,
\qquad
\mathfrak v_k=\frac{\mathcal U(1)}{2^{k+1}}.
\]
When \(\mathfrak v_k>0\), each edge satisfies
\(\sum_j\gamma_{k,i,j}^2\le\mathfrak v_k\). For every
\(\chi>0\), independence of the signs and the elementary bound
\(\cosh z\le e^{z^2/2}\) give
\[
\mathbb E_\eta\exp\!\left(\chi\sum_j\eta_j\gamma_{k,i,j}\right)
=\prod_j\cosh(\chi\gamma_{k,i,j})
\le\exp\!\left(\frac{\chi^2\mathfrak v_k}{2}\right).
\]
Applying this bound to both signs, choosing
\(\chi=\sqrt{\mathfrak s}/\mathfrak v_k\), and taking the union over
\(2^{k+1}\) edges gives, for \(\mathfrak s>0\),
\[
\mathbb P_\eta\{\mathcal Y_k\ge\mathfrak s\}
\le 2^{k+2}\exp\!\left(-\frac{\mathfrak s}{2\mathfrak v_k}\right).
\]
Set \(\mathfrak t_k=2\mathfrak v_k\log(2^{k+2})\). Integrating the
tail above \(\mathfrak t_k\) yields
\[
\mathbb E_\eta\mathcal Y_k
\le\mathfrak t_k+
 \int_{\mathfrak t_k}^\infty
 2^{k+2}e^{-\mathfrak s/(2\mathfrak v_k)}\,d\mathfrak s
=2\mathfrak v_k\{\log(2^{k+2})+1\}
\le32\bigl(1+\log(2^{k+1})\bigr)\mathfrak v_k.
\]
When \(\mathfrak v_k=0\), \cref{eq:obj:bv-maximal-cell-paths-control}
makes every edge increment zero. Thus in every case
\begin{equation}
\mathbb E_\eta M_k(\eta)^2
\le
32\bigl(1+\log(2^{k+1})\bigr)
   \frac{\mathcal U(1)}{2^{k+1}}.
\label{eq:obj:bv-maximal-cell-paths-level}
\end{equation}
The numerical factors in \cref{eq:obj:bv-maximal-cell-paths-level}
have a summable square root:
\begin{equation}
\sum_{k=0}^\infty
 \sqrt{\frac{32(1+\log(2^{k+1}))}{2^{k+1}}}
\le
\sum_{k=0}^\infty 8\left(\frac34\right)^k
=32\le40.
\label{eq:obj:bv-maximal-cell-paths-series}
\end{equation}
Here the termwise bound follows from \(\log 2\le1\) and
\(k+2\le4(9/8)^k\).

To reconstruct the signed path, for each \(t\in[0,1]\) and
\(k\ge0\) define
\[
\xi(t)=
\begin{cases}
\mathcal U(t)/\mathcal U(1),&\mathcal U(1)>0,\\
0,&\mathcal U(1)=0,
\end{cases}
\qquad
\iota_k(t)=\min\{2^k-1,\lfloor2^k\xi(t)\rfloor\}.
\]
These lower dyadic indices start at zero, satisfy
\(\iota_{k+1}(t)\in\{2\iota_k(t),2\iota_k(t)+1\}\), and obey
\[
\frac{\iota_k(t)}{2^k}\le\xi(t)
\le\frac{\iota_k(t)+1}{2^k}.
\]
The parent grid point equals the even child by compatibility; the odd
child differs by one adjacent edge. At level zero the grid point and
\(0\) have the same control value, so \cref{eq:obj:bv-maximal-cell-paths-control}
makes their path coordinates equal. Telescoping therefore gives, for
every \(k\ge0\),
\[
\left|\sum_{j=1}^d\eta_j
 \Delta_j(t_{k,\iota_k(t)})\right|
\le
\left|\sum_{j=1}^d\eta_j\Delta_j(0)\right|
+\sum_{\nu=0}^{k-1}M_\nu(\eta),
\]
where the sum is empty at \(k=0\). The control values of
\(t_{k,\iota_k(t)}\) and \(t\) differ by at most
\(\mathcal U(1)/2^k\). Applying
\cref{eq:obj:bv-maximal-cell-paths-control} in whichever order these
two times occur gives
\[
\sum_{j=1}^d
 \bigl(\Delta_j(t_{k,\iota_k(t)})-\Delta_j(t)\bigr)^2
\le \frac{\mathcal U(1)}{2^k}.
\]
Hence the finite signed sums along the chain converge to their value
at \(t\), including on flat parts of the control. If
\(\mathcal U(1)=0\), the same control bound makes all \(\Delta_j\)
constant and the conclusion is immediate. Passing to the limit and
taking the supremum over \(t\) gives
\[
\left\|\sum_{j=1}^d\eta_j\Delta_j\right\|_\infty
\le
\left|\sum_{j=1}^d\eta_j\Delta_j(0)\right|
 +\sum_{k=0}^\infty M_k(\eta).
\]
Sign orthogonality, the \(L_2\) triangle inequality, and
\cref{eq:obj:bv-maximal-cell-paths-level,eq:obj:bv-maximal-cell-paths-series}
imply
\[
\mathcal E_\eta(\Delta)
\le
4096\left\{
 \sum_{j=1}^d\Delta_j(0)^2+\mathcal U(1)
\right\}
\le
4096\sum_{j=1}^d
 \left(\|\Delta_j\|_\infty+
 \operatorname{TV}_{[0,1]}(\Delta_j)\right)^2.
\]
Finally, the triangle inequalities for the supremum norm and total
variation give
\[
\mathcal E_\eta(\Delta)
\le
8192\sum_{j=1}^d
 \bigl(\rho_j(\omega)^2+\rho_j(\omega')^2\bigr).
\]
% lean: Causalean.Stat.Concentration.BoundedVariation.rademacherEnergy_le_pathSize_sq

\noindent\textit{4. Average the pointwise bound.}
Integrating the last inequality over the two independent count tables
and using their identical laws gives
\[
\mathbb E_{\omega,\omega'}\mathcal E_\eta(\Delta)
\le
8192\sum_{j=1}^d
 \mathbb E_{\omega,\omega'}
 \bigl(\rho_j(\omega)^2+\rho_j(\omega')^2\bigr)
=
16384\sum_{j=1}^d\mathbb E\rho_j^2.
\]
Combining this with the two relations in Step 2 proves the stated
inequality. % lean: Causalean.Stat.Concentration.BoundedVariation.centered_path_sum_maximal
\end{proof}

By \cref{obj:lem:bv-maximal-cell-paths}, cellwise path moments yield control uniform in the shadow price. The reported estimator averages over auxiliary draws from an observed sample; the following identity relates that average to the uncapped count experiment.

\begin{lemmav}{lem:capped-marked-count-identity}[Capped marked-count integral identity]
Fix \(n,d\in\mathbb N\) and an observed categorical law \(P\) on the \(d\)-cell alphabet, with category-cell masses \(q_{\zeta j}\). Let \(O^{(n)}=(O_1,\ldots,O_n)\sim P^{\otimes n}\). For \(M\in\{0,\ldots,n\}\), a permutation \(\sigma\) of the \(n\) record indices, and pilot/evaluation marks \(\omega\in\{\mathrm{pilot},\mathrm{evaluation}\}^n\), write \((N',N)(O^{(n)},\sigma,M,\omega)\) for the pilot and evaluation count tables formed from the first \(M\) permuted records according to \cref{obj:def:jf-estimator}. Define
\[
w_n(M)=\frac{\Pr\{\operatorname{Pois}(n/4)=M\}}{n!\,2^n}.
\]
In the uncapped marked-Poisson experiment, \(M\sim\operatorname{Pois}(n/4)\); conditional on \(M\), draw \(M\) independent records from \(P\) and mark each independently as pilot or evaluation with probability \(1/2\), yielding count tables \((N',N)\).

For every real-valued function \(g\) of pairs of count tables that is integrable under the ideal law of \cref{obj:def:process-handle}, whose pilot and evaluation category-cell counts are independent \(\operatorname{Pois}((n/8)q_{\zeta j})\) variables,
\[
\mathbb E_{O^{(n)}\sim P^{\otimes n}}\left[
\sum_{M=0}^{n}\sum_{\sigma}\sum_{\omega\in\{\mathrm{pilot},\mathrm{evaluation}\}^n}
w_n(M)\,
g\bigl((N',N)(O^{(n)},\sigma,M,\omega)\bigr)
\right]
=
\mathbb E_{\mathrm{uncapped}}\left[
\mathbf 1\{M\leq n\}\,g(N',N)
\right].
\]
\end{lemmav}

\begin{proof}[Proof of \cref{obj:lem:capped-marked-count-identity}]
\begin{enumerate}
\item Let \(\mathcal O_d=\{1,\ldots,d\}\times\{0,1\}^2\), and let \(\mathcal S_d\) be the space of finite sequences of records in \(\mathcal O_d\), each carrying a pilot or evaluation mark. Write
\[
\mathbf S=(M,((O_i^\circ,\chi_i))_{i=1}^M)\in\mathcal S_d
\]
for a sequence with law \(\mathsf Q\): \(M\sim\operatorname{Pois}(n/4)\), the \(O_i^\circ\) are independent with law \(P\), and the independent marks \(\chi_i\) are pilot or evaluation with probability \(1/2\) each. Define \(\Phi:\mathcal S_d\to(\mathbb N^{\mathcal Z\times\{1,\ldots,d\}})^2\) by \(\Phi(\mathbf S)=(N'(\mathbf S),N(\mathbf S))\), where
\[
N'_{\zeta j}(\mathbf S)
=\sum_{i=1}^M\mathbf1\{O_i^\circ=(j,\zeta),\ \chi_i=\mathrm{pilot}\},
\qquad
N_{\zeta j}(\mathbf S)
=\sum_{i=1}^M\mathbf1\{O_i^\circ=(j,\zeta),\ \chi_i=\mathrm{evaluation}\}.
\]
These are the count tables used in \cref{obj:def:jf-estimator}. Partition the marked observation space into the singleton types \((j,\zeta,\chi)\), where \(\chi\in\{\mathrm{pilot},\mathrm{evaluation}\}\). Each type has probability \(q_{\zeta j}/2\). For any array \(\boldsymbol\nu=(\nu_{\zeta j,\chi})\in\mathbb N^{\mathcal Z\times\{1,\ldots,d\}\times\{\mathrm{pilot},\mathrm{evaluation}\}}\), the Poisson probability of the total count and the conditional multinomial probability of the singleton counts give
\[
\begin{aligned}
&\Pr_{\mathsf Q}\!\left\{
\#\{i\leq M:O_i^\circ=(j,\zeta),\ \chi_i=\chi\}
=\nu_{\zeta j,\chi}\ \text{for all }(\zeta,j,\chi)\right\}\\
&\quad=e^{-n/4}
\frac{(n/4)^{\sum_{\zeta,j,\chi}\nu_{\zeta j,\chi}}}
{\left(\sum_{\zeta,j,\chi}\nu_{\zeta j,\chi}\right)!}
\frac{\left(\sum_{\zeta,j,\chi}\nu_{\zeta j,\chi}\right)!}
{\prod_{\zeta,j,\chi}\nu_{\zeta j,\chi}!}
\prod_{\zeta,j,\chi}\left(\frac{q_{\zeta j}}2\right)^{\nu_{\zeta j,\chi}}\\
&\quad=\prod_{\zeta,j,\chi}
 e^{-(n/8)q_{\zeta j}}
 \frac{\bigl((n/8)q_{\zeta j}\bigr)^{\nu_{\zeta j,\chi}}}
 {\nu_{\zeta j,\chi}!}.
\end{aligned}
\]
Here \(\sum_{\zeta,j}q_{\zeta j}=1\), and zero powers use the usual convention \(0^0=1\). Grouping the singleton counts by mark gives \(\Phi(\mathbf S)\); the factorization shows that its pilot and evaluation coordinates are independent Poisson variables, each with mean \((n/8)q_{\zeta j}\). Its law is the ideal count law specified in \cref{obj:def:process-handle}; the identity also covers zero cell masses and \(n=0\).
% lean: CausalSmith.Stat.DiscreteBudgetvalueCurve.uncappedMarkedCountLaw_table_eq_idealCountLaw

\item Set \(\varphi:\mathcal S_d\to\mathbb R\) by \(\varphi(\mathbf S)=g(\Phi(\mathbf S))\). The preceding identification and the assumed integrability of \(g\) give
\[
\mathbb E_{\mathsf Q}|\varphi(\mathbf S)|
=\mathbb E_{\mathrm{ideal}}|g(N',N)|<\infty.
\]
% lean: CausalSmith.Stat.DiscreteBudgetvalueCurve.CappedMarkedCountIntegralIdentity

\item For \(0\leq M\leq n\), define the marked prefix
\[
\mathbf S_M(O^{(n)},\sigma,\omega)
=\bigl(M,((O_{\sigma(i)},\omega_i))_{i=1}^M\bigr).
\]
A permutation of an independent \(P\)-sample remains an independent \(P\)-sample, while a uniform mark vector gives each mark probability \(1/2\). Consequently, at every \(M\) whose Poisson probability is positive,
\[
\mathbb E_{O^{(n)}}\!\left[
 \frac1{n!\,2^n}\sum_{\sigma}\sum_{\omega}
 \varphi\bigl(\mathbf S_M(O^{(n)},\sigma,\omega)\bigr)
 \right]
=
\mathbb E\!\left[
 \varphi\bigl(M,((O_i^\circ,\chi_i))_{i=1}^M\bigr)\right].
\]
Integrability under \(\mathsf Q\) supplies integrability on each such count fibre; fibres with zero Poisson probability contribute zero. Multiplying by the respective Poisson probabilities and summing over \(M\leq n\) therefore yields
\[
\mathbb E_{O^{(n)}}\!\left[
 \sum_{M=0}^n\sum_{\sigma}\sum_{\omega}
 \frac{\Pr\{\operatorname{Pois}(n/4)=M\}}{n!\,2^n}\,
 \varphi\bigl(\mathbf S_M(O^{(n)},\sigma,\omega)\bigr)
 \right]
=
\mathbb E_{\mathsf Q}\!\left[
 \mathbf1\{M\leq n\}\varphi(\mathbf S)\right].
\]
This is the finite count-fibre decomposition of the marked-Poisson expectation.
% lean: Causalean.Mathlib.Probability.Poisson.FinitePartition.CappedPrefix.capped_marked_prefix_integral

\item For every admissible prefix, the indicator sums defining \(\Phi\) agree coordinatewise with the permuted pilot and evaluation counts:
\[
\Phi\bigl(\mathbf S_M(O^{(n)},\sigma,\omega)\bigr)
=(N',N)(O^{(n)},\sigma,M,\omega).
\]
Substitute this identity and \(\varphi=g\circ\Phi\) into the preceding equality. Its left side has weight \(w_n(M)\), and its right side is the stated nonoverflow expectation under the uncapped marked-Poisson experiment. When \(d=0\), there is no probability law on the empty observed alphabet, so the assertion quantified over \(P\) is vacuous.
% lean: CausalSmith.Stat.DiscreteBudgetvalueCurve.permutedMarkedCountTable_eq_uncapped_prefix
\end{enumerate}
\end{proof}

The identity in \cref{obj:lem:capped-marked-count-identity} supplies the count-law link for the conditional average. Together with the radius, approximation, bad-pilot, and maximal-path bounds, it supports a common constant for the process criterion and the threshold-process bounds.

\begin{theoremv}{thm:integrated-process-certificate}[Uniform integrated process bounds]
Fix an overlap parameter \(0<\epsilon<1/2\). There is a single constant \(C_\epsilon>0\), belonging to the uniform process handle in \cref{obj:def:process-handle}, for which the following conclusions hold whenever:

\begin{itemize}
\item \textbf{(Sample and alphabet.)} \(n\geq1\), \(d\geq16\), and \(d/\log(ed)\leq n\).
\item \textbf{(Sampling.)} The law \(\mu\) of the \(n\) observed records satisfies \cref{obj:ass:iid-sampling} for the observed margin of \(P\).
\item \textbf{(Causal law.)} The full-data law \(P\) belongs to \(\mathcal M_{d,\epsilon}\) of \cref{obj:def:model-class}.
\end{itemize}

Write \(L=\log(ed)\), \(m=n/8\), and \(v_j=\sqrt{p_jL/m}+L/m\). Simultaneously for every cell \(j\) and every \(\lambda\in[0,1]\), the absolute bias of the uncapped Jackson cell estimate of \(f_\lambda(q_j)\) is at most
\[
C_\epsilon\left\{\sqrt{\frac{p_j}{nL}}+\frac{1}{nL}\right\}.
\]
The expected squared supremum of the centered good-pilot error sum is at most \(C_\epsilon d/(nL)\). The uncentered bad-pilot cell envelope satisfies
\[
\left\|\mathbf 1_{G_j^c}\sup_{\lambda\in[0,1]}
\left|T_j(\lambda)-f_\lambda(q_j)\right|\right\|_{L_2}
\leq C_\epsilon d^{1/4}e^{-10L}v_j,
\]
and the expected squared supremum of the centered bad-pilot error sum is at most \(C_\epsilon d/(mL)\). Finally, the ideal and averaged versions of the dual threshold process in \cref{obj:def:jf-estimator} each have expected squared supremum error relative to \(F_P\) at most \(C_\epsilon d/(nL)\). The same \(C_\epsilon\) applies to all these bounds.
\end{theoremv}

The final process bounds in \cref{obj:thm:integrated-process-certificate} apply to an entire interval of shadow prices. Through the dual representation in \cref{obj:prop:dual-representation}, they provide the simultaneous control used in \cref{obj:thm:curve-upper}, including budgets selected after the curve is estimated. The average is formed at the threshold-process level before optimization.

Two boundedness facts establish that the empirical and conditionally averaged threshold processes have finite values throughout the shadow-price interval.

\begin{lemmav}{lem:empirical-threshold-process-bounded}[Bounded empirical threshold process]
For any \(n,d\in\mathbb N\), any real overlap parameter \(\epsilon\), and any sample \(O^{(n)}\) of \(n\) observations on the \(d\)-cell alphabet, let
\[
\widehat F(\lambda)=\sum_{j=1}^{d}f_\lambda(\widehat q_j)
\]
be the empirical threshold process of \cref{obj:def:jf-estimator}, with \(f_\lambda\) as in \cref{obj:def:dual-process}. There exists \(B\in\mathbb R\) such that
\[
|\widehat F(\lambda)|\leq B
\qquad\text{for every }\lambda\in[0,1].
\]
\end{lemmav}

\begin{proof}[Proof of \cref{obj:lem:empirical-threshold-process-bounded}]
Fix \(n,d,\epsilon\) and the sample. For each cell \(j\), \cref{obj:def:dual-process,obj:def:jf-estimator} give, for \(\lambda\in[0,1]\),
\[
f_\lambda(\widehat q_j)
=
g_0(\widehat q_j)
+
\max\!\left\{0,\,
g_1(\widehat q_j)-g_0(\widehat q_j)
-\lambda s(\widehat q_j)\right\}.
\]
The empirical vector \(\widehat q_j\) is fixed as \(\lambda\) varies. Thus this expression is continuous in \(\lambda\), and so is its finite sum \(\widehat F(\lambda)\). When \(n=0\), the empirical vectors are zero under the convention \(x/0=0\); when \(d=0\), the sum is empty and \(\widehat F\) is identically zero. The same continuity argument covers both cases. % lean: CausalSmith.Stat.DiscreteBudgetvalueCurve.thresholdFunReal_continuous_time

Since \([0,1]\) is compact, the continuous real-valued function \(\widehat F\) is bounded there. Hence some \(B\in\mathbb R\) satisfies
\[
|\widehat F(\lambda)|\leq B
\qquad\text{for every }\lambda\in[0,1].
\]
% lean: CausalSmith.Stat.DiscreteBudgetvalueCurve.empiricalThresholdProcess_bounded
\end{proof}

Conditional averaging over the auxiliary Poisson count preserves boundedness for the threshold process used by the reported estimator.

\begin{lemmav}{lem:averaged-threshold-process-bounded}[Bounded averaged threshold process]
Fix any observed sample \(O^{(n)}\) of \(n\) records from a \(d\)-cell alphabet. Suppose:
\begin{itemize}
\item \textbf{(Overlap parameter.)} \(\epsilon>0\).
\item \textbf{(Sample size.)} \(n\geq 1\).
\item \textbf{(Alphabet size.)} \(d\geq 1\).
\end{itemize}
The conditionally averaged threshold process \(\widehat F(\lambda)\) of \cref{obj:def:jf-estimator}, whose Poisson\((n/4)\) weights are summed over \(M=0,\ldots,n\), is bounded on \([0,1]\): there exists \(B\in\mathbb R\) such that
\[
|\widehat F(\lambda)|\leq B
\qquad\text{for every }\lambda\in[0,1].
\]
\end{lemmav}

\begin{proof}[Proof of \cref{obj:lem:averaged-threshold-process-bounded}]
\noindent\textbf{1. Continuity of the cell threshold.}
For \(u,u'\in\mathbb R_+^4\), write
\[
s_a(u)=u_{a0}+u_{a1}\quad(a\in\{0,1\}),
\qquad
s(u)=s_0(u)+s_1(u)=\sum_{a,y\in\{0,1\}}u_{ay}.
\]
Assign the quotient in \cref{obj:def:dual-process} the value \(g_a(0)=0\) at its zero denominator. The bounds below show that this is its continuous extension on the nonnegative cone. Since \(u\) is nonnegative, \(s(u)=0\) implies \(u=0\). When \(s(u)>0\), its two branches are \(s(u)u_{a1}/s_a(u)\) and \(u_{a1}/\epsilon\), according as \(s_a(u)\geq\epsilon s(u)\) or \(s_a(u)<\epsilon s(u)\). On each branch, its absolute coordinate derivatives are at most \(1+\epsilon^{-1}\); at zero, \(0\leq g_a(u)\leq s(u)\). Comparing along the segment between two nonnegative vectors therefore gives, for \(a\in\{0,1\}\),
\[
|g_a(u)-g_a(u')|
\leq(1+\epsilon^{-1})\sum_{\zeta\in\mathcal Z}|u_\zeta-u'_\zeta|.
\]
The positive-part map is \(1\)-Lipschitz. Thus, for \(\lambda\in[0,1]\),
\[
\begin{aligned}
|f_\lambda(u)-f_\lambda(u')|
&\leq 2|g_0(u)-g_0(u')|
   +|g_1(u)-g_1(u')|
   +\lambda|s(u)-s(u')|\\
&\leq \bigl[3(1+\epsilon^{-1})+1\bigr]
       \sum_{\zeta\in\mathcal Z}|u_\zeta-u'_\zeta|.
\end{aligned}
\]
Hence \(f_\lambda\) is continuous on the nonnegative cone for every fixed \(\lambda\in[0,1]\).
% lean: CausalSmith.Stat.DiscreteBudgetvalueCurve.thresholdFunReal_continuousOn_nonnegative

\noindent\textbf{2. The pilot rectangle.}
Fix any pilot counts and cell \(j\). The quantities in \cref{obj:def:jf-estimator} satisfy
\[
m=\frac n8>0,\qquad L=\log(ed)\geq1,\qquad
\delta=\frac Lm>0,\qquad h_{\zeta j}>0.
\]
For every \(\boldsymbol z\in[-1,1]^{\mathcal Z}\) and \(\zeta\in\mathcal Z\),
\[
0\leq\ell_{\zeta j}
\leq c^0_{\zeta j}+R_{\zeta j}z_\zeta
\leq u_{\zeta j}.
\]
The map \(\boldsymbol z\mapsto c_j^0+R_j\odot\boldsymbol z\) is affine and continuous. Since \(f_\lambda\) is continuous on the nonnegative cone, its composition with this map is continuous on \([-1,1]^{\mathcal Z}\) for every \(\lambda\in[0,1]\). This includes pilot counts arising when \(M=0\).
% lean: CausalSmith.Stat.DiscreteBudgetvalueCurve.pilotRectangle_threshold_pullback_continuousOn

\noindent\textbf{3. Continuity of each clipped cell estimate.}
Keep the counts fixed. The Jackson order and degree in \cref{obj:def:jf-estimator} satisfy \(K\geq2\) and \(D=2(K-1)\geq2\). For \(r=0,\ldots,D\), set
\[
\xi_r=\frac{r}{D+1},
\qquad
\boldsymbol\xi_\gamma=(\xi_{\gamma_\zeta})_{\zeta\in\mathcal Z},
\qquad
\vartheta_{\gamma,\zeta}=\arccos(\xi_{\gamma_\zeta}),
\quad
\gamma\in\{0,\ldots,D\}^{\mathcal Z}.
\]
At a fixed grid point, the defining Jackson integral and the threshold formula give
\[
\begin{aligned}
&\left|
P_{j,\lambda}(c_j^0+R_j\odot\boldsymbol\xi_\gamma)
-P_{j,\mu}(c_j^0+R_j\odot\boldsymbol\xi_\gamma)
\right|\\
&\qquad\leq |\lambda-\mu|
\int_{[-\pi,\pi]^{\mathcal Z}}
s\!\left(c_j^0+R_j\odot
\cos(\boldsymbol\vartheta_\gamma+\boldsymbol t)\right)
\prod_{\zeta\in\mathcal Z}J_K(t_\zeta)\,d\boldsymbol t,
\qquad \lambda,\mu\in[0,1],
\end{aligned}
\]
where cosine acts coordinatewise. The integral is finite: its argument lies in the pilot rectangle, and the integrand is continuous on a compact box. Also,
\[
|f_\lambda(c_j^0)-f_\mu(c_j^0)|
\leq |\lambda-\mu|\,s(c_j^0).
\]

For \(\alpha,\gamma\in\{0,\ldots,D\}^{\mathcal Z}\), define the fixed tensor interpolation weights by
\[
\boldsymbol z^\alpha:=\prod_{\zeta\in\mathcal Z}z_\zeta^{\alpha_\zeta},
\qquad
w_{\alpha,\gamma}:=
[\boldsymbol z^\alpha]\,
\prod_{\zeta\in\mathcal Z}
\prod_{\substack{0\leq r\leq D\\r\ne\gamma_\zeta}}
\frac{z_\zeta-\xi_r}{\xi_{\gamma_\zeta}-\xi_r},
\]
where \([\boldsymbol z^\alpha]\) extracts the indicated coefficient. Tensor Lagrange interpolation of the degree-\(D\) polynomial yields
\[
\beta_{j,\alpha}(\lambda)
=
\sum_{\gamma\in\{0,\ldots,D\}^{\mathcal Z}}
w_{\alpha,\gamma}
\left[
P_{j,\lambda}(c_j^0+R_j\odot\boldsymbol\xi_\gamma)
-f_\lambda(c_j^0)
\right].
\]
The preceding bounds show that every \(\beta_{j,\alpha}(\lambda)\) is continuous in \(\lambda\). All count-dependent factors in
\[
Z_j(\lambda)=f_\lambda(c_j^0)+
\sum_{\alpha\in\{0,\ldots,D\}^{\mathcal Z}}
\beta_{j,\alpha}(\lambda)
\prod_{\zeta\in\mathcal Z}
\frac{U_{\alpha_\zeta}(N_{\zeta j};c^0_{\zeta j})}
     {R_{\zeta j}^{\alpha_\zeta}}
\]
are fixed and have positive radii in their denominators. Thus \(Z_j\) is continuous. Finally,
\[
T_j(\lambda)=f_\lambda(c_j^0)+
\min\!\left\{d^{1/4}S_j,\,
\max\!\left\{-d^{1/4}S_j,\,
Z_j(\lambda)-f_\lambda(c_j^0)\right\}\right\}
\]
is continuous, since finite minima and maxima preserve continuity.
% lean: CausalSmith.Stat.DiscreteBudgetvalueCurve.jacksonCellStatistic_continuous_time

\noindent\textbf{4. Averaging and boundedness.}
Write \(\mathfrak S_n\) for the permutations of the \(n\) record positions. Conditional on the fixed sample, the average is the finite sum
\[
\widehat F(\lambda)
=
\sum_{M=0}^{n}
\frac{e^{-n/4}(n/4)^M}{M!}\,
\frac{1}{n!\,2^n}
\sum_{\sigma\in\mathfrak S_n}
\sum_{\omega\in\{0,1\}^n}
\left.\left(\sum_{j=1}^{d}T_j(\lambda)\right)
\right|_{M,\sigma,\omega}.
\]
Only the first \(M\) marks affect the count table; each marking of those retained records occurs \(2^{n-M}\) times in the displayed sum. Thus the \(2^n\) normalization gives the same average as \(2^M\) markings of the retained records. The vertical bar means that the counts are computed from that permutation and marking. Each cell summand \(T_j(\lambda)\) is continuous in \(\lambda\). The capped value for \(M>n\) is the constant zero. Hence \(\widehat F\) is continuous on \([0,1]\). Compactness of that interval gives the finite real number
\[
B=\max_{\lambda\in[0,1]}|\widehat F(\lambda)|,
\qquad
|\widehat F(\lambda)|\leq B
\quad\text{for every }\lambda\in[0,1].
\]
% lean: CausalSmith.Stat.DiscreteBudgetvalueCurve.averagedThresholdProcess_bounded
\end{proof}

For the lower bound, a simplex records the masses assigned to pairs of covariate cells. The paired family then isolates the effect of allocating a fixed treatment capacity across unequal benefits.

\begin{definitionv}{def:paired-family}[Paired causal family $\mathcal P_k^{\mathrm{pair}}$]
The paired causal family \(\mathcal P_k^{\mathrm{pair}}\) consists of causal completions constructed as follows. For \(k\geq1\), choose \(r=(r_1,\ldots,r_k)\) in the simplex and \(\theta_i\in[-1/8,1/8]\). For each \(i\), create cells \((i,+)\) and \((i,-)\) with masses \(r_i/2\), and set
\[
e_{i,\sigma}=\frac12,\qquad
\mu_{0,i,\sigma}=\frac14,\qquad
\mu_{1,i,\sigma}=\frac12+\sigma\theta_i,
\qquad
\sigma\in\{+1,-1\}.
\]
Generate the potential outcomes independently of \(A\) conditional on \(X\).
\end{definitionv}

The simplex \(\Delta_k\) in \cref{obj:synth_10} supplies the pair-mass vector for \cref{obj:def:paired-family}.

\begin{definitionv}{synth_10}[Probability simplex on \(k\) categories]
For an integer \(k\geq1\), define
\[
\Delta_k=\left\{r=(r_1,\ldots,r_k)\in[0,1]^k:\sum_{i=1}^k r_i=1\right\}.
\]
\end{definitionv}

With the pair masses specified, the causal family fixes the propensity and control mean while varying the treatment mean within each pair. Treatment effects remain positive while their ranking depends on the contrast. The capacity and sampling identities make this allocation problem explicit.

\begin{propositionv}{prop:paired-capacity-identity}[Paired capacity and sample identity]
Let \(k\geq 1\) and \(0<\epsilon<1/2\). The following statements hold.
\begin{itemize}
\item \textbf{(Capacity identity.)} For every \(P\in\mathcal P^{\mathrm{pair}}_k\) and every pair-mass vector \(r\in\Delta_k\) and contrast vector \(\theta\in[-1/8,1/8]^k\) for which \(P\) satisfies the paired-completion conditions of \cref{obj:def:paired-family}, the law \(P\) belongs to \(\mathcal M_{2k,\epsilon}\) of \cref{obj:def:model-class}. Every occupied cell has treatment effect \(\tau_j\in[1/8,3/8]\). The half-population budget is attained by a policy \(\pi\in\Pi_{1/2}(P)\) satisfying
\[
\sum_j p_j\pi_j=\frac12,
\qquad
V_{1/2}(P)=B(P)+\sum_j p_j\pi_j\tau_j.
\]
For the budget values of \cref{obj:def:budget-value},
\[
V_{1/2}(P)
=\frac38+\frac12\sum_{i=1}^k r_i|\theta_i|
=\frac18+F_P\!\left(\frac14\right),
\qquad
V_1(P)=\frac12.
\]

\item \textbf{(Exact sample reduction.)} For every integer \(n\geq0\), there is a Markov kernel from pairs of \(n\)-record samples on the \(k\)-point alphabet to \(n\) observed records on the \(2k\)-cell alphabet with the following property. For every pair \(R,S\in\Delta_k\), the simplex of \(k\)-point laws, set
\[
r_i=\frac{R_i+S_i}{2},
\qquad
\theta_i=
\begin{cases}
0, & R_i+S_i=0,\\[2pt]
\dfrac{R_i-S_i}{8(R_i+S_i)}, & R_i+S_i>0.
\end{cases}
\]
The paired law determined by \((r,\theta)\) belongs to \(\mathcal P^{\mathrm{pair}}_k\). Applying the kernel to independent \(n\)-record samples from \(R\) and \(S\) yields exactly the law of \(n\) independent observed records from the observed margin of that paired law.
\end{itemize}
\end{propositionv}

For two sampling laws \leanref{sym:R}{\ensuremath{R}} and \leanref{sym:S}{\ensuremath{S}}, the exact sample reduction gives a half-budget value of \(3/8+\|R-S\|_1/32\), while the full-budget value is \(1/2\). Thus the half-budget comparison carries a two-sample distance problem within the causal model. The statistical lower-bound input is a squared-error bound for the two-sample distance.

The following bound states the two-sample risk at the scale used by the reduction.

\begin{lemmav}{lem:two-sample-l1-lower}[Two-sample \(L_1\) lower bound]*
There is a universal constant \(c>0\) such that, for every integer \(k\geq 2\) and \(1\leq n\leq k^2\), the minimax squared-error risk for estimating \(\|R-S\|_1\) from \(n\) independent observations from each of \(R,S\in\Delta_k\), the probability simplex on \(k\) categories, satisfies
\[
\inf_{\widehat L}\sup_{R,S\in\Delta_k}
E_{R,S}\!\left[(\widehat L-\|R-S\|_1)^2\right]
\geq c\min\left\{1,\frac{k}{n\log(ek)}\right\}.
\]
\end{lemmav}
% CAUSALSMITH-CITED-SCOPE-BEGIN lem:two-sample-l1-lower
\verificationfootnotetext{\textbf{Formalization scope.} This result uses the published conclusion \citep{JiaoHanWeissman2018} (Theorem 3 and the unknown-both-distributions reduction following Theorem 6). The cited source proof is not formalized here: Lean verifies its use and all remaining steps. If that cited conclusion is false, the portions of this result that depend on it are not certified.}
% CAUSALSMITH-CITED-SCOPE-END lem:two-sample-l1-lower

\begin{proof}[Proof of \cref{obj:lem:two-sample-l1-lower}]
Write
\[
\mathcal R_{n,k}:=
\inf_{\widehat L}\sup_{R,S\in\Delta_k}
E_{R,S}\!\left[(\widehat L-\|R-S\|_1)^2\right].
\] % lean: CausalSmith.Stat.DiscreteBudgetvalueCurve.two_sample_l1_lower

We first transfer the published known-\(S\) bound to this two-sample experiment. Given a two-sample estimator \(\widehat L\) and a fixed \(S\), average over its second sample:
\[
\overline L_S(X^n):=
E_{Y^n\sim S^{\otimes n}}\!\left[\widehat L(X^n,Y^n)\right].
\]
Conditional Jensen inequality gives, for every \(R\),
\[
E_R\!\left[(\overline L_S-\|R-S\|_1)^2\right]
\leq E_{R,S}\!\left[(\widehat L-\|R-S\|_1)^2\right].
\]
Taking the relevant suprema and infima therefore yields
\[
\sup_{S\in\Delta_k}\inf_{\widetilde L_S}\sup_{R\in\Delta_k}
E_R\!\left[(\widetilde L_S-\|R-S\|_1)^2\right]
\leq \mathcal R_{n,k}.
\] % lean: CausalSmith.Stat.DiscreteBudgetvalueCurve.twoSampleL1_known_minimax_le

The published bound \citep[Theorem 3 and the unknown-both-distributions reduction following Theorem 6]{JiaoHanWeissman2018}, applied with \(a_0=1\) and \(C_0=2\), consequently supplies constants \(\gamma_{\mathrm{pub}}>0\) and \(k_0\) such that
\[
\mathcal R_{n,k}\geq \gamma_{\mathrm{pub}}\,\frac{k}{n\log n}
\quad\text{if}\quad
k\geq k_0,\qquad n\geq\frac{k}{\log k},
\qquad \log n\leq2\log k.
\]
For \(n\geq2\), the last condition and \(\log(ek)=1+\log k\) imply
\[
\mathcal R_{n,k}\geq
\frac{\gamma_{\mathrm{pub}}}{2}\,\frac{k}{n\log(ek)}.
\] % lean: CausalSmith.Stat.DiscreteBudgetvalueCurve.twoSampleL1_lower_dense_logAlphabet

In particular, enlarging a threshold \(k_D\) if necessary, there is a constant \(\gamma_D>0\) such that, whenever \(k\geq k_D\), \(k\geq2\), \(2\leq n\leq k^2\), and \(n\geq k/\log k\),
\[
\mathcal R_{n,k}\geq
\gamma_D\min\left\{1,\frac{k}{n\log(ek)}\right\}.
\]
Indeed, \(n\leq k^2\) gives \(\log n\leq2\log k\), and the minimum is at most its second argument. % lean: CausalSmith.Stat.DiscreteBudgetvalueCurve.twoSampleL1_lower_large_dense

For the smaller sample sizes, use the same dense bound at the integer
\[
n_\circ(k):=\left\lceil\frac{k}{\log k}\right\rceil.
\]
When \(k\geq3\), the inequalities \(1\leq\log k\leq k-1\) show that \(k/\log k>1\), while
\[
\frac{k}{\log k}\leq n_\circ(k)
<\frac{k}{\log k}+1
\leq\frac{2k}{\log k},
\qquad
n_\circ(k)\leq k^2,
\qquad
\log n_\circ(k)\leq2\log k.
\]
Also \(\log(ek)\leq2\log k\), so \(n_\circ(k)\log(ek)\leq4k\). Applying the dense bound at \(n_\circ(k)\) thus gives, for some \(\gamma_S>0\) and threshold \(k_S\),
\[
\mathcal R_{n_\circ(k),k}\geq\gamma_S
\quad (k\geq\max\{k_S,3\}).
\]
If \(1\leq n<k/\log k\), then \(n\leq n_\circ(k)\). An estimator with \(n_\circ(k)\) observations can use their first \(n\) coordinates, whose law is exactly the \(n\)-sample law. Hence
\[
\mathcal R_{n,k}\geq\mathcal R_{n_\circ(k),k}
\geq\gamma_S.
\] % lean: CausalSmith.Stat.DiscreteBudgetvalueCurve.twoSampleL1_lower_large_sparse

Set \(k_\star:=\max\{k_D,k_S,3\}\). To cover \(2\leq k<k_\star\), consider the two parameter pairs
\[
(R_0,S)=((1,0,\ldots,0),(1,0,\ldots,0)),
\qquad
(R_1,S)=((1/2,1/2,0,\ldots,0),(1,0,\ldots,0)).
\]
Their targets are \(0\) and \(1\). Let \(\omega_0\) be the data point at which every observation in both samples has the first category, and put \(\rho_n:=2^{-n}\). The probabilities of \(\{\omega_0\}\) under the two pairs are \(1\) and \(\rho_n\), respectively. Thus, for every estimator \(\widehat L\), its worst-case risk satisfies
\[
\sup_{R,S\in\Delta_k}E_{R,S}\!\left[(\widehat L-\|R-S\|_1)^2\right]
\geq
\max\!\left\{\widehat L(\omega_0)^2,\,
\rho_n\bigl(\widehat L(\omega_0)-1\bigr)^2\right\}
\geq \frac{\rho_n}{4},
\]
since \(\rho_n\leq1\) and \(x^2+(x-1)^2\geq1/2\) for every real \(x\). Taking the infimum over estimators gives
\[
\mathcal R_{n,k}\geq\frac{\rho_n}{4}.
\] % lean: CausalSmith.Stat.DiscreteBudgetvalueCurve.twoSampleL1_finiteAlphabetLower

Minimax risk decreases as the sample size increases by the same first-coordinate argument. Since \(n\leq k^2\leq k_\star^2\), the preceding bound at \(k_\star^2\) gives
\[
\mathcal R_{n,k}\geq
\mathcal R_{k_\star^2,k}\geq
\frac{2^{-k_\star^2}}4
\qquad(2\leq k<k_\star).
\] % lean: CausalSmith.Stat.DiscreteBudgetvalueCurve.twoSampleL1_minimax_antitone

Finally, choose
\[
c:=\min\left\{\gamma_D,\gamma_S,\frac{2^{-k_\star^2}}4\right\}>0.
\]
For \(k\geq k_\star\), split according as \(n<k/\log k\) or \(n\geq k/\log k\); the sparse and dense bounds, respectively, give the claim. For \(2\leq k<k_\star\), use the finite-alphabet bound. In each case the claimed minimum is at most \(1\). % lean: CausalSmith.Stat.DiscreteBudgetvalueCurve.two_sample_l1_lower
\end{proof}

The published intermediate-regime lower-bound input for \cref{obj:lem:two-sample-l1-lower} is given by \citet[Theorem 3 and the unknown-both-distributions reduction following Theorem 6]{JiaoHanWeissman2018}. Combined with the exact sample reduction in \cref{obj:prop:paired-capacity-identity}, the lemma supplies the half-budget lower bound in \cref{obj:thm:capacity-active-lower}. That bound and the upper bound in \cref{obj:thm:curve-upper} yield the matched risk statement in \cref{obj:thm:matched-curve-frontier} under its stated conditions.

The theorem statements and internal proofs are machine-checked in Lean 4. The sampling and causal conditions in \cref{obj:ass:iid-sampling,obj:ass:consistency,obj:ass:conditional-exchangeability,obj:ass:overlap} are model assumptions. The published conclusion cited for \cref{obj:lem:two-sample-l1-lower} is a mathematical input from the literature.


\section{Proofs of the main results}\label{sec:deferred-proofs}

\begin{proof}[Proof of \cref{obj:prop:dual-representation}]
Write
\[
p_j=P(X=j),\qquad q_{(a,y)j}=P(X=j,A=a,Y=y),\qquad
\mu_{aj}=E_P[Y(a)\mid X=j],\qquad
\tau_j=\mu_{1j}-\mu_{0j},\qquad
B(P)=\sum_{j=1}^d p_j\mu_{0j}.
\]
For a cell with \(p_j=0\), set \(\mu_{aj}=0\); its contributions below vanish. The cell masses are nonnegative, and the binary outcomes give \(\mu_{aj}\in[0,1]\), hence \(\tau_j\leq1\). % lean: CausalSmith.Stat.DiscreteBudgetvalueCurve.poRegression_unit_interval_budget

For an occupied cell, \cref{obj:ass:consistency,obj:ass:conditional-exchangeability,obj:ass:overlap} give \(s_a(q_j)=P(X=j,A=a)\geq\epsilon p_j>0\) and
\[
g_a(q_j)
=p_j\frac{P(X=j,A=a,Y=1)}{P(X=j,A=a)}
=p_j\mu_{aj}.
\]
For a null cell, \(q_j=0\) and the same identity has both sides zero. Substituting these identities into \cref{obj:def:dual-process}, and using \(p_j\geq0\), yields, for \(0\leq\lambda\leq1\),
\[
f_\lambda(q_j)=p_j\mu_{0j}+p_j[\tau_j-\lambda]_+,
\qquad
F_P(\lambda)=B(P)+\sum_{j=1}^d p_j[\tau_j-\lambda]_+.
\] % lean: CausalSmith.Stat.DiscreteBudgetvalueCurve.budget_dualProcess_identified

Fix \(b\in[0,1]\). If \(\pi\in\Pi_b(P)\), then \(0\leq\pi_j\leq1\) and \(\sum_jp_j\pi_j\leq b\) by \cref{obj:def:budget-policy-class}. For \(\lambda\in[0,1]\), the inequality
\(p_j\pi_j(\tau_j-\lambda)\leq p_j[\tau_j-\lambda]_+\)
holds whether \(\tau_j-\lambda\) is nonnegative or negative. Since \(\lambda\geq0\),
\[
\sum_{j=1}^d p_j\pi_j\tau_j
=\lambda\sum_{j=1}^d p_j\pi_j
 +\sum_{j=1}^d p_j\pi_j(\tau_j-\lambda)
\leq b\lambda+\sum_{j=1}^d p_j[\tau_j-\lambda]_+.
\] % lean: CausalSmith.Stat.DiscreteBudgetvalueCurve.weighted_budget_weak_duality
Taking the maximum over feasible policies and using \cref{obj:def:budget-value} gives
\[
V_b(P)\leq b\lambda+F_P(\lambda)
\quad(0\leq\lambda\leq1),
\qquad
V_b(P)\leq\inf_{0\leq\lambda\leq1}\{b\lambda+F_P(\lambda)\}.
\] % lean: CausalSmith.Stat.DiscreteBudgetvalueCurve.budget_dual_weak_bound

For the reverse inequality, the feasible policy set is a compact subset of \([0,1]^d\), contains the zero policy, and has a continuous gain. Thus a feasible maximizer \(\pi^\star\) exists. Applied to the affine budget and box constraints, the polyhedral normal-cone formula gives nonnegative multipliers \(\lambda^\star,\eta_j^-,\eta_j^+\) satisfying
\[
\begin{gathered}
\sum_jp_j\pi^\star_j\leq b,\qquad
\lambda^\star\Bigl(b-\sum_jp_j\pi^\star_j\Bigr)=0,\\
\eta_j^-\pi^\star_j=0,\qquad
\eta_j^+(1-\pi^\star_j)=0,\qquad
p_j(\lambda^\star-\tau_j)-\eta_j^-+\eta_j^+=0.
\end{gathered}
\] % lean: Causalean.Mathlib.Optimization.finite_knapsack_kkt
These conditions give the hinge identity cell by cell. It is immediate when \(p_j=0\). When \(p_j>0\), a coordinate with \(\pi^\star_j=0\) has \(\eta_j^+=0\) and \(\tau_j\leq\lambda^\star\); one with \(\pi^\star_j=1\) has \(\eta_j^-=0\) and \(\tau_j\geq\lambda^\star\); an interior coordinate has both box multipliers zero and \(\tau_j=\lambda^\star\). Consequently,
\[
p_j\pi^\star_j\tau_j
=\lambda^\star p_j\pi^\star_j
 +p_j[\tau_j-\lambda^\star]_+,
\qquad
\sum_jp_j\pi^\star_j\tau_j
=b\lambda^\star+\sum_jp_j[\tau_j-\lambda^\star]_+,
\]
where the second equality uses budget complementary slackness. % lean: Causalean.Mathlib.Optimization.finite_knapsack_dual_witness

If \(\lambda^\star\leq1\), this is a witness with price in \([0,1]\). If \(\lambda^\star>1\), then \(\tau_j\leq1\) makes every hinge at \(\lambda^\star\) zero. Also,
\[
b\lambda^\star
=\sum_jp_j\pi^\star_j\tau_j
\leq\sum_jp_j\pi^\star_j
\leq b.
\]
As \(b\geq0\), this forces \(b=0\). Every hinge at price \(1\) is then zero as well, and the displayed witness equality holds with price \(1\). Hence in every case there is \(\lambda_b\in[0,1]\) such that
\[
b\lambda_b+F_P(\lambda_b)
=B(P)+\sum_jp_j\pi^\star_j\tau_j
=V_b(P).
\] % lean: Causalean.Mathlib.Optimization.finite_knapsack_dual_witness
The process \(F_P\) is a finite sum of continuous hinge functions on \([0,1]\), so it is bounded there. Thus the dual infimum is well defined and is at most its value at \(\lambda_b\). Together with weak duality, this proves the representation. % lean: CausalSmith.Stat.DiscreteBudgetvalueCurve.dual_representation

Finally, define
\[
\omega_{\widehat F,F_P}=\sup_{\lambda\in[0,1]}|\widehat F(\lambda)-F_P(\lambda)|.
\]
Both processes are bounded on \([0,1]\), so \(\omega_{\widehat F,F_P}\) is finite and both dual objectives are bounded below for \(b\in[0,1]\). For every such \(b\) and every \(\lambda\in[0,1]\), the two inequalities
\[
b\lambda+\widehat F(\lambda)\leq b\lambda+F_P(\lambda)+\omega_{\widehat F,F_P},
\qquad
b\lambda+F_P(\lambda)\leq b\lambda+\widehat F(\lambda)+\omega_{\widehat F,F_P}
\]
give, on taking infima,
\[
\left|
\inf_{0\leq\lambda\leq1}\{b\lambda+\widehat F(\lambda)\}
-\inf_{0\leq\lambda\leq1}\{b\lambda+F_P(\lambda)\}
\right|\leq \omega_{\widehat F,F_P}.
\] % lean: CausalSmith.Stat.DiscreteBudgetvalueCurve.dualMinimum_uniform_contraction
Substitute the representation just proved and take the supremum over \(b\in[0,1]\). % lean: CausalSmith.Stat.DiscreteBudgetvalueCurve.dual_representation
\end{proof}

\begin{proof}[Proof of \cref{obj:thm:matched-curve-frontier}]
Set
\[
\varrho_{n,d}:=\min\left\{1,\frac{d}{n\log(ed)}\right\}.
\]
% lean: CausalSmith.Stat.DiscreteBudgetvalueCurve.curveRate
Choose \(c_\epsilon>0\) from \cref{obj:thm:capacity-active-lower} and \(C_\epsilon^{\mathrm{up}}>0\) from \cref{obj:thm:curve-upper}, and define
\[
C_\epsilon:=C_\epsilon^{\mathrm{up}}+c_\epsilon.
\]
% lean: CausalSmith.Stat.DiscreteBudgetvalueCurve.matched_curve_frontier
The discrete \(L_1\) lower bound underlying \cref{obj:thm:capacity-active-lower} is supplied by \citep[Theorem 3 and the unknown-both-distributions reduction following Theorem 6]{JiaoHanWeissman2018}.

Fix \(n\ge1\), \(d\ge2\), and \(b_0\in[0,1/2]\), and use the independent observed-record laws of \cref{obj:ass:iid-sampling}. Since \(b_0\le1/2\le1-b_0\), evaluation of any admissible curve estimator \(\widehat V\) at \(1/2\) gives a measurable scalar estimator
\[
T_{\widehat V}(O^{(n)}):=\widehat V_{1/2}(O^{(n)}).
\]
% lean: CausalSmith.Stat.DiscreteBudgetvalueCurve.scalar_coordinate_of_curve_estimator
By \cref{obj:thm:capacity-active-lower}, there is a law \(P\in\mathcal M_{d,\epsilon}\) for which
\[
\begin{aligned}
c_\epsilon\varrho_{n,d}
&\le E_P\!\left[\bigl(T_{\widehat V}-V_{1/2}(P)\bigr)^2\right]\\
&\le E_P\!\left[\sup_{b\in[b_0,1-b_0]}
       \bigl|\widehat V_b-V_b(P)\bigr|^2\right]\\
&\le \sup_{P'\in\mathcal M_{d,\epsilon}}
 E_{P'}\!\left[\sup_{b\in[b_0,1-b_0]}
       \bigl|\widehat V_b-V_b(P')\bigr|^2\right].
\end{aligned}
\]
% lean: CausalSmith.Stat.DiscreteBudgetvalueCurve.scalar_coordinate_risk_le_curve_risk
The middle inequality holds samplewise because \(1/2\) belongs to the budget interval. The lower-bound result also supplies a law in the model class when applied to the zero scalar estimator, while \cref{obj:def:jf-estimator} supplies an admissible curve estimator. Thus taking the infimum over curve estimators in \cref{obj:def:curve-risk} gives
\[
c_\epsilon\varrho_{n,d}
\le \mathfrak R_{n,d,\epsilon,b_0}^{\mathrm{curve}}.
\]
% lean: CausalSmith.Stat.DiscreteBudgetvalueCurve.matched_curve_frontier

For the upper bound, insert the estimator of \cref{obj:def:jf-estimator} into that infimum. \Cref{obj:thm:curve-upper} bounds its worst-case curve risk, so
\[
\begin{aligned}
\mathfrak R_{n,d,\epsilon,b_0}^{\mathrm{curve}}
&\le \sup_{P\in\mathcal M_{d,\epsilon}}
 E_P\!\left[\sup_{b\in[b_0,1-b_0]}
 \bigl|\widehat V_b^{\mathrm{JF}}-V_b(P)\bigr|^2\right]\\
&\le C_\epsilon^{\mathrm{up}}\varrho_{n,d}
 \le C_\epsilon\varrho_{n,d}.
\end{aligned}
\]
% lean: CausalSmith.Stat.DiscreteBudgetvalueCurve.matched_curve_frontier
Here \(\varrho_{n,d}\ge0\), and the chosen constants satisfy \(0<c_\epsilon\le C_\epsilon<\infty\).

Now let \(T\) be any measurable scalar estimator. Applying \cref{obj:thm:capacity-active-lower} to the same independent observed-record laws supplies a law \(P\in\mathcal M_{d,\epsilon}\) with \(V_1(P)=1/2\), a binding half-budget policy \(\pi\in\Pi_{1/2}(P)\), and
\[
\sum_{j=1}^d p_j\pi_j=\frac12,\qquad
V_{1/2}(P)=B(P)+\sum_{j=1}^d p_j\pi_j\tau_j,\qquad
E_P\!\left[(T-V_{1/2}(P))^2\right]\ge c_\epsilon\varrho_{n,d}.
\]
% lean: CausalSmith.Stat.DiscreteBudgetvalueCurve.capacity_active_lower
For every occupied cell, the same result gives
\[
p_j>0\quad\Longrightarrow\quad
\frac18\le\tau_j\le\frac38
\quad\Longrightarrow\quad \tau_j>0.
\]
% lean: CausalSmith.Stat.DiscreteBudgetvalueCurve.matched_curve_frontier

Finally, suppose \(b_0=0\) and take any \(P\in\mathcal M_{d,\epsilon}\). Since \(p_j\ge0\) and \(\sum_jp_j=1\), every policy in \([0,1]^d\) is feasible at budget \(1\). For any such policy,
\[
\sum_{j=1}^d p_j\pi_j\tau_j
\le \sum_{j=1}^d p_j\max\{\tau_j,0\},
\]
% lean: CausalSmith.Stat.DiscreteBudgetvalueCurve.full_budget_value
and equality is attained by the policy
\[
\pi_j^*:=\mathbf1\{\tau_j>0\},\qquad j=1,\ldots,d.
\]
% lean: CausalSmith.Stat.DiscreteBudgetvalueCurve.full_budget_value
Consequently, \cref{obj:def:budget-value} gives
\[
V_1(P)=B(P)+\sum_{j=1}^d p_j\max\{\tau_j,0\}.
\]
% lean: CausalSmith.Stat.DiscreteBudgetvalueCurve.full_budget_value
\end{proof}

\begin{proof}[Proof of \cref{obj:thm:capacity-active-lower}]
Write
\[
\mathcal R_{n,k}
=\inf_{\widehat H}\sup_{R,S\in\Delta_k}
E_{R,S}\!\left[(\widehat H-\|R-S\|_1)^2\right].
\] % lean: CausalSmith.Stat.DiscreteBudgetvalueCurve.twoSampleL1MinimaxRisk
Here \(\widehat H\) ranges over measurable estimators based on \(n\) observations from each of \(R\) and \(S\). We first establish the two-sample bound used in the proof:
\[
\mathcal R_{n,k}\ge c_{\mathrm L}
\min\left\{1,\frac{k}{n\log(ek)}\right\}
\qquad(k\ge2,\ n\ge1)
\]
for a universal \(c_{\mathrm L}>0\). % lean: CausalSmith.Stat.DiscreteBudgetvalueCurve.two_sample_l1_lower_all_samples

For \(n\le k^2\), this is \cref{obj:lem:two-sample-l1-lower}. Its dense-range input is the published known-\(S\) bound: when \(n\ge a_0k/\log k\) and \(\log n\le C_0\log k\), the worst-case squared risk is at least a positive constant times \(k/(n\log n)\), with constants depending on \(a_0,C_0\) \citep[Theorem 3 and the unknown-both-distributions reduction following Theorem 6]{JiaoHanWeissman2018}. % lean: CausalSmith.Stat.DiscreteBudgetvalueCurve.two_sample_l1_lower

For completeness, consider the remaining range \(n>k^2\). At sufficiently large \(k\), set
\[
h=\lfloor k/2\rfloor,\qquad
\ell=16\bigl(\lfloor\log_2 k\rfloor+1\bigr),\qquad
\xi=\sqrt{\frac{h\ell}{200n}}>0,\qquad
\omega=\frac{\xi}{50\ell}.
\] % lean: Causalean.Stat.Minimax.Multinomial.TwoSampleL1.exists_pairedLargeParameters
To obtain the finite priors used here, let
\[
\mathcal E_\ell=\inf_{\deg\varphi\le\ell}\sup_{|x|\le1}\bigl||x|-\varphi(x)\bigr|.
\]
Polynomial alternation gives nodes \(x_0<\cdots<x_{\ell+1}\) in \([-1,1]\) and signed weights \(a_i\), normalized so that
\[
\sum_i|a_i|=1,\qquad
\sum_i a_i x_i^{\ell'}=0\quad(0\le\ell'\le\ell),\qquad
\sum_i a_i|x_i|=\mathcal E_\ell.
\]
Indeed, the weights proportional to \(\prod_{i'\ne i}(x_i-x_{i'})^{-1}\) annihilate polynomials of degree at most \(\ell\) by Lagrange interpolation; their alternating signs can be oriented to agree with the alternating approximation errors. The required approximation error has the bound
\[
\mathcal E_\ell\ge\frac1{100\ell}.
\]
For this bound, let \(\mathscr F_\ell\) denote the normalized Fejér kernel of order \(\ell\) on the circle of period \(\pi\), and define
\[
g(\phi)=|\sin(2\phi)|,\qquad
\mathscr D_\ell f=f(0)-\bigl((2\mathscr F_{2\ell}-\mathscr F_\ell)*f\bigr)(0).
\]
The Fejér kernels have integral one, so \(\|\mathscr D_\ell\|\le4\). Their Fourier multipliers give \(\mathscr D_\ell[\varphi(\sin(2\phi))]=0\) when \(\deg\varphi\le\ell\). The Fourier series of \(g\) has nonnegative contributions to the magnitude of this functional; retaining the frequencies indexed from \(\ell\) through \(2\ell-1\) yields
\[
|\mathscr D_\ell g|
\ge\sum_{\ell'=\ell}^{2\ell-1}
 \frac{4}{\pi(4{\ell'}^2-1)}
\ge\frac1{16\ell}.
\]
Thus \(4\mathcal E_\ell\ge1/(16\ell)\), which gives the displayed bound. Writing \(a_i^+=\max(a_i,0)\) and \(a_i^-=\max(-a_i,0)\), define
\[
\nu_0=2\sum_i a_i^-\delta_{x_i},\qquad
\nu_1=2\sum_i a_i^+\delta_{x_i}.
\]
The degree-zero identity makes each measure a probability measure. The remaining identities show that their moments agree through degree \(\ell\) and that
\[
\int|x|\,d\nu_1(x)-\int|x|\,d\nu_0(x)
=2\mathcal E_\ell\ge\frac1{50\ell}.
\] % lean: Causalean.Stat.Minimax.Multinomial.TwoSampleL1.exists_scalarMomentPriors

Use \(2h\) cells and, for \(z=(z_1,\ldots,z_h)\in[-1,1]^h\), define the two probability vectors by
\[
R_{i,\pm}=\frac1{2h},\qquad
S_{i,\pm}(z)=\frac{1\pm\xi z_i}{2h},
\qquad
\|R-S(z)\|_1=\frac{\xi}{h}\sum_{i=1}^h|z_i|.
\] % lean: Causalean.Stat.Minimax.Multinomial.TwoSampleL1.pairedProductTarget_mean
Under prior \(\nu_s^{\otimes h}\), the target has center
\(\gamma_s=\xi\int|z|\,d\nu_s(z)\). Independence gives
\[
|\gamma_1-\gamma_0|\ge\omega,\qquad
\operatorname{Var}_{\nu_s^{\otimes h}}\!\bigl(\|R-S(z)\|_1\bigr)
\le\frac{\xi^2}{h}.
\] % lean: Causalean.Stat.Minimax.Multinomial.TwoSampleL1.pairedProductTarget_variance_le
For the stated sufficiently large alphabet threshold, \(128\xi^2\le h\omega^2\). Chebyshev’s inequality therefore places at most \(1/8\) of either prior outside the interval of radius \(\omega/4\) around its center. % lean: Causalean.Stat.Minimax.Multinomial.TwoSampleL1.pairedProductTarget_bad_mass_le

To compare the sample laws, Poissonize the \(S\)-sample at mean \(2n\). Conditional on \(z_i\), the two counts in pair \(i\) are independent Poisson variables with means \((n/h)(1+\xi z_i)\) and \((n/h)(1-\xi z_i)\). Relative to \(z_i=0\), their likelihood ratios obey
\[
E_0[L_zL_{z'}]
=\exp\!\left(\frac{2n\xi^2}{h}zz'\right).
\] % lean: Causalean.Stat.Minimax.Multinomial.TwoSampleL1.scalarPoissonPairLikelihood_inner
For \(s\in\{0,1\}\), define the one-pair mixture, its null reference law, its density, and the exponential scale by
\[
\mathbb Q_s=\int\!\left[
 \operatorname{Pois}\!\left(\frac nh(1+\xi z)\right)
 \otimes\operatorname{Pois}\!\left(\frac nh(1-\xi z)\right)
 \right]d\nu_s(z),\quad
\mathbb Q_\circ=\operatorname{Pois}(n/h)^{\otimes2},\quad
\rho_s=\frac{d\mathbb Q_s}{d\mathbb Q_\circ},\quad
\Delta\nu=\nu_1-\nu_0,\quad
\chi=\frac{2n\xi^2}{h}\le\frac\ell{50}.
\]
The likelihood identity above and the matching moments give
\[
\int(\rho_1-\rho_0)^2\,d\mathbb Q_\circ
=\sum_{\ell'>\ell}\frac{\chi^{\ell'}}{\ell'!}
 \left(\int z^{\ell'}\,d\Delta\nu(z)\right)^2
\le4\sum_{\ell'>\ell}\frac{\chi^{\ell'}}{\ell'!}.
\]
Here \(\bigl|\int z^{\ell'}d\Delta\nu(z)\bigr|\le2\) because both priors are supported on \([-1,1]\). For \(\ell'>\ell\), the factorial bound \(\ell'!\ge(\ell'/e)^{\ell'}\), together with \(e<3\), yields \(\chi^{\ell'}/\ell'!\le(1/4)^{\ell'}\). Cauchy–Schwarz therefore gives
\[
\operatorname{TV}(\mathbb Q_0,\mathbb Q_1)
\le\frac12\left(\int(\rho_1-\rho_0)^2\,d\mathbb Q_\circ\right)^{1/2}
\le\left(\sum_{\ell'>\ell}4^{-\ell'}\right)^{1/2}
\le2^{-\ell}\le2^{-\ell/4}.
\] % lean: Causalean.Stat.Minimax.Multinomial.TwoSampleL1.scalarPoissonPredictive_tv_le
The pair mixtures are independent under either prior. Telescoping their product densities and using the one-pair bound gives
\[
\operatorname{TV}(\mathbb Q_0^{\otimes h},\mathbb Q_1^{\otimes h})
\le h\operatorname{TV}(\mathbb Q_0,\mathbb Q_1)
\le h2^{-\ell/4}.
\]
When a \(\operatorname{Pois}(2n)\) count reaches \(n\), its first \(n\) marks have the fixed-sample law. For each prior, replacing the sample on the complementary event changes its predictive law in total variation by at most \(\Pr\{\operatorname{Pois}(2n)<n\}\). The triangle inequality therefore adds twice this probability. Exponential Markov inequality at \(\log 2\) gives
\[
\Pr\{\operatorname{Pois}(2n)<n\}
\le \exp\{-(1-\log 2)n\}.
\]
Since \(h\le k/2\), \(2^{-\ell/4}\le k^{-4}\), and \(n>k^2\), these bounds give, for the stated sufficiently large alphabet threshold, the fixed-sample predictive bound
\[
\operatorname{TV}(\text{prior-0 samples},\text{prior-1 samples})
\le h\,2^{-\ell/4}
 +2\Pr\{\operatorname{Pois}(2n)<n\}
\le\frac1{16}.
\] % lean: Causalean.Stat.Minimax.Multinomial.TwoSampleL1.pairedFixedPredictive_tv_le
The first sample has the same law \(R\) under both priors, so it leaves this comparison unchanged.

Choose the nearer of \(\gamma_0,\gamma_1\) from an arbitrary estimator of \(\|R-S\|_1\). The sum of its testing errors is at least \(15/16\) by the predictive bound. The two prior tail probabilities contribute at most \(1/4\); on each remaining error event, squared estimation error is at least \(\omega^2/16\). Consequently
\[
\mathcal R_{n,2h}\ge\frac{11}{512}\omega^2,
\qquad
\omega^2=\frac{h}{500000\,n\ell}
\ge c\,\frac{k}{n\log(ek)}
\]
for a universal \(c>0\). Padding the unused cell when \(k\) is odd preserves the target and sample law. % lean: Causalean.Stat.Minimax.Multinomial.TwoSampleL1.largeSample_fuzzyCertificate

The finitely many alphabet sizes below that threshold are covered by a binary comparison. Take \(R=(1/2,1/2)\), compare \(S=R\) with
\(S=((1+\xi)/2,(1-\xi)/2)\), and set \(\xi=1/(64\sqrt n)\). The targets differ by \(\xi\), while independence and the one-record chi-squared identity give
\[
\chi^2(S^{\otimes n},R^{\otimes n})
=(1+\xi^2)^n-1\le\frac1{64}.
\] % lean: Causalean.Stat.Minimax.Multinomial.TwoSampleL1.binary_fuzzyCertificate
Thus their total variation is at most \(1/16\), and the same testing argument yields a constant multiple of \(1/n\). Zero-mass padding and the finite bound on \(k/\log(ek)\) give the asserted rate for these alphabet sizes. Combining the ranges, and using \(k/(n\log(ek))<1\) when \(n>k^2\), proves the displayed two-sample bound. % lean: CausalSmith.Stat.DiscreteBudgetvalueCurve.two_sample_l1_lower_all_samples

Now fix \(T,n,d\) as in the statement. Suppose first that \(d\ge4\), and put \(k=\lfloor d/2\rfloor\). Then \(k\ge2\) and \(2k\le d\). For each \(R,S\in\Delta_k\), choose the paired law of \cref{obj:def:paired-family} with
\[
r_i=\frac{R_i+S_i}{2},\qquad
\theta_i=
\begin{cases}
0,&R_i+S_i=0,\\[2pt]
\dfrac{R_i-S_i}{8(R_i+S_i)},&R_i+S_i>0.
\end{cases}
\] % lean: CausalSmith.Stat.DiscreteBudgetvalueCurve.pairedL1_lower_transport_capacity
The capacity and sample identities in \cref{obj:prop:paired-capacity-identity} give
\[
V_{1/2}(P_{R,S})=\frac38+\frac1{32}\|R-S\|_1,
\qquad V_1(P_{R,S})=\frac12,
\] % lean: CausalSmith.Stat.DiscreteBudgetvalueCurve.pairedL1_lower_transport_capacity
and an exact Markov kernel taking the two independent \(n\)-samples to \(n\) observed records from \(P_{R,S}\). Apply \(T\) after this kernel and condition on the two input samples. The resulting two-sample estimator of \(\|R-S\|_1\) has squared risk at most \(32^2\) times the squared risk of \(T\), by conditional Jensen’s inequality. Since the two-sample lower bound is positive, a parameter pair attains at least half its lower-bound value. Hence some paired law satisfies
\[
E_{P_{R,S}}\!\left[(T-V_{1/2}(P_{R,S}))^2\right]
\ge \frac{c_{\mathrm L}}{2048}
 \min\left\{1,\frac{k}{n\log(ek)}\right\}.
\] % lean: CausalSmith.Stat.DiscreteBudgetvalueCurve.pairedL1_lower_transport_capacity
Pad the law with zero-mass cells to reach \(d\) labels. The observed sample law, both budget values, and the occupied-cell properties are preserved; \cref{obj:ass:iid-sampling} identifies the stipulated sampling law with that product law. % lean: CausalSmith.Stat.DiscreteBudgetvalueCurve.pairedL1_lower_transport_capacity_padded

Because \(k\le d\le3k\) and \(\log(ek)\le\log(ed)\),
\[
\min\left\{1,\frac{d}{n\log(ed)}\right\}
\le3\min\left\{1,\frac{k}{n\log(ek)}\right\}.
\] % lean: CausalSmith.Stat.DiscreteBudgetvalueCurve.curveRate_half_floor
Thus the required bound holds in this case with coefficient \(c_{\mathrm L}/6144\). % lean: CausalSmith.Stat.DiscreteBudgetvalueCurve.capacity_active_lower

For \(d=2,3\), use one occupied pair and pad any remaining cell by zero mass. Compare paired contrasts \(0\) and
\[
t=\frac1{8\sqrt n}.
\] % lean: CausalSmith.Stat.DiscreteBudgetvalueCurve.small_capacity_lower_padded
Both laws satisfy the class, value, occupied-effect, and binding-capacity properties by \cref{obj:prop:paired-capacity-identity}. Their half-budget values are \(3/8\) and \(3/8+t/2\). The one-record chi-squared divergence is \(2t^2\); tensorization and \(2nt^2=1/32\le\log2\) give total variation at most \(1/2\) for the \(n\)-record laws. Testing at the midpoint of the two values then puts error probability at least \(1/4\) under one law, on an event where absolute estimation error is at least \(t/4\). Therefore one of the laws satisfies
\[
E_P\!\left[(T-V_{1/2}(P))^2\right]
\ge\frac{(t/4)^2}{4}
=\frac1{4096n}.
\] % lean: CausalSmith.Stat.DiscreteBudgetvalueCurve.small_capacity_lower_padded
For \(2\le d\le3\), \(\min\{1,d/[n\log(ed)]\}\le3/n\). Taking
\[
c_\epsilon=\min\left\{\frac{c_{\mathrm L}}{6144},
                         \frac1{12288}\right\}>0
\] % lean: CausalSmith.Stat.DiscreteBudgetvalueCurve.capacity_active_lower
completes both cases. % lean: CausalSmith.Stat.DiscreteBudgetvalueCurve.capacity_active_lower
\end{proof}

\begin{proof}[Proof of \cref{obj:prop:fixed-alphabet-reduction}]
Fix \(d\ge2\) and write
\[
\log(ed)=1+\log d,\qquad
\alpha_d:=\frac{d}{\log(ed)}.
\]
Since \(\log d\le d-1\) and \(d>1\), we have \(0<\log(ed)\le d\), so \(\alpha_d\ge1\). Thus, for every \(n\ge1\),
\[
\frac1n
\le \min\left\{1,\frac{d}{n\log(ed)}\right\}
=\min\left\{1,\frac{\alpha_d}{n}\right\}
\le\frac{\alpha_d}{n}.
\]
% lean: CausalSmith.Stat.DiscreteBudgetvalueCurve.fixed_alphabet_reduction

The lower-bound input to \cref{obj:thm:matched-curve-frontier} is the known-distribution \(L_1\) minimax result of \citep[Theorem 3 and the unknown-both-distributions reduction following Theorem 6]{JiaoHanWeissman2018}. Applying the stated frontier bound and then the preceding comparison gives, uniformly over \(b_0\in[0,1/2]\),
\[
\frac{c_\epsilon}{n}
\le c_\epsilon\min\left\{1,\frac{d}{n\log(ed)}\right\}
\le \mathfrak R_{n,d,\epsilon,b_0}^{\mathrm{curve}}
\le C_\epsilon\min\left\{1,\frac{d}{n\log(ed)}\right\}
\le\frac{C_\epsilon\alpha_d}{n}.
\]
% lean: CausalSmith.Stat.DiscreteBudgetvalueCurve.fixed_alphabet_reduction
The risk-bound constants can be chosen as \(c_\epsilon\) and \(C_\epsilon\alpha_d\), respectively; they are positive, ordered, and fixed across \(n\) and \(b_0\).

For the identity, denote the common effect in occupied cells by \(\gamma\), so that
\[
\tau_j=\gamma>0\qquad\text{whenever }p_j>0.
\]
If \(b_0\in[0,1/2]\) and \(b\in[b_0,1-b_0]\), then \(0\le b\le1\). Cell masses satisfy \(p_j\ge0\) and \(\sum_jp_j=1\). In a zero-mass cell both sides of the following identity vanish. Hence, for every \(\pi\in\Pi_b(P)\),
\[
\sum_{j=1}^d p_j\pi_j\tau_j
=\gamma\sum_{j=1}^d p_j\pi_j
\le \gamma b,
\]
where the last inequality is the budget constraint in \cref{obj:def:budget-policy-class}.
% lean: CausalSmith.Stat.DiscreteBudgetvalueCurve.constant_effect_budget_value

Define the constant policy by \(\pi^{(b)}_j=b\) for every cell \(j\). Because \(0\le b\le1\), its coordinates are valid, and
\[
\sum_{j=1}^d p_j\pi^{(b)}_j
=b\sum_{j=1}^d p_j=b,
\qquad
\sum_{j=1}^d p_j\pi^{(b)}_j\tau_j
=\gamma b.
\]
% lean: CausalSmith.Stat.DiscreteBudgetvalueCurve.constant_effect_budget_value
Thus \(\pi^{(b)}\in\Pi_b(P)\) attains the upper bound. The maximization in \cref{obj:def:budget-value} equals \(\gamma b\), and therefore \(V_b(P)=B(P)+\gamma b\), as claimed.
\end{proof}

\begin{proof}[Proof of \cref{obj:thm:curve-upper}]
\textit{1. Constants and unit bound.} Write
\[
L:=\log(ed)=1+\log d,\qquad
\rho_{n,d}:=\frac{d}{nL},\qquad
\gamma_\epsilon:=3(1+\epsilon^{-1})+1.
\]
Let \(C_\epsilon^{\mathrm{proc}}>0\) be the constant supplied by \cref{obj:thm:integrated-process-certificate}, and set
\[
C_\epsilon:=\max\{1,C_\epsilon^{\mathrm{proc}},4096\gamma_\epsilon^2\}>0.
\]
Since \(d\ge2\), the inequality \(\log d\le d-1\) gives \(0<L\le d\), so \(\rho_{n,d}\ge0\).

The interval \(I=[b_0,1-b_0]\) is nonempty and contained in \([0,1]\). Write
\[
p_j=P(X=j),\qquad
\mu_{aj}=\begin{cases}E_P[Y(a)\mid X=j],&p_j>0,\\0,&p_j=0,\end{cases}
\qquad
\tau_j=\mu_{1j}-\mu_{0j},\qquad
B(P)=\sum_{j=1}^d p_j\mu_{0j}.
\]
The binary potential-outcome condition in \cref{obj:def:model-class} gives \(0\le\mu_{aj}\le1\), including the stated convention at empty cells. For any \(b\in I\), the zero policy is feasible under \cref{obj:def:budget-policy-class}, and the value of every feasible policy is
\[
B(P)+\sum_{j=1}^d p_j\pi_j\tau_j
=\sum_{j=1}^d p_j\bigl\{(1-\pi_j)\mu_{0j}+\pi_j\mu_{1j}\bigr\}\in[0,1].
\]
Hence \(V_b(P)\in[0,1]\) by \cref{obj:def:budget-value}. The estimator in \cref{obj:def:jf-estimator} is either \(1/2\) or clipped to \([0,1]\). Under \cref{obj:ass:iid-sampling}, its squared supremum loss therefore satisfies
\[
E_P\!\left[\sup_{b\in I}
 \bigl|\widehat V_b^{\mathrm{JF}}-V_b(P)\bigr|^2\right]\le1.
\]
% lean: CausalSmith.Stat.DiscreteBudgetvalueCurve.jfEstimator_curveRisk_le_one

\textit{2. Alphabets with \(2\le d<16\).} In this branch, \(\widehat F\) is the empirical threshold process of \cref{obj:def:jf-estimator}. It is bounded on \([0,1]\) by \cref{obj:lem:empirical-threshold-process-bounded}. The dual representation and its uniform perturbation bound in \cref{obj:prop:dual-representation}, together with contraction by clipping toward \(V_b(P)\in[0,1]\), give, sample by sample,
\[
\sup_{b\in I}
 \bigl|\widehat V_b^{\mathrm{JF}}-V_b(P)\bigr|^2
\le
\left(\sup_{\lambda\in[0,1]}
 \bigl|\widehat F(\lambda)-F_P(\lambda)\bigr|\right)^2.
\]
% lean: CausalSmith.Stat.DiscreteBudgetvalueCurve.jfEstimator_curveLoss_le_processError

Here is the pointwise estimate used for the empirical process. Let \(\mathcal Z=\{0,1\}^2\), index \(\zeta=(a,y)\), and write
\[
q_{\zeta j}=P(X=j,A=a,Y=y),\qquad
s_a(u)=u_{(a,0)}+u_{(a,1)},\qquad
s(u)=s_0(u)+s_1(u)
\]
for every nonnegative four-vector \(u=(u_\zeta)_{\zeta\in\mathcal Z}\). For nonnegative four-vectors \(u,v\), define
\[
\Delta(u,v):=\sum_{\zeta\in\mathcal Z}|u_\zeta-v_\zeta|.
\]
For the formula in \cref{obj:def:dual-process}, explicitly define its value at the zero vector by \(g_a(0)=0\). This is its continuous extension, since \(0\le g_a(u)\le s(u)\) for every nonnegative \(u\) with \(s(u)>0\). Thus \(f_\lambda(0)=0\), and the stabilized arm contributions, including those of empty covariate cells, obey
\[
|g_a(u)-g_a(v)|\le(1+\epsilon^{-1})\Delta(u,v),
\qquad
|s(u)-s(v)|\le\Delta(u,v).
\]
To see the first inequality, an arm with \(s_a(u)\le\epsilon s(u)\) has \(g_a(u)=u_{a1}/\epsilon\); an arm with \(s_a(u)\ge\epsilon s(u)>0\) has \(g_a(u)=s(u)u_{a1}/s_a(u)\). When both vectors are in the latter regime, use
\[
\left|\frac{u_{a1}}{s_a(u)}-\frac{v_{a1}}{s_a(v)}\right|
\le
\frac{|u_{a1}-v_{a1}|+|u_{a0}-v_{a0}|}{s_a(u)},
\qquad
\frac{s(u)}{s_a(u)}\le\epsilon^{-1},
\]
and \(u_{a1}/s_a(u),v_{a1}/s_a(v)\in[0,1]\). When both vectors are in the first regime, \(|g_a(u)-g_a(v)|=|u_{a1}-v_{a1}|/\epsilon\le\epsilon^{-1}\Delta(u,v)\). When \(u\) is in the first regime and \(v\) in the second, use
\[
g_a(u)-g_a(v)
=
\frac{u_{a1}-v_{a1}}{\epsilon}
+
\frac{v_{a1}}{s_a(v)}
 \frac{s_a(v)-\epsilon s(v)}{\epsilon}.
\]
The second term is nonnegative and \(0\le v_{a1}/s_a(v)\le1\). Since \(s_a(u)\le\epsilon s(u)\), its numerator satisfies
\[
s_a(v)-\epsilon s(v)
\le s_a(v)-s_a(u)+\epsilon\{s(u)-s(v)\}.
\]
Combining these inequalities cancels the \(a1\) coordinate in the upper bound and gives
\[
-\epsilon^{-1}\Delta(u,v)
\le g_a(u)-g_a(v)
\le \epsilon^{-1}(v_{a0}-u_{a0})+s(u)-s(v)
\le(1+\epsilon^{-1})\Delta(u,v).
\]
Swapping \(u\) and \(v\) handles the reverse case. If either total mass is zero, that vector is zero and \(0\le g_a(w)\le s(w)\) gives the bound directly. Since \(x\mapsto[x]_+\) is one-Lipschitz and \(0\le\lambda\le1\), it follows that
\[
|f_\lambda(u)-f_\lambda(v)|
\le
2|g_0(u)-g_0(v)|+|g_1(u)-g_1(v)|
 +\lambda|s(u)-s(v)|
\le\gamma_\epsilon\Delta(u,v).
\]
% lean: CausalSmith.Stat.DiscreteBudgetvalueCurve.budget_thresholdFunReal_pointwise_lipschitz

Applying this cellwise and using Cauchy--Schwarz across the \(4d\) category-cell coordinates yields
\[
\left(\sup_{\lambda\in[0,1]}
 |\widehat F(\lambda)-F_P(\lambda)|\right)^2
\le
4d\,\gamma_\epsilon^2
\sum_{j=1}^d\sum_{\zeta\in\mathcal Z}
 (\widehat q_{\zeta j}-q_{\zeta j})^2.
\]
Each \(\widehat q_{\zeta j}\) is the empirical mass of one observed atom. Its Bernoulli second moment under \cref{obj:ass:iid-sampling} gives
\[
E_P(\widehat q_{\zeta j}-q_{\zeta j})^2
=\frac{q_{\zeta j}(1-q_{\zeta j})}{n}\le\frac1n.
\]
Consequently,
\[
E_P\!\left[\sup_{b\in I}
 |\widehat V_b^{\mathrm{JF}}-V_b(P)|^2\right]
\le
\frac{16d^2\gamma_\epsilon^2}{n}.
\]
% lean: CausalSmith.Stat.DiscreteBudgetvalueCurve.jfEstimator_empiricalCurveRisk_le

Because \(L\le d<16\), we have \(dL\le256\), and hence
\[
\frac{16d^2\gamma_\epsilon^2}{n}
\le4096\gamma_\epsilon^2\rho_{n,d}.
\]
If \(\rho_{n,d}\ge1\), the unit bound above is at most \(C_\epsilon\min\{1,\rho_{n,d}\}\). If \(\rho_{n,d}<1\), the displayed empirical bound is at most the same quantity.
% lean: CausalSmith.Stat.DiscreteBudgetvalueCurve.empiricalCurveRate_algebra

\textit{3. Alphabets with \(d\ge16\) and \(d/L\le n\).} Here \(\rho_{n,d}\le1\). The estimator uses the conditionally averaged threshold process \(\widehat F\) specified in \cref{obj:def:jf-estimator}. It is bounded on \([0,1]\) by \cref{obj:lem:averaged-threshold-process-bounded}. The same dual perturbation and clipping argument bounds its curve risk by the averaged process risk:
\[
E_P\!\left[\sup_{b\in I}
 |\widehat V_b^{\mathrm{JF}}-V_b(P)|^2\right]
\le
E_P\!\left[
 \left(\sup_{\lambda\in[0,1]}
 |\widehat F(\lambda)-F_P(\lambda)|\right)^2\right].
\]
% lean: CausalSmith.Stat.DiscreteBudgetvalueCurve.jfEstimator_averagedCurveRisk_le

Under the stated sampling and model conditions, \cref{obj:thm:integrated-process-certificate} bounds the right-hand side by \(C_\epsilon^{\mathrm{proc}}\rho_{n,d}\). Thus the curve risk is at most
\[
C_\epsilon^{\mathrm{proc}}\rho_{n,d}
\le C_\epsilon\min\{1,\rho_{n,d}\}.
\]
% lean: CausalSmith.Stat.DiscreteBudgetvalueCurve.integrated_process_certificate

\textit{4. Alphabets with \(d\ge16\) and \(n<d/L\).} Now \(\rho_{n,d}>1\), and \cref{obj:def:jf-estimator} sets the entire estimated curve equal to \(1/2\). The unit bound above gives
\[
E_P\!\left[\sup_{b\in I}
 |\widehat V_b^{\mathrm{JF}}-V_b(P)|^2\right]
\le1\le C_\epsilon\min\{1,\rho_{n,d}\}.
\]
% lean: CausalSmith.Stat.DiscreteBudgetvalueCurve.jfEstimator_curveRisk_le_one

These cases give the claimed inequality for every \(P\in\mathcal M_{d,\epsilon}\), with one constant depending only on \(\epsilon\). Consider the set of expected squared losses indexed by these laws. If this set is nonempty, its supremum obeys the same bound because every member does. If it is empty, define its supremum as zero; the bound then follows from \(C_\epsilon>0\) and \(\rho_{n,d}\ge0\). In the present model the set is nonempty: take \(X\) uniform on \([d]\), take \(A\) independent of \(X\) with treatment probability \(1/2\), and set \(Y(0)=Y(1)=Y=0\) almost surely. This law satisfies consistency and conditional exchangeability, and its treatment propensity in every cell is \(1/2\), as required by \(0<\epsilon<1/2\). Thus the nonempty case gives the stated supremum bound.
% lean: CausalSmith.Stat.DiscreteBudgetvalueCurve.curve_upper
\end{proof}

\begin{proof}[Proof of \cref{obj:thm:integrated-process-certificate}]
Write \(L=\log(ed)\), \(m=n/8\), and \(K=\max\{2,\lfloor L/512\rfloor\}\). We first establish estimates for an arbitrary observed categorical law in the independent pilot and evaluation count experiment of \cref{obj:def:process-handle}. All constants introduced below are positive and depend at most on \(\epsilon\), unless called numerical.

For each pilot realization, every point \(u\) in the rectangle \(\prod_{\zeta\in\mathcal Z}[\ell_{\zeta j},u_{\zeta j}]\) is nonnegative and satisfies \(\|u-c_j^0\|_1\leq S_j\). Positivity and normalization of the Jackson convolution, together with \cref{obj:lem:bv-lipschitz}, therefore give the uniform bound
\[
\sup_{u\in\prod_{\zeta\in\mathcal Z}[\ell_{\zeta j},u_{\zeta j}]}
\|P_{j,\cdot}(u)-f_\cdot(c_j^0)\|_{BV([0,1])}
\leq c_{\mathrm{BV}}S_j.
\]
Write the polynomial on this rectangle in tensor Chebyshev form, with coefficients \(b_{j,\nu}\):
\[
P_{j,\lambda}(u)-f_\lambda(c_j^0)
 =\sum_{\nu\in\{0,\ldots,D\}^4}b_{j,\nu}(\lambda)
   \prod_{\zeta\in\mathcal Z}
   T_{\nu_\zeta}\!\left(\frac{u_\zeta-c^0_{\zeta j}}{R_{\zeta j}}\right),
\qquad D=2(K-1).
\]
Each coefficient is a normalized tensor Chebyshev integral of the uniformly bounded path, so \(\|b_{j,\nu}\|_{BV([0,1])}\leq16c_{\mathrm{BV}}S_j\). If \(A_h\) is the sum of the absolute monomial coefficients of \(T_h\), its recurrence gives \(A_0=A_1=1\), \(A_{h+1}\leq2A_h+A_{h-1}\), and hence \(A_h\leq(1+\sqrt2)^h\). Conversion to the centered monomial coefficients of \cref{obj:def:jf-estimator} now yields
\[
\sum_\alpha\|\beta_{j,\alpha}\|_{BV([0,1])}
 \leq16c_{\mathrm{BV}}S_j(D+1)^4(1+\sqrt2)^{4D}
 \leq c_{\mathrm{coeff}}S_j e^{12K}.
\]
The estimate holds for every pilot realization, including those outside \(G_j\): the pilot radii are strictly positive, and the convolution rectangle lies in \(\mathbb R_+^4\). Also \(L=1+\log d\geq1\) and \(K\leq2+L/512\). Direct expansion of this last bound yields
\[
24K+\frac{16K^2}{L}\leq114+\frac L{16}
 \leq115+\frac{\log d}{16},
\qquad
S_j^2e^{24K+16K^2/L}\leq e^{115}d^{1/16}S_j^2.
\]
These are the two coefficient bounds required by \cref{obj:def:process-handle}. % lean: CausalSmith.Stat.DiscreteBudgetvalueCurve.jacksonCoefficientBV_pilot_le
% lean: CausalSmith.Stat.DiscreteBudgetvalueCurve.jacksonDegree_exponential_budget

Here are the cell estimates used below. Put
\[
a_j=\sqrt{\frac{p_j}{mL}}+\frac1{mL},
\qquad
v_j=\sqrt{\frac{p_jL}{m}}+\frac Lm.
\]
Conditioning on the pilot counts, factorial moments of the independent evaluation counts give
\[
\mathbb E\!\left[U_h(N_{\zeta j};c^0_{\zeta j})\mid N'\right]
 =(q_{\zeta j}-c^0_{\zeta j})^h.
\]
On \(G_j\), the localization interval contains \(q_{\zeta j}\), and its radius satisfies
\(q_{\zeta j}/(mR_{\zeta j}^{\,2})\leq1/L\). The second factorial-moment identity is
\[
\mathbb E\!\left[U_h(N_{\zeta j};c^0_{\zeta j})^2\mid N'\right]
 =\sum_{r=0}^{h}{h\choose r}^{\!2}r!
   \left(\frac{q_{\zeta j}}m\right)^r
   (q_{\zeta j}-c^0_{\zeta j})^{2h-2r}.
\]
Using \(|q_{\zeta j}-c^0_{\zeta j}|\leq R_{\zeta j}\), \(q_{\zeta j}/(mR_{\zeta j}^2)\leq1/L\), and \({h\choose r}^2r!\leq h^{2r}/r!\), it gives
\[
\mathbb E\!\left[\left.
 \frac{U_h(N_{\zeta j};c^0_{\zeta j})^2}{R_{\zeta j}^{2h}}
 \right|N'\right]
 \leq\sum_{r=0}^{h}\frac{(h^2/L)^r}{r!}
 \leq e^{h^2/L}.
\]
For each factorial index \(\alpha\in\{0,\ldots,D\}^{\mathcal Z}\), independence of the four evaluation coordinates and the preceding moment bound give, on \(G_j\),
\[
\mathbb E\!\left[\left.
\prod_{\zeta\in\mathcal Z}
\frac{U_{\alpha_\zeta}(N_{\zeta j};c^0_{\zeta j})^2}
     {R_{\zeta j}^{2\alpha_\zeta}}
\right|N'\right]
\leq \exp\!\left(\frac{\sum_\zeta\alpha_\zeta^2}{L}\right)
\leq e^{4D^2/L}.
\]
Cauchy–Schwarz over the \((D+1)^4\) factorial indices and the coefficient envelope therefore give
\[
\mathbb E\!\left[\|Z_j-f_\cdot(c_j^0)\|_{BV([0,1])}^2\mid N'\right]
\leq 4\left(\sum_\alpha\|\beta_{j,\alpha}\|_{BV([0,1])}\right)^2
       \sum_{\alpha\in\{0,\ldots,D\}^{\mathcal Z}}e^{4D^2/L}
\leq 4c_{\mathrm{coeff}}^2S_j^2(D+1)^4
       e^{24K+16K^2/L}.
\]
Here \(D+1\leq2K\leq e^K\). Using \(K\leq2+L/512\) gives
\[
28K+\frac{16K^2}{L}\leq122+\frac L{16}
=123+\frac{\log d}{16},
\qquad
(D+1)^4e^{24K+16K^2/L}\leq e^{123}d^{1/16}.
\]
Consequently the conditional squared bounded-variation norm is at most \(4c_{\mathrm{coeff}}^2e^{123}d^{1/16}S_j^2\). Every radius is positive, so \(S_j>0\). Apply \(|\Pi_{[-t,t]}(x)-x|\leq x^2/t\) with \(t=d^{1/4}S_j\). Uniformly in \(\lambda\),
\[
\mathbb E\!\left[|T_j(\lambda)-Z_j(\lambda)|\mid N'\right]
 \leq 4c_{\mathrm{coeff}}^2e^{123}S_jd^{-3/16}
 \qquad\text{on }G_j.
\]
% lean: CausalSmith.Stat.DiscreteBudgetvalueCurve.goodPilot_localJacksonFourPoissonFactorialPath_sq_integral_le The factorial-moment identity also gives
\[
\mathbb E[Z_j(\lambda)\mid N']=P_{j,\lambda}(q_j).
\]
To bound its approximation error, put \(\delta=L/m\). On \(G_j\), the vector \(q_j\) lies in the pilot rectangle, so \cref{obj:lem:jackson-boundary-adaptive} gives
\[
|P_{j,\lambda}(q_j)-f_\lambda(q_j)|
\leq32\{3(1+\epsilon^{-1})+1\}
\left\{
\frac1K\sum_{\zeta\in\mathcal Z}
\sqrt{(q_{\zeta j}-\ell_{\zeta j})(u_{\zeta j}-q_{\zeta j})}
+\frac{S_j}{K^2}\right\}.
\]
The good-pilot inequality and the truncated endpoints imply, coordinatewise,
\[
(q_{\zeta j}-\ell_{\zeta j})(u_{\zeta j}-q_{\zeta j})
\leq8H_0^2q_{\zeta j}\delta.
\]
Indeed, when \(\ell_{\zeta j}>0\), the product is at most \(h_{\zeta j}^2\); the inequalities \(c_{\zeta j}>h_{\zeta j}\) and \(|c_{\zeta j}-q_{\zeta j}|\leq h_{\zeta j}/4\), inserted into the formula for \(h_{\zeta j}\), give the displayed bound. When \(\ell_{\zeta j}=0\), that formula and \(c_{\zeta j}\leq h_{\zeta j}\) give \(u_{\zeta j}\leq8H_0^2\delta\), so the product is at most \(q_{\zeta j}u_{\zeta j}\). Since each \(q_{\zeta j}\leq p_j\) and there are four coordinates,
\[
\sum_{\zeta\in\mathcal Z}
\sqrt{(q_{\zeta j}-\ell_{\zeta j})(u_{\zeta j}-q_{\zeta j})}
\leq4\sqrt{8H_0^2p_jL/m}.
\]
Thus the approximation contribution has the asserted \(\sqrt{p_jL/m}/K+S_j/K^2\) scale, including when \(p_j=0\). Combining it with the conditional clipping contribution gives, on \(G_j\),
% lean: CausalSmith.Stat.DiscreteBudgetvalueCurve.goodPilot_localJacksonPolynomial_boundary_sharp
\[
\sup_{\lambda\in[0,1]}
\left|\mathbb E\{T_j(\lambda)\mid N'\}-f_\lambda(q_j)\right|
 \leq c_\epsilon\left\{
 S_jd^{-3/16}
 +\frac{\sqrt{p_jL/m}}K+\frac{S_j}{K^2}\right\}.
\]
The radius definition and truncation at zero give the pointwise estimates
\[
R_{\zeta j}^{2}
\leq2H_0^2\left\{\frac{N'_{\zeta j}L}{m^2}
+\left(\frac Lm\right)^2\right\},
\qquad
S_j^2\leq4\sum_{\zeta\in\mathcal Z}R_{\zeta j}^2.
\]
Taking expectations and using \(\mathbb E N'_{\zeta j}=m q_{\zeta j}\) and \(\sum_\zeta q_{\zeta j}=p_j\) yields
\[
\mathbb E S_j^2
\leq8H_0^2\left\{\frac{p_jL}{m}
+4\left(\frac Lm\right)^2\right\}
\leq32H_0^2v_j^2.
\]
This is the second-moment bound in \cref{obj:lem:pilot-radius-second-moment}; Cauchy–Schwarz also gives \(\mathbb E S_j\leq\sqrt{32}\,H_0v_j\).
% lean: CausalSmith.Stat.DiscreteBudgetvalueCurve.pilotRadius_sum_sq_integral_le_cellScale Finally,
\(K^{-1}\leq768/L\) and \(L^2d^{-3/16}\leq(35/3)^2\) for \(d\geq16\). After averaging over pilots, these inequalities convert the displayed good-pilot bound to the scale \(a_j\).

For the complementary pilot event, put
\[
B_j=\sum_{\zeta\in\mathcal Z}
 \bigl(|c_{\zeta j}-q_{\zeta j}|+h_{\zeta j}\bigr).
\]
Applying \cref{obj:lem:pilot-bad-score-moment} with \(t=2\) gives the weighted second-moment estimate
\[
\mathbb E[\mathbf1_{G_j^c}B_j^2]
\leq 16\cdot10^{240}e^{-20L}v_j^2
=:c_{\mathrm{mom}}e^{-20L}v_j^2.
\]
% lean: CausalSmith.Stat.DiscreteBudgetvalueCurve.pilotCell_badScore_sq_integral_le
The pilot endpoints give \(|c^0_{\zeta j}-c_{\zeta j}|\leq h_{\zeta j}/2\) and \(R_{\zeta j}\leq h_{\zeta j}\). Clipping bounds \(|T_j(\lambda)-f_\lambda(c_j^0)|\) by \(d^{1/4}S_j\), while \cref{obj:lem:bv-lipschitz} bounds \(\sup_\lambda|f_\lambda(c_j^0)-f_\lambda(q_j)|\) by \(c_{\mathrm{BV}}\|c_j^0-q_j\|_1\). Hence, on \(G_j^c\),
\[
\sup_{\lambda\in[0,1]}|T_j(\lambda)-f_\lambda(q_j)|
 \leq(c_{\mathrm{BV}}+d^{1/4})
 \sum_{\zeta\in\mathcal Z}
 \bigl(|c_{\zeta j}-q_{\zeta j}|+h_{\zeta j}\bigr).
\]
Since \(d^{1/4}\geq1\), taking square roots yields, for a positive \(c_{\mathrm{env}}\),
\[
\left\|\mathbf1_{G_j^c}
 \sup_{\lambda\in[0,1]}|T_j(\lambda)-f_\lambda(q_j)|
 \right\|_{L_2}
 \leq c_{\mathrm{env}}d^{1/4}e^{-10L}v_j.
\]
Splitting the expectation over \(G_j\) and \(G_j^c\), and using the preceding conditional and bad-event bounds, gives a positive \(c_{\mathrm{bias}}\) such that
\[
\sup_{\lambda\in[0,1]}
 \left|\mathbb ET_j(\lambda)-f_\lambda(q_j)\right|
 \leq c_{\mathrm{bias}}a_j.
\]
In particular, the exponentially weighted bad-event contribution is absorbed at this scale because \(v_j\leq L^2a_j\) and
\(d^{1/4}e^{-10L}L^2\) is bounded for \(d\geq16\). % lean: CausalSmith.Stat.DiscreteBudgetvalueCurve.uncappedJacksonCellBias_le_rate
% lean: CausalSmith.Stat.DiscreteBudgetvalueCurve.badPilotCellEnvelope_le

Define the good- and bad-pilot error paths by
\[
W_j^{\mathrm{good}}(\lambda)
 =\mathbf1_{G_j}\{T_j(\lambda)-f_\lambda(q_j)\},
\qquad
W_j^{\mathrm{bad}}(\lambda)
 =\mathbf1_{G_j^c}\{T_j(\lambda)-f_\lambda(q_j)\}.
\]
The pilot and evaluation counts are independent across cells. On \(G_j\), clipping is one-Lipschitz and cannot increase total variation. The preceding conditional factorial-moment bound, the coefficient envelope, and \cref{obj:lem:bv-lipschitz} therefore give
\[
\mathbb E\!\left[\|W_j^{\mathrm{good}}\|_{BV([0,1])}^2\mid N'\right]
 \leq c\,\mathbf1_{G_j}d^{1/16}S_j^2.
\]
Indeed, the clipped factorial path has squared bounded-variation norm at most a constant times \(d^{1/16}S_j^2\); its correction from the pilot center to \(q_j\) has bounded-variation norm at most \(c_\epsilon S_j\), since \(q_j\) lies in the pilot rectangle. The second-moment calculation above gives \(\mathbb ES_j^2\leq32H_0^2v_j^2\). Consequently,
\[
\mathbb E\|W_j^{\mathrm{good}}\|_{BV([0,1])}^2
 \leq c\,d^{1/16}v_j^2.
\]
The paths are independent across cells, continuous, and have the finite variation and second moments established above. By \cref{obj:lem:bv-maximal-cell-paths}, and because \(\|W_j^{\mathrm{good}}\|_\infty+\operatorname{TV}(W_j^{\mathrm{good}})\leq2\|W_j^{\mathrm{good}}\|_{BV([0,1])}\),
\[
\begin{aligned}
\mathbb E\left\|\sum_j
 (W_j^{\mathrm{good}}-\mathbb EW_j^{\mathrm{good}})
 \right\|_\infty^2
&\leq 16384\sum_j\mathbb E
 \left[\bigl(\|W_j^{\mathrm{good}}\|_\infty+
 \operatorname{TV}(W_j^{\mathrm{good}})\bigr)^2\right]\\
&\leq65536\sum_j\mathbb E
 \|W_j^{\mathrm{good}}\|_{BV([0,1])}^2\\
&\leq c\,d^{1/16}
 \left(\frac{2L}{m}+2d\frac{L^2}{m^2}\right).
\end{aligned}
\]
The last scale follows from \(\sum_jp_j=1\) and
\(\sum_jv_j^2\leq2L/m+2d(L/m)^2\). On the branch \(d/L\leq n=8m\), the inequalities \(L\leq17d^{1/16}\) and \(d/m\leq8L\) turn the last display into
\[
\mathbb E\left\|\sum_j
 (W_j^{\mathrm{good}}-\mathbb EW_j^{\mathrm{good}})
 \right\|_\infty^2
 \leq c_{\mathrm{good}}\frac d{mL}.
\]
% lean: CausalSmith.Stat.DiscreteBudgetvalueCurve.centeredIdealGoodPilotErrorRisk_le_rate

For the bad-pilot paths, centering, the triangle inequality, and Cauchy–Schwarz give
\[
\mathbb E\left\|\sum_j
 (W_j^{\mathrm{bad}}-\mathbb EW_j^{\mathrm{bad}})
 \right\|_\infty^2
 \leq4d\sum_j\mathbb E\|W_j^{\mathrm{bad}}\|_\infty^2
 \leq c\,d^{3/2}e^{-20L}
 \left(\frac{2L}{m}+2d\frac{L^2}{m^2}\right).
\]
Here \(L\leq d\), \(dL/m\leq8L^2\), and
\(d^{1/2}e^{-20L}\leq d^{-4}\). They absorb the remaining factors and give a positive \(c_{\mathrm{bad}}\) with
\[
\mathbb E\left\|\sum_j
 (W_j^{\mathrm{bad}}-\mathbb EW_j^{\mathrm{bad}})
 \right\|_\infty^2
 \leq c_{\mathrm{bad}}\frac d{mL}.
\]
% lean: CausalSmith.Stat.DiscreteBudgetvalueCurve.centeredIdealBadPilotErrorRisk_le_rate

It remains to combine the centered terms with their mean. Define
\[
\mathfrak b=c_{\mathrm{bias}}\sum_j a_j.
\]
The good and bad error paths partition \(\sum_j(T_j-f_\cdot(q_j))\), so the absolute value of its mean at every \(\lambda\) is at most \(\mathfrak b\). Since \(\sum_jp_j=1\), Cauchy–Schwarz and \(d/(nL)\leq1\) give
\[
\mathfrak b^2
 \leq144c_{\mathrm{bias}}^2\frac d{nL}.
\]
Applying \((x+y+z)^2\leq3(x^2+y^2+z^2)\) to the centered good path, centered bad path, and mean, then using \(m=n/8\), proves
\[
\mathbb E_{\mathrm{id}}
 \left\|\sum_jT_j-F_P\right\|_\infty^2
 \leq\Gamma_{\mathrm{id}}\frac d{nL},
\qquad
\Gamma_{\mathrm{id}}
 =24c_{\mathrm{good}}+24c_{\mathrm{bad}}
   +432c_{\mathrm{bias}}^2.
\]
% lean: CausalSmith.Stat.DiscreteBudgetvalueCurve.idealProcessRisk_le_rate

For the averaged process of \cref{obj:def:jf-estimator}, define the nonnegative count-table function
\[
\mathcal G(N',N)
 =\left\|\sum_jT_j(\cdot;N',N)-F_P\right\|_\infty^2.
\]
Its ideal expectation is finite by the preceding process bound. Conditional Jensen bounds the averaged-process risk by the capped auxiliary risk. Apply \cref{obj:lem:capped-marked-count-identity} to \(\mathcal G\); in the uncapped marked-Poisson experiment, \(M\) is the total number of marked records, and the count-table marginal is the ideal law. Hence
\[
\mathbb E_{O,\mathrm{aux}}[\mathbf1_{\{M\leq n\}}\mathcal G(N',N)]
 =\mathbb E_{\mathrm{uncap}}[\mathbf1_{\{M\leq n\}}\mathcal G(N',N)]
 \leq\mathbb E_{\mathrm{id}}\mathcal G(N',N).
\]
On \(M>n\), the auxiliary process is zero; \(0\leq f_\lambda(q_j)\leq p_j\) implies \(\|F_P\|_\infty\leq1\). The Poisson cap-tail bound therefore gives
\[
\begin{aligned}
\mathbb E\|\widehat F-F_P\|_\infty^2
&\leq\mathbb E_{O,\mathrm{aux}}\|\widetilde F^{\mathrm{cap}}-F_P\|_\infty^2\\
&=\mathbb E_{O,\mathrm{aux}}[\mathbf1_{\{M\leq n\}}\mathcal G(N',N)]
  +\Pr(M>n)\|F_P\|_\infty^2\\
&\leq\mathbb E_{\mathrm{id}}
       \left\|\sum_jT_j-F_P\right\|_\infty^2
       +e^{-n/16}.
\end{aligned}
\]
The elementary bounds \(e^{-n/16}\leq16/n\) and
\(L=1+\log d\leq d\) show that \(e^{-n/16}\leq16d/(nL)\). Hence
\[
\mathbb E\|\widehat F-F_P\|_\infty^2
 \leq(\Gamma_{\mathrm{id}}+16)\frac d{nL}.
\]
Under \cref{obj:ass:iid-sampling}, the product law used in this expectation is the law of the observed sample. % lean: CausalSmith.Stat.DiscreteBudgetvalueCurve.capTransferRisk
% lean: CausalSmith.Stat.DiscreteBudgetvalueCurve.exponential_cap_remainder_le_rate

Choose one constant larger than every coefficient used above; for example, with the cap-transfer coefficient written \(c_{\mathrm{cap}}=1\), take
\[
C_\epsilon
 =1+c_{\mathrm{coeff}}+e^{115}
   +8c_{\mathrm{bias}}+8c_{\mathrm{good}}
   +c_{\mathrm{env}}+c_{\mathrm{bad}}
   +\Gamma_{\mathrm{id}}
   +c_{\mathrm{cap}}(\Gamma_{\mathrm{id}}+16).
\]
All terms are nonnegative, so \(C_\epsilon>0\). The five estimates established for arbitrary observed categorical laws and every pilot realization give \(C_\epsilon\in\mathsf H_{n,d}\) as defined in \cref{obj:def:process-handle}.

Finally, take the observed margin of the full-data law \(P\in\mathcal M_{d,\epsilon}\) from \cref{obj:def:model-class}. Summing its four observed category masses in cell \(j\) gives the same covariate mass \(p_j\) as summing the corresponding full-data masses. Moreover,
\[
\sqrt{\frac{p_j}{mL}}+\frac1{mL}
 =\sqrt{\frac{8p_j}{nL}}+\frac8{nL}
 \leq8\left\{\sqrt{\frac{p_j}{nL}}+\frac1{nL}\right\},
\qquad
\frac d{mL}=8\frac d{nL}.
\]
Thus the common choice \(C_\epsilon\) gives the stated \(nL\)-scale cell-bias and good-pilot bounds, the stated \(mL\)-scale bad-pilot bound, the cellwise bad-pilot envelope, and both process-risk bounds simultaneously. % lean: CausalSmith.Stat.DiscreteBudgetvalueCurve.integrated_process_certificate
\end{proof}

\begin{proof}[Proof of \cref{obj:prop:paired-capacity-identity}]
Fix a paired completion \(P\) with parameters \(r,\theta\). By \cref{obj:def:paired-family}, each cell \((i,\sigma)\), where \(\sigma\in\{+1,-1\}\), has mass \(p_{i,\sigma}=r_i/2\), treatment propensity \(1/2\) when occupied, and weighted outcome means
\[
p_{i,\sigma}\mu_{0,i,\sigma}=\frac{r_i}{2}\frac14,
\qquad
p_{i,\sigma}\mu_{1,i,\sigma}
=\frac{r_i}{2}\left(\frac12+\sigma\theta_i\right).
\] % lean: CausalSmith.Stat.DiscreteBudgetvalueCurve.PairedCompletion
These identities also hold when \(r_i=0\). The paired completion satisfies consistency and conditional exchangeability. Since \(k\geq1\), its alphabet has at least two cells; since \(0<\epsilon<1/2\), propensity \(1/2\) satisfies overlap on every occupied cell. Thus \(P\in\mathcal M_{2k,\epsilon}\) by \cref{obj:def:model-class}. On an occupied cell, division by \(p_{i,\sigma}>0\) gives
\[
\tau_{i,\sigma}
=\mu_{1,i,\sigma}-\mu_{0,i,\sigma}
=\frac14+\sigma\theta_i
\in\left[\frac18,\frac38\right].
\] % lean: CausalSmith.Stat.DiscreteBudgetvalueCurve.pairedCompletion_effect_bounds

Summing the weighted control means and using \(\sum_i r_i=1\) gives
\[
B(P)=\sum_{i,\sigma}p_{i,\sigma}\mu_{0,i,\sigma}=\frac14.
\] % lean: CausalSmith.Stat.DiscreteBudgetvalueCurve.pairedCompletion_controlValue
Choose the policy that treats the \(+\) cell when \(\theta_i\geq0\) and the \(-\) cell when \(\theta_i<0\):
\[
\pi^\star_{i,+}=\mathbf1\{\theta_i\geq0\},
\qquad
\pi^\star_{i,-}=\mathbf1\{\theta_i<0\}.
\]
It selects exactly one cell per pair, including at a tie. Consequently,
\[
\sum_{i,\sigma}p_{i,\sigma}\pi^\star_{i,\sigma}
=\sum_i\frac{r_i}{2}=\frac12,
\qquad
B(P)+\sum_{i,\sigma}p_{i,\sigma}\pi^\star_{i,\sigma}\tau_{i,\sigma}
=\frac38+\frac12\sum_i r_i|\theta_i|.
\] % lean: CausalSmith.Stat.DiscreteBudgetvalueCurve.pairedSignPolicy_budget
% lean: CausalSmith.Stat.DiscreteBudgetvalueCurve.pairedSignPolicy_value

To bound the value above, every \(\pi\in\Pi_b(P)\) and every \(\lambda\geq0\) satisfy
\[
\sum_j p_j\pi_j\tau_j
\leq b\lambda+\sum_j p_j[\tau_j-\lambda]_+.
\]
Indeed, \(p_j\pi_j(\tau_j-\lambda)\leq p_j[\tau_j-\lambda]_+\) cell by cell, while \(\lambda\sum_jp_j\pi_j\leq b\lambda\). This bound also makes the feasible gains bounded above in the supremum defining \(V_b(P)\). % lean: CausalSmith.Stat.DiscreteBudgetvalueCurve.weighted_budget_weak_duality

For the observed margin of \(P\), the arm mass in an occupied cell is \(p_j/2\). Hence the stabilized contributions in \cref{obj:def:dual-process} equal \(p_j\mu_{aj}\). In a null cell all four observed masses vanish, so \(q_j=0\). To specify the quotient at this point, interpret \(g_a\) by its zero-mass extension: \(g_a(0)=0\) for each arm \(a\). The formula in \cref{obj:def:dual-process} then gives \(f_{1/4}(0)=0\). At \(\lambda=1/4\), this yields
\[
F_P\!\left(\frac14\right)
=B(P)+\sum_jp_j\left[\tau_j-\frac14\right]_+
=\frac14+\frac12\sum_i r_i
 \bigl([\theta_i]_++[-\theta_i]_+\bigr)
=\frac14+\frac12\sum_i r_i|\theta_i|.
\] % lean: CausalSmith.Stat.DiscreteBudgetvalueCurve.pairedCompletion_dualQuarter
Applying the preceding bound at \(b=1/2\) and \(\lambda=1/4\), and then using the feasible policy \(\pi^\star\) for the reverse inequality, gives
\[
V_{1/2}(P)
=\frac38+\frac12\sum_i r_i|\theta_i|
=\frac18+F_P\!\left(\frac14\right)
=B(P)+\sum_jp_j\pi^\star_j\tau_j.
\] % lean: CausalSmith.Stat.DiscreteBudgetvalueCurve.pairedCompletion_bindingValue
Thus \(\pi^\star\) attains the half-population value while using treatment mass exactly \(1/2\).

At budget one, treating every cell is feasible. Every occupied-cell effect is positive, so no feasible policy has greater gain. The weighted treated means cancel pairwise:
\[
V_1(P)
=\sum_{i=1}^k\frac{r_i}{2}
 \left[\left(\frac12+\theta_i\right)
       +\left(\frac12-\theta_i\right)\right]
=\frac12.
\] % lean: CausalSmith.Stat.DiscreteBudgetvalueCurve.pairedCompletion_budgetValue_one

For the sample reduction, fix \(R,S\in\Delta_k\) and define \(r,\theta\) as in the statement. Nonnegativity and the simplex equalities for \(R,S\) give \(r\in\Delta_k\). If \(R_i+S_i>0\), then \(|R_i-S_i|\leq R_i+S_i\), so \(|\theta_i|\leq1/8\); when the sum is zero, \(\theta_i=0\). Thus \(\theta\) is an admissible contrast vector. The construction in \cref{obj:def:paired-family} supplies a paired law with these parameters. Moreover, for every \(i\), including a zero-mass pair,
\[
r_i\theta_i=\frac{R_i-S_i}{16}.
\] % lean: CausalSmith.Stat.DiscreteBudgetvalueCurve.normalizedPairedParameters
% lean: CausalSmith.Stat.DiscreteBudgetvalueCurve.normalizedPairedMassContrast
% lean: CausalSmith.Stat.DiscreteBudgetvalueCurve.pairedLaw_completion

Here is a kernel that depends on the input records but not on \(R,S\). Denote the two independent input samples by
\[
(\xi_t)_{t=1}^n\sim R^{\otimes n},\qquad
(\eta_t)_{t=1}^n\sim S^{\otimes n},\qquad
(\xi_t)_{t=1}^n\perp(\eta_t)_{t=1}^n.
\]
Pair their labels coordinatewise as \((\xi_t,\eta_t)\), \(1\leq t\leq n\). Independently across coordinates, choose \(\xi_t\) or \(\eta_t\) with equal probability to obtain a label, choose a uniform sign \(\sigma\), and choose an independent fair treatment indicator \(a\). Conditional on these choices, generate a Bernoulli outcome with parameter \(1/4\) if \(a=0\), with parameter \(1/2+\sigma/8\) if \(a=1\) and the chosen label came from \(\xi_t\), and with parameter \(1/2-\sigma/8\) if it came from \(\eta_t\). At \(n=0\), the kernel assigns probability one to the empty output sequence.

For clarity, define the Bernoulli atom mass and the two outcome parameters, for \(y,a\in\{0,1\}\) and \(\sigma\in\{+1,-1\}\), by
\[
\operatorname{Ber}_{\rho}(y)
=\begin{cases}1-\rho,&y=0,\\ \rho,&y=1,\end{cases}
\quad(0\leq\rho\leq1),
\qquad
\nu^\pm_{a,\sigma}
=\begin{cases}\frac14,&a=0,\\[2pt]\frac12\pm\frac{\sigma}{8},&a=1.\end{cases}
\]
The conditional mass of one output record is therefore
\[
\Pr\{O_t=((i,\sigma),a,y)\mid \xi_t,\eta_t\}
=\frac{
 \mathbf1\{\xi_t=i\}\operatorname{Ber}_{\nu^+_{a,\sigma}}(y)
+\mathbf1\{\eta_t=i\}\operatorname{Ber}_{\nu^-_{a,\sigma}}(y)
}{8}.
\] % lean: CausalSmith.Stat.DiscreteBudgetvalueCurve.pairedKernelAtom
The factors \(1/8\) come from the three fair choices. Summing over labels, signs, treatments, and outcomes gives one, so this specifies a probability kernel.

Mixing this conditional mass over the independent input labels gives
\[
\Pr\{O_t=((i,\sigma),a,y)\}
=\frac{
 R_i\operatorname{Ber}_{\nu^+_{a,\sigma}}(y)
+S_i\operatorname{Ber}_{\nu^-_{a,\sigma}}(y)
}{8}
=\frac{r_i}{4}\operatorname{Ber}_{\mu_{a,i,\sigma}}(y),
\qquad
\mu_{a,i,\sigma}
=\begin{cases}\frac14,&a=0,\\[2pt]\frac12+\sigma\theta_i,&a=1.\end{cases}
\] % lean: CausalSmith.Stat.DiscreteBudgetvalueCurve.pairedKernelAtom_mixed_observed
For \(a=0\), the second equality uses \(r_i=(R_i+S_i)/2\). For \(a=1\), it also uses \(r_i\theta_i=(R_i-S_i)/16\) and the fact that Bernoulli atom masses are affine in their parameter. The equality covers zero-mass pairs as both sides then vanish. Its right-hand side is exactly the observed atom mass of the paired law in \cref{obj:def:paired-family}.

Finally, the input pairs and the kernel randomizations are independent across \(t\). Thus, for every output sequence \(o_t=((i_t,\sigma_t),a_t,y_t)\), \(1\leq t\leq n\),
\[
\Pr\{O_1=o_1,\ldots,O_n=o_n\}
=\prod_{t=1}^n
 \frac{r_{i_t}}4\operatorname{Ber}_{\mu_{a_t,i_t,\sigma_t}}(y_t),
\]
the product law of \(n\) observed records from that paired law. The empty product gives the same conclusion when \(n=0\). % lean: CausalSmith.Stat.DiscreteBudgetvalueCurve.pairedSampleKernel_reproduces
\end{proof}

\bibliographystyle{plainnat}

\bibliography{references}

\end{document}
